\documentclass[11pt,letterpaper]{article}
\usepackage[T1]{fontenc}
\usepackage[utf8]{inputenc}
\usepackage[margin=0.88in,top=0.84in,bottom=0.83in,headheight=13pt,headsep=20pt,footskip=28pt]{geometry}
\usepackage{newpxtext}
\usepackage{amsmath,amsthm}
\usepackage{newpxmath}
\usepackage{microtype}
\usepackage{xcolor}
\definecolor{Forest}{HTML}{30392C}
\definecolor{Olive}{HTML}{687144}
\definecolor{Muted}{HTML}{65685F}
\definecolor{Sage}{HTML}{CBD0BC}
\definecolor{Pale}{HTML}{F2F4EB}
\usepackage{mathtools}
\usepackage{booktabs,tabularx,longtable,array}
\usepackage{enumitem}
\usepackage{titlesec}
\usepackage{fancyhdr}
\usepackage[most]{tcolorbox}
\usepackage{needspace}
\PassOptionsToPackage{obeyspaces}{url}
\usepackage{xurl}
\usepackage[colorlinks=true,linkcolor=Olive,citecolor=Olive,urlcolor=Olive,bookmarksnumbered=true,pdfencoding=auto,psdextra]{hyperref}
\hypersetup{pdftitle={Large cardinals at the inconsistency frontier: synthesis of the continuation reports},pdfauthor={Prepared for Vladimir Reshetnikov},pdfsubject={Deduplicated results of thirty-six research continuations on cover-exacting, exacting and ultraexacting cardinals},pdfkeywords={large cardinals, cover exacting, ultraexacting, HCD, strongly compact, Ground Axiom, I0, I3}}
\urlstyle{same}
\setlength{\parindent}{0pt}
\setlength{\parskip}{5.1pt plus 1pt minus .4pt}
\linespread{1.035}
\setlength{\emergencystretch}{2em}
\setcounter{tocdepth}{1}
\setcounter{secnumdepth}{2}
\titleformat{\section}{\Large\bfseries\color{Forest}}{\thesection}{.6em}{}
\titleformat{\subsection}{\large\bfseries\color{Forest}}{\thesubsection}{.55em}{}
\titleformat{\subsubsection}{\normalsize\bfseries\color{Forest}}{}{0pt}{}
\titlespacing*{\section}{0pt}{18pt plus 3pt minus 2pt}{8pt}
\titlespacing*{\subsection}{0pt}{12pt plus 2pt minus 1pt}{5pt}
\titlespacing*{\subsubsection}{0pt}{9pt}{3pt}
\setlist[itemize]{leftmargin=1.4em,itemsep=2pt,topsep=3pt,parsep=0pt}
\setlist[enumerate]{leftmargin=1.7em,itemsep=3pt,topsep=4pt,parsep=0pt}
\pagestyle{fancy}
\fancyhf{}
\fancyhead[L]{\sffamily\fontsize{9}{11}\selectfont\color{Muted}LARGE CARDINALS AT THE INCONSISTENCY FRONTIER}
\fancyhead[R]{\sffamily\fontsize{9}{11}\selectfont\color{Muted}SYNTHESIS\quad 18 SEP 2026}
\fancyfoot[L]{\small\color{Muted}Synthesis of the continuation reports}
\fancyfoot[R]{\small\color{Muted}\thepage}
\renewcommand{\headrulewidth}{.35pt}
\renewcommand{\footrulewidth}{0pt}
\fancypagestyle{plain}{\fancyhf{}\fancyfoot[L]{\small\color{Muted}Synthesis of the continuation reports}\fancyfoot[R]{\small\color{Muted}\thepage}\renewcommand{\headrulewidth}{0pt}}
\tcbset{colback=Pale,colframe=Sage,colbacktitle=Sage,coltitle=Forest,fonttitle=\bfseries,boxrule=.5pt,arc=3pt,left=9pt,right=9pt,top=6pt,bottom=6pt,toptitle=3pt,bottomtitle=3pt,before skip=8pt,after skip=8pt,breakable}
\newtcolorbox{assessment}[1]{title={#1}}
\theoremstyle{definition}
\newtheorem{definition}{Definition}[section]
\newtheorem{theorem}[definition]{Theorem}
\newtheorem{proposition}[definition]{Proposition}
\newtheorem{lemma}[definition]{Lemma}
\newtheorem{corollary}[definition]{Corollary}
\newtheorem{remark}[definition]{Remark}
\newtheorem{question}[definition]{Question}
\newtheorem{inputthm}{Input}
\renewcommand{\theinputthm}{\Alph{inputthm}}
\numberwithin{equation}{section}
\newcommand{\ZFC}{\mathsf{ZFC}}
\newcommand{\ZF}{\mathsf{ZF}}
\newcommand{\AC}{\mathsf{AC}}
\newcommand{\GA}{\mathsf{GA}}
\newcommand{\HOD}{\mathrm{HOD}}
\newcommand{\OD}{\mathrm{OD}}
\newcommand{\CD}{\mathrm{CD}}
\newcommand{\HCD}{\mathrm{HCD}}
\newcommand{\UF}{\mathrm{UF}}
\newcommand{\Ex}{\operatorname{Ex}}
\newcommand{\UEx}{\operatorname{UEx}}
\newcommand{\Ext}{\operatorname{Ext}}
\newcommand{\SC}{\operatorname{SC}}
\newcommand{\CEx}{\operatorname{CEx}}
\newcommand{\REx}{\operatorname{REx}}
\newcommand{\Con}{\operatorname{Con}}
\newcommand{\crit}{\operatorname{crit}}
\newcommand{\cf}{\operatorname{cf}}
\newcommand{\ot}{\operatorname{ot}}
\newcommand{\ran}{\operatorname{ran}}
\newcommand{\rank}{\operatorname{rank}}
\newcommand{\lcm}{\operatorname{lcm}}
\newcommand{\Ord}{\mathrm{Ord}}
\newcommand{\Pow}{\mathcal P}
\newcommand{\id}{\mathrm{id}}
\newcommand{\restr}{\mathbin{\upharpoonright}}
\newcommand{\Izero}{\mathrm{I0}}
\newcommand{\Itwo}{\mathrm{I2}}
\newcommand{\Ithree}{\mathrm{I3}}
\newcommand{\Iwf}{\mathrm{I3}_{\mathrm{wf}}(0)}
\newcommand{\Sel}{\operatorname{Sel}}
\newcommand{\FF}{\mathcal F}
\newcommand{\PP}{\mathbb P}
\newcommand{\Clam}{\mathcal C_\lambda}
\newcommand{\Rig}{\operatorname{Rig}}
\newcommand{\Spec}{\operatorname{Spec}}
\newcommand{\ch}{\operatorname{ind}}
\newcommand{\CF}{\mathsf{CF}}
\newcommand{\src}[1]{\unskip\ {\normalfont\small\sffamily\color{Olive}[#1]}\ }
\newcommand{\R}[1]{\textsf{R#1}}
\newcolumntype{Y}{>{\raggedright\arraybackslash}X}
\newcolumntype{P}[1]{>{\raggedright\arraybackslash}p{#1}}
\newcolumntype{C}{>{\centering\arraybackslash}p{.034\textwidth}}
\newcommand{\yes}{$\bullet$}
\newcommand{\prt}{$\circ$}
\newcommand{\arxiv}[1]{\href{https://arxiv.org/abs/#1}{arXiv:#1}}
\newcommand{\doi}[1]{\href{https://doi.org/#1}{doi:\nolinkurl{#1}}}

\begin{document}
\thispagestyle{empty}
{\sffamily\bfseries\small\color{Olive}SET THEORY\quad /\quad RESEARCH SYNTHESIS}

\vspace{13pt}
{\fontsize{28}{31}\selectfont\bfseries\color{Forest}Large cardinals at the\\ inconsistency frontier:\\ a synthesis\par}
\vspace{10pt}
{\fontsize{14}{18}\selectfont Deduplicated results of the thirty-six continuation reports on\\ cover-exacting, exacting, and ultraexacting cardinals\par}
\vspace{14pt}
{\color{Muted}Prepared for Vladimir Reshetnikov\hfill 18 September 2026\par}
\vspace{3pt}
{\color{Sage}\hrule height .6pt}
\vspace{12pt}

\textbf{Purpose.} Thirty-six independent continuation reports were written in four rounds of nine. The first nine (\R1--\R9, in \nolinkurl{docs/cardinals/research-reports/01}--\nolinkurl{09}) continue the research plan, the \emph{Unified Report}; seven of them prove essentially the same central theorem with different names, thresholds, and proof routes. The second nine (\R{10}--\R{18}, in \nolinkurl{10}--\nolinkurl{18}) continue the first edition of this synthesis; all nine prove the same sharp-width theorem for definable sections of the cofinal quotient, again with different refinements. The third nine (\R{19}--\R{27}, in \nolinkurl{19}--\nolinkurl{27}) continue the second edition; four of them prove the same theorem on internal normal measures and Prikry cores, and one proves an outright non-coexistence theorem for cover exactingness. The fourth nine (\R{28}--\R{36}, in \nolinkurl{28}--\nolinkurl{36}) continue the third edition; eight of them answer the same question of that edition, by the same device, in the Prikry extension of a model with one measurable cardinal. This document merges them. Each result is stated once, in the strongest form reached by any report, with a compact proof and a source tag such as {\small\sffamily\color{Olive}[R1, R6]}. Section~\ref{sec:compare} compares the reports and records the discrepancies found.

\begin{assessment}{The eight result groups}
\textbf{A. Completeness barrier} (\R1--\R4, \R6--\R8). If $\lambda$ is $\gamma$-cover exacting then $\lambda$ is regular in $\HCD(\gamma^+)$, and in $\HCD(\gamma)$ when $\gamma$ is singular, although $\cf^V(\lambda)=\omega$. A strongly compact cardinal \emph{below} $\lambda$ forces a small, explicit witness to non-stabilization of the $\HCD$ hierarchy; a strongly compact \emph{above} $\gamma$ forces a proper ground and contradicts the Ground Axiom.

\textbf{B. Exacting cardinals} (\R9; \R{10}, \R{14}, \R{17}, \R{18}). No $\OD_{V_\lambda}$ family of fewer than $\lambda$ short cofinal subsets of $\lambda$ exists; small definable families of subsets of $V_\lambda$ are separated at a bounded rank; no definable countably additive law on cofinal sequences exists.

\textbf{C. Ultraexacting cardinals} (\R5, \R7). Definable finite choice fails on the quotient of cofinal $\omega$-subsets of $\lambda$ modulo finite difference: no finite $\OD_{V_\lambda}$ family of selectors, no ordinal-definable finite-valued transversal, and a forbidden Icarus enrichment.

\textbf{D. Consistency} (\R1--\R3, \R5, \R7--\R9). $\Con(\ZFC+\Iwf)$ implies the consistency of a $\lambda$-cover-exacting $\lambda$; and one $\Izero$-equiconsistent theory carries all the boundary phenomena of groups A--C simultaneously.

\textbf{E. Class embeddings} (\R3). For $j:V\to M$ and a fixed regular $\theta$, no unbounded subset of $j``\theta$ belongs to $M$.

\textbf{F. Sections of the cofinal quotient} (\R{10}--\R{18}). At an ultraexacting $\lambda$ there are at least $\lambda$ \emph{rigid} classes $q$ of cofinal $\omega$-sets modulo finite difference: every set of representatives definable from $q$ and parameters in $V_\lambda$ meets $q$ in $0$ or $\lambda$ points. Hence every definable complete section has $\lambda$ fibres of full size $\lambda$, and the definable width of the quotient is exactly $\lambda$. The equivariant finite-label problem of group~C is solved exactly: a definable rule exists iff every single permutation fixes a label (\R{13}). With a strongly compact cardinal below $\lambda$, an inner model in which $\lambda$ is regular has its successor of $\lambda$ collapsed or a canonical stationary set destroyed (\R{14}).

\textbf{G. Third round} (\R{19}--\R{27}). \emph{No strongly compact sandwich:} in $\ZFC$, a $\gamma$-cover-exacting $\lambda$ cannot have strongly compact cardinals both below $\lambda$ and above $\gamma$ (\R{24}); this removes the Ground Axiom from Corollary~\ref{cor:ga} in that situation and makes the two-sided configurations of groups A and F vacuous. \emph{Hidden measures:} for the critical-sequence class $q$ of an ultraexacting witness, the tail filter of $q$ is a normal measure on $\lambda$ inside $\HOD_{\{q\}}$, every representative of $q$ is Prikry generic over that model, and $(\lambda^+)^V$ is measurable there (\R{19}, \R{22}, \R{23}, \R{26}, \R{27}). \emph{Exact strength:} the existence of a rigid class, and hence all of group F except the finite-label necessity, is equiconsistent with one measurable cardinal; the normal-trace principle with concentration on measurables, with Mitchell order two (\R{25}, \R{27}). \emph{Charges:} a finitely additive definable kernel exists, an ultrafilter-valued one does not (\R{21}).

\textbf{H. Fourth round: the Prikry model} (\R{28}--\R{36}). In the extension $W[c]$ of a model $W$ with a normal measure on $\kappa$ by Prikry forcing, with $q=[c]$ and definitions allowed to use $q$, ordinals and arbitrary ground-model sets: the exact finite-label classification of group~F holds, and there is no nonempty finite definable family of $r$-of-$n$ selectors on $Q_\kappa$ (eight reports; the device is to insert one point into the generic sequence, which changes neither the extension nor $q$ but rotates coded tuples of classes). So the last conclusions of groups C and F that were tied to an embedding already hold at the strength of one measurable cardinal. The least size of a nonempty definable family of ultrafilters on $q$ is exactly $\aleph_0$, and with a ground-model parameter there are countably infinite definable families of total selectors and of total ultrafilter-valued kernels (\R{34}). There are $2^\kappa$ jointly rigid classes (\R{31}, \R{32}, \R{35}). Separately, at an ultraexacting $\lambda$ the core $\HOD_{\{b,q\}}$ regards $\lambda$ as a measurable Vop\v enka cardinal (\R{29}).
\end{assessment}

\textbf{Status.} Nothing here refutes a standard large-cardinal hypothesis; all thirty-six reports agree on that. The strongest negative result is Theorem~\ref{thm:sandwich}, a $\ZFC$ theorem that answers the Blue--Goldberg Problem~5.2 negatively \emph{when there is a further strongly compact cardinal above the cover bound}; without one the problem remains open. The proofs are conventional, unrefereed, and not machine-checked. The synthesis re-derived the arguments while merging them, but did not re-verify the external sources; theorem numbers for imported results are those quoted by the reports (those used in the Lean formalization were later checked against the sources; see \nolinkurl{Cardinals/README.md}). Priority is not claimed for any statement. \emph{Second edition:} Sections~\ref{sec:above}, \ref{sec:width}, \ref{sec:cons}, \ref{sec:compare} and \ref{sec:open} were amended and Section~\ref{sec:quot} was added. \emph{Third edition:} Section~\ref{sec:above} was amended again and Section~\ref{sec:measures} was added. \emph{Fourth edition:} Section~\ref{sec:measures} gained Subsection~\ref{sec:vopenka} and Section~\ref{sec:prikry} was added; Section~\ref{sec:countable}, on countable analogues, is an addition by the synthesis and does not come from any report. The numbering of all earlier statements is unchanged.

\newpage
\tableofcontents

\newpage
\section{The source reports}\label{sec:sources}

The original archive folders had colliding names (several were called \nolinkurl{Large_Cardinals_Research_Continuation}, with browser-style suffixes). They were renamed to numbered directories in order of file modification time, and one redundant level of nesting was removed.

\begin{center}
\small
\renewcommand{\arraystretch}{1.18}
\begin{tabularx}{\textwidth}{@{}lP{.37\textwidth}Y@{}}
\toprule
&\textbf{Title (main file)}&\textbf{Distinctive content}\\\midrule
\R1&\emph{Frontier II: Cofinality barriers for cover-exacting cardinals} (\nolinkurl{Large_Cardinals_Continuation})&Barrier for all $\eta>\gamma$; promotion principle $\mathsf{LP}$; ground theorem with forcing-size bound; failed-stationarity certificate; $\Itwo$ calibration.\\
\R2&\emph{Cover exactingness, definability, and conditional inconsistency} (\nolinkurl{Cover_Exacting_Definability_Followup})&No-small-cover form of the gap; weak covering principle $\mathsf{Cov}$; first $\Izero$ boundary package.\\
\R3&\emph{Local barriers at the large-cardinal frontier} (\nolinkurl{Large_Cardinals_Research_Continuation})&Chain-condition form of the ground obstruction; internal unbounded-range theorem for $j:V\to M$; first-HOD-gap theory.\\
\R4&\emph{Cover-exacting cardinals and the Ground Axiom} (\nolinkurl{Large_Cardinals_Research_Continuation})&Barrier proved \emph{without} the stationary characterization (uniform covers, envelopes, absorption); ultrafilter-code promotion failure.\\
\R5&\emph{Finite choice at the ultraexacting boundary} (\nolinkurl{Large_Cardinals_Finite_Choice_Continuation})&Finite-family selector obstruction with cycle amplification; selector-coded Icarus corollary; bounded/cofinal quotient separation; finite checks (Python).\\
\R6&\emph{Cover-exacting cardinals and high-completeness cores} (\nolinkurl{Cover_Exacting_High_Completeness_Cores})&Singular endpoint $q(\gamma)$; proof of $\gamma\ge\lambda$; cofinal-refinement principle; first-regularization theorem.\\
\R7&\emph{Local complete-definability barriers, finite symmetry, and an $\Izero$-calibrated consistency result} (\nolinkurl{Large_Cardinals_Continuation})&Finite-label symmetry theorem; finite-valued transversal theorem; small-ground formulation; finite checks (Python).\\
\R8&\emph{Cover-exacting cardinals and completely definable grounds} (\nolinkurl{Large_Cardinals_Continuation})&$\lambda$ strongly inaccessible in the ground; rank-embedding witnesses escape $\CD(\gamma^+)$; first-disagreement profile.\\
\R9&\emph{Frontier: a Prikry consistency refinement and local definability obstructions} (\nolinkurl{Large_Cardinals_Research_Continuation})&Prikry construction of a cover-exacting cardinal from $\Iwf$; projection-width theorem for exacting cardinals.\\
\bottomrule
\end{tabularx}
\end{center}

The second round arrived as nine zip archives with overlapping internal folder names (three of them contained a folder \nolinkurl{Cardinals3_Continuation} or \nolinkurl{Cardinals_Continuation}). They were extracted into \nolinkurl{10}--\nolinkurl{18}, again in order of source modification time and without the extra nesting level.

\begin{center}
\small
\renewcommand{\arraystretch}{1.18}
\begin{tabularx}{\textwidth}{@{}lP{.33\textwidth}Y@{}}
\toprule
&\textbf{Title (main file)}&\textbf{Distinctive content}\\\midrule
\R{10}&\emph{Beyond finite choice} (\nolinkurl{Beyond_Finite_Choice})&Small invariant families at \emph{exacting} cardinals; no conversion of classes into short cofinal information; probability laws and kernels; bounded-separation characterization.\\
\R{11}&\emph{Quotient parameters and sharp uniformization barriers} (\nolinkurl{Large_Cardinals_Quotient_Continuation})&The models $N_q$; classes that do and do not reveal cofinality; $\lambda$ recodings with the same $N_q$; minimum-order-type compression.\\
\R{12}&\emph{Maximal-width uniformization} (\nolinkurl{Maximal_Width_Uniformization})&Roots and orbits of $e$; coordinate width; integer phase spectrum; envelopes and edit budgets; two sharpness examples. Finite checks (Python).\\
\R{13}&\emph{Cyclic symmetry at the ultraexacting boundary} (\nolinkurl{Cyclic_Symmetry_Ultraexacting})&Exact finite-label classification via cyclic stabilizers of incidence words; the five-label rule; both output types occur $\lambda$ times. Finite checks (Python).\\
\R{14}&\emph{Small definable sections and stationary-correctness barriers} (\nolinkurl{Cardinals_Continuation})&Small collections of sections; successor stationary-set obstruction; $(\lambda^+)$-chain-condition bound; localization of the destroyed stationary set.\\
\R{15}&\emph{Orbit saturation} (\nolinkurl{Orbit_Saturation})&Zero-or-full width for small families of \emph{partial} sections; thin covers; choiceless inner models; classification and count of periodic classes; $|Q_\lambda|=2^\lambda$.\\
\R{16}&\emph{Small fibres at the ultraexacting boundary} (\nolinkurl{Small_Fibres_Ultraexacting})&Tail-invariant hulls; ordinal colourings; coarser equivalence relations; predicate-expanded $\mathrm{I1}$ embeddings.\\
\R{17}&\emph{Thin sections and simultaneous saturation} (\nolinkurl{Thin_Sections_and_Saturation})&One set of $\lambda$ classes works for \emph{all} $\OD_{V_\lambda}$ sets at once; fixed bijection $\lambda\to V_\lambda$; sharpness of the parameter rank.\\
\R{18}&\emph{Cofinal-orbit rigidity} (\nolinkurl{Cofinal_Orbit_Rigidity})&Rigidity relative to the class $q$ itself as a parameter; the minimal-rank-parameter device; homomorphisms into relations with small classes.\\
\bottomrule
\end{tabularx}
\end{center}

The third round arrived in the same way and was extracted into \nolinkurl{19}--\nolinkurl{27}.

\begin{center}
\small
\renewcommand{\arraystretch}{1.18}
\begin{tabularx}{\textwidth}{@{}lP{.33\textwidth}Y@{}}
\toprule
&\textbf{Title (main file)}&\textbf{Distinctive content}\\\midrule
\R{19}&\emph{Fixed tail parameters and stationary collapse} (\nolinkurl{Fixed_Tail_Parameters})&$(\lambda^+)^V$ is measurable in $\HOD_{b,\vec q_F}$ for every finite set of orbit classes; one club of simultaneous inaccessibles; no $\lambda^+$-cc reconstruction.\\
\R{20}&\emph{Cofinality cascades above strongly compact cardinals} (\nolinkurl{Cofinality_Cascades})&A general forcing theorem: singularizing $\lambda$ below cofinality $\delta$ above a strongly compact $\delta$ forces a cascade of $\delta$ ground-regular cardinals to change cofinality, and antichains of size $(\lambda^{+\delta})^W$.\\
\R{21}&\emph{Finitely additive kernels} (\nolinkurl{Finitely_Additive_Kernels})&The definable tail-mean kernel; common null partitions; phase balance of definable charges; no definable ultrafilter-valued kernel.\\
\R{22}&\emph{Canonical tail measures and Prikry cores} (\nolinkurl{Canonical_Tail_Measures})&Normal tail measure in $\HOD_{\{q\}}$ and in $\HOD_{V_\lambda\cup\{q\}}$; Prikry core; three mutually definable classes with normal, non-normal and non-ultra tail filters.\\
\R{23}&\emph{Normal measures and Prikry layers} (\nolinkurl{Normal_Measures_and_Prikry_Layers})&One class $q$ serving every parameter $p\in V_\lambda$; coherence of the measures $U_{q,p}$; collapse-or-destroy for the Prikry layer.\\
\R{24}&\emph{Cover exactingness between strongly compact cardinals} (\nolinkurl{Cover_Exacting_Stationary_Seeds})&The no-sandwich theorem via $\omega$-club amenability of $\HCD(\eta)$ and stationary seeds; stationary-width bound; measurability of large regulars in high $\HCD$ models; a weaker lower hypothesis.\\
\R{25}&\emph{The measurable strength of cofinal-orbit rigidity} (\nolinkurl{Cofinal_Orbit_Consistency})&Rigidity $\Leftrightarrow$ regularity in $N_q$; equiconsistency with a measurable (Prikry forcing; Dodd--Jensen); the Vop\v enka algebra of a class.\\
\R{26}&\emph{Canonical Prikry cores} (\nolinkurl{Ultraexacting_Prikry_Cores})&Cores with and without the fixed bijection $b$; the $\ZF$ core $\HOD_{V_\lambda\cup\{q\}}$; definable-splitting principle $\mathsf{DS}$.\\
\R{27}&\emph{Normal measures hidden in critical-sequence quotients} (\nolinkurl{Normal_Traces_and_Prikry_Factorization})&Concentration on measurables; exact strength of the trace principles $\mathsf{NT}$ (a measurable) and $\mathsf{NT}^+$ ($o(\kappa)\ge2$).\\
\bottomrule
\end{tabularx}
\end{center}

The fourth round was extracted into \nolinkurl{28}--\nolinkurl{36}. Two archives had the same internal folder name (\nolinkurl{Cardinals5_Prikry_Symmetry}) with different contents; they are \R{30} and \R{31}.

\begin{center}
\small
\renewcommand{\arraystretch}{1.18}
\begin{tabularx}{\textwidth}{@{}lP{.33\textwidth}Y@{}}
\toprule
&\textbf{Title (main file)}&\textbf{Distinctive content}\\\midrule
\R{28}&\emph{Finite symmetry at measurable strength} (\nolinkurl{Finite_Symmetry_at_Measurable_Strength})&Deletion of a generic point; sets of finite restriction tables; the countable ultrafilter orbit in $\ZFC$; finite checks (Python).\\
\R{29}&\emph{Vop\v enka reflection in hidden Prikry cores} (\nolinkurl{Vopenka_Reflection})&The only report on ultraexacting cardinals: the tail measure contains every internal predicate-extendibility set; $\HOD_{\{b,q\}}$ sees a measurable Vop\v enka cardinal.\\
\R{30}&\emph{Prikry symmetry at the measurable threshold} (\nolinkurl{Prikry_Symmetry})&Decide first, then insert; block coding $\omega\cdot a_n+\pi^{-n}(i)$; separation from rank-into-rank at the least measurable.\\
\R{31}&\emph{Prikry symmetry and maximal quotient rigidity} (\nolinkurl{Prikry_Symmetry_and_Maximal_Rigidity})&$2^\kappa$ rigid classes $q_A$, $A\in\Pow(\kappa)^W$, with measurable cores; finite joint parameters.\\
\R{32}&\emph{Maximal rigidity and finite symmetry from one measurable} (\nolinkurl{Maximal_Rigidity_and_Finite_Symmetry})&Tail-invariant parameters; one indexed family of all classes; rank capture: the core computes $\kappa^+$ and its Prikry layer contains $V_{\kappa+1}$.\\
\R{33}&\emph{Finite definable choice from one measurable cardinal} (\nolinkurl{Prikry_Finite_Symmetry})&Finite families of label rules; local constancy of definable maps on $q$; full phase and coordinate saturation with ground parameters.\\
\R{34}&\emph{Prikry symmetry and the finite--countable choice gap} (\nolinkurl{Prikry_Choice_Gap})&Countably infinite definable families of \emph{total} selectors and kernels from a ground-model name presentation; a coded model where they are ordinal definable.\\
\R{35}&\emph{Finite choice and maximal rigidity at one measurable cardinal} (\nolinkurl{Prikry_Choice_and_Maximal_Rigidity})&Invariant parameters $z(c)$; $2^\kappa$ classes coded in one parameter with a normal measure in the joint core; the ultrafilter cloud as a $\mathbb Z$-torsor.\\
\R{36}&\emph{Finite choice at measurable strength} (\nolinkurl{Cardinals5_Prikry_Finite_Choice})&Profinite density of phases; compact families need $2^{2^{\aleph_0}}$ ultrafilters; residue laws for finite families of charges, sharp; the core's least measurable.\\
\bottomrule
\end{tabularx}
\end{center}

A result-by-report concordance is given in Appendix~\ref{app:concordance}.

\section{Conventions and imported results}\label{sec:setup}

\subsection{Conventions}

We work in $\ZFC$. Unadorned cardinalities, completeness, clubs and stationarity are computed in $V$; $|x|^N$ and $\cf^N$ are computed in an inner model $N$. For an inner model $N\models\ZF$, ``$\cf^N(\lambda)=\lambda$'' means that $N$ contains no cofinal map from an ordinal below $\lambda$ into $\lambda$; no Choice in $N$ is needed. All embeddings in Sections~\ref{sec:core}--\ref{sec:uex} are sets, and elementarity is expressed through satisfaction for set structures.

\begin{definition}[Exacting, ultraexacting]\label{def:ex}
A cardinal $\lambda$ is \emph{exacting}, $\Ex(\lambda)$, if for every cardinal $\zeta>\lambda$ there are $X\prec V_\zeta$ with $V_\lambda\cup\{\lambda\}\subseteq X$ and an elementary $j:X\to V_\zeta$ with $j(\lambda)=\lambda$ and $j\restr\lambda\neq\id$. It is \emph{ultraexacting}, $\UEx(\lambda)$, if moreover $j\restr V_\lambda\in X$ can be arranged. The domain $X$ need not be transitive.
\end{definition}

\begin{definition}[Relative and cover exactingness]\label{def:cex}
For $\alpha>\lambda$ and $Y\subseteq V_\alpha$, a \emph{relative $\lambda$-exacting witness} is an elementary $j:(X,\in,X\cap Y)\to(V_\alpha,\in,Y)$ with $(X,\in,X\cap Y)\prec(V_\alpha,\in,Y)$, $V_\lambda\cup\{\lambda\}\subseteq X$, $\crit(j)<\lambda$, $j(\lambda)=\lambda$. For a set $Y$, write $\REx(\lambda,Y)$ if such witnesses exist at every $\alpha>\lambda$ with $Y\subseteq V_\alpha$. For infinite cardinals $\lambda\le\gamma$, $\lambda$ is \emph{$\gamma$-cover exacting}, $\CEx_\gamma(\lambda)$, if
\[
\forall\alpha>\lambda\ \forall y\in V_\alpha\ \exists Y\subseteq V_\alpha\ \bigl(y\in Y\wedge|Y|\le\gamma\wedge\text{there is a relative witness for }(V_\alpha,Y)\bigr).
\]
The bound $\gamma$ is fixed before $\alpha$ and $y$; the condition is $y\in Y$, not $y\subseteq Y$; the witness need not fix $y$.
\end{definition}

Some reports define $\REx$ only at all sufficiently large heights. By Lemma~\ref{lem:descent} the two conventions agree for every use below.

\begin{definition}[Complete definability]\label{def:hcd}
Let $\UF_\eta(\Ord)$ be the class of proper $\eta$-complete ultrafilters on ordinals (closed under intersections of fewer than $\eta$ members). $\CD(\eta)$ is the class of sets definable from finitely many ordinals and finitely many members of $\UF_\eta(\Ord)$, and $\HCD(\eta)=\{x:\operatorname{tc}(\{x\})\subseteq\CD(\eta)\}$; $\HCD=\bigcap_\eta\HCD(\eta)$. For $\eta\le\xi$,
\begin{equation}\label{eq:mono}
\CD(\xi)\subseteq\CD(\eta),\qquad\HCD(\xi)\subseteq\HCD(\eta),
\end{equation}
and for a set of ordinals $a$: $a\in\CD(\eta)\iff a\in\HCD(\eta)$.
\end{definition}

Put $q(\gamma)=\gamma$ if $\gamma$ is singular and $q(\gamma)=\gamma^+$ otherwise. A \emph{ground} is an inner model $W\models\ZFC$ with $V=W[G]$ for a set forcing $\PP\in W$; the \emph{Ground Axiom} $\GA$ says there is no proper ground.

\subsection{Imported results}

\begin{inputthm}[Local Kunen]\label{in:kunen}
In $\ZFC$ there is no nontrivial elementary $V_{\mu+2}\to V_{\mu+2}$. \cite{kunen,hkp}
\end{inputthm}

\begin{inputthm}[Goldberg's $\HCD$ theorems]\label{in:hcd}
$\HCD(\eta)\models\ZF$ for every infinite cardinal $\eta$. If $\delta$ is strongly compact, then $\HCD(\delta)\models\ZFC$, it is a ground of $V$ by a forcing of size below $(2^{2^\delta})^+$, and it has the $\delta$-cover and $\delta$-approximation properties in $V$. If $\delta$ is extendible, $\HCD(\delta)=\HCD$. \cite[Prop.~4.2--4.3, Thm.~4.4, 4.10, Prop.~4.11, Thm.~4.14--4.15]{unique}
\end{inputthm}

Several reports stress that the ground and covering statements need only \emph{strong compactness}; the supercompact hypothesis in the shorter lecture-note summary is not needed. \R1 notes that the arXiv~v2 numbering is shifted by one relative to the author version used here.

\begin{inputthm}[Blue--Goldberg cover-exacting notes]\label{in:bg}
(i) $\CEx_\gamma(\lambda)$ implies $\gamma\ge\lambda$. (ii) Lemma~3.3: $\REx(\lambda,Y)$ and $Z\in\OD_{\{Y\}}$ imply $\REx(\lambda,Z)$. (iii) Proposition~3.4: $\CEx_\gamma(\lambda)$ iff every set belongs to some $Y$ with $|Y|\le\gamma$ and $\REx(\lambda,Y)$, iff $\{\sigma\in P_{\gamma^+}(V_\alpha):\REx(\lambda,\sigma)\}$ is stationary for all large $\alpha$. (iv) Theorem~3.5: $\Con(\ZFC+\Itwo)\Rightarrow\Con(\ZFC+\exists\lambda\,\CEx_\lambda(\lambda))$. (v) Theorem~3.6: no $\gamma$-cover-exacting cardinal lies above an extendible. (vi) Problem~5.2: can one lie above a strongly compact? \cite{cover}
\end{inputthm}

Items (i)--(iii) are lecture-note results; \R4 and \R6 reprove everything needed from them (Section~\ref{sec:hulls}), so group~A depends on the notes only through the definitions.

\begin{inputthm}[Exacting and ultraexacting cardinals]\label{in:abl}
(i) An exacting $\lambda$ is a strong limit of cofinality $\omega$ and is regular in $\HOD_{V_\lambda}$ \cite[Thm.~2.10]{abl}. (ii) Ultraexacting witnesses can be chosen with critical point above any $\alpha<\lambda$ \cite[Lemma~3.2]{abl}. (iii) $\lambda$ is ultraexacting iff every ordinal-definable $A\subseteq V_{\lambda+1}$ admits an elementary $e:(V_{\lambda+1},A)\to(V_{\lambda+1},A)$ with $\crit(e)<\lambda$ \cite[Thm.~B/3.1]{abgl}. (iv) $\ZFC+\exists\lambda\,\UEx(\lambda)$ and $\ZFC+\Izero$ are equiconsistent \cite[Thm.~A/4.5]{abgl}. (v) An ultraexacting $\lambda$ is preserved by a $\chi$-directed closed forcing (any $\chi>\lambda$) forcing $V_\lambda\subseteq\HOD$; then no exacting cardinal lies below $\lambda$ \cite[Prop.~3.9, Cor.~3.10]{abgl}.
\end{inputthm}

\begin{inputthm}[Iteration interface for $\Iwf$]\label{in:iter}
For $j:V_\Lambda\to V_\Lambda$ with critical sequence $\langle\kappa_n\rangle$, finite iterates $j_{m,n}$, direct limit $M^j_\omega$ with well-founded part $W^j_\omega$: $j$ is an $\Iwf$-embedding if $j_{0,\omega}[\kappa_0^+]\subseteq W^j_\omega$. Then $V_\Lambda\cup\{\Lambda\}\subseteq W^j_\omega$; $j^+(j_{0,\omega}(A))=j_{0,\omega}(A)$ for $A\in V_{\kappa_0+1}$; if $j^+(E)=E\subseteq\Lambda^2$ then $j\restr\Lambda$ is elementary on $(\Lambda,E)$; and $A\in j_{0,\omega}(U_0)$ iff $A$ contains a tail of $\langle\kappa_n\rangle$, for $A\in\Pow(\Lambda)\cap W^j_\omega$. Moreover $\Itwo$ is strictly stronger than $\Iwf$, and $\exists\lambda\,\Ex(\lambda)$ is strictly stronger than $\Ithree$, in consistency strength. \cite[\S3]{ad}, \cite[Prop.~5.1, Def.~5.2, Thm.~5.4, 5.10, Cor.~5.11]{abgl}
\end{inputthm}

\section{Core lemmas on exacting witnesses}\label{sec:core}

\begin{lemma}[Critical-sequence normalization]\label{lem:crit}\src{R1--R8}
Let $\lambda$ be an uncountable cardinal and let $e$ be a nontrivial elementary self-embedding of $V_\lambda$, or of $V_{\lambda+1}$ with $\crit(e)<\lambda$. With $\kappa_n=e^n(\crit(e))$ we have $\sup_n\kappa_n=\lambda$, and $e(\xi)>\xi$ for every $\xi\in[\kappa_0,\lambda)$. In particular every exacting witness $j$ restricts to such an $e=j\restr V_\lambda$, so $\cf^V(\lambda)=\omega$ and every ordinal below $\lambda$ fixed by $j$ lies below $\crit(j)$.
\end{lemma}

\begin{proof}
If $\mu=\sup_n\kappa_n<\lambda$, the sequence lies in $V_\lambda$, $e(\mu)=\mu$, and $e\restr V_{\mu+2}$ contradicts Input~\ref{in:kunen}. For $\kappa_n\le\xi<\kappa_{n+1}$, $e(\xi)\ge\kappa_{n+1}>\xi$. The restriction $j\restr V_\lambda$ is elementary because $V_\lambda$ is definable from the fixed parameter $\lambda$ and $V_\lambda\subseteq X$; transitivity of $X$ is not used.
\end{proof}

\begin{lemma}[No preserved short cofinal set]\label{lem:short}\src{R1, R2, R4, R6, R7}
If $a\subseteq\lambda$ is cofinal with $\ot(a)<\lambda$, then $\REx(\lambda,a)$ fails.
\end{lemma}

\begin{proof}
In $(V_\alpha,\in,a)$ the set $a$ is definable from its predicate, so $a\in X$ and $j(a)=a$; hence $\tau=\ot(a)$ and the increasing enumeration $h:\tau\to a$ are in $X$ and fixed. By Lemma~\ref{lem:crit}, $\tau<\crit(j)$, so $j(h(\xi))=j(h)(j(\xi))=h(\xi)$ for all $\xi<\tau$. Thus the cofinal set $a$ consists of fixed ordinals, contradicting Lemma~\ref{lem:crit}.
\end{proof}

\begin{lemma}[Predicate bookkeeping]\label{lem:descent}\src{R4, R6}
Let $j$ be a relative witness for $(V_\theta,Y)$ with $Y\in V_\theta$. (a) $Y\in X$ and $j(Y)=Y$; members of $Y$ need not be fixed. (b) If $\cf(\theta)>\omega$, there are relative witnesses for $(V_\beta,Y)$ at every $\beta<\theta$ with $\beta>\lambda$ and $Y\subseteq V_\beta$.
\end{lemma}

\begin{proof}
(a) $Y$ is the unique $z$ with $\forall u(u\in z\leftrightarrow Y(u))$. (b) Existence of a witness at a height $\beta<\theta$ is computed correctly in $V_\theta$. The least failing $\beta$ is definable from $\lambda$ and $Y$, hence lies in $X$ and is fixed; restricting $j$ to $X\cap V_\beta$ gives a witness at $\beta$.
\end{proof}

\begin{lemma}[Definability transfer]\label{lem:transfer}\src{R4, R6; focused form R1--R3, R8}
If $\REx(\lambda,Y)$ and $Z\in\OD_{\{Y\}}$, then $\REx(\lambda,Z)$. In particular no cofinal subset of $\lambda$ of order type below $\lambda$ belongs to $\OD_{\{Y\}}$, and $\lambda$ is regular in $\HOD_{\{Y\}}$.
\end{lemma}

\begin{proof}
Let $Z$ be the least counterexample in the canonical well-order of $\OD_{\{Y\}}$ and $\beta$ its least failing height; both are definable from $\lambda$ and $Y$ \emph{without further ordinal parameters}. Take a witness for $Y$ at a height $\theta$ reflecting these definitions. Then $Z,\beta\in X$ are fixed, the predicate $u\in Z$ is a definitional expansion, and restriction to $V_\beta$ gives the excluded witness. The second assertion follows from Lemma~\ref{lem:short}.
\end{proof}

\begin{assessment}{The parameter issue, as every report records it}
An arbitrary ordinal parameter in a definition need not be fixed by an exacting embedding. All proofs therefore first pass to a canonical least object definable from the fixed pair $(Y,\lambda)$ and only then choose a witness at a reflecting height. The same device reappears in Lemma~\ref{lem:canonical} for ultraexacting cardinals.
\end{assessment}

\begin{lemma}[The cover bound]\label{lem:bound}\src{R6; known from the notes}
$\CEx_\gamma(\lambda)$ implies $\gamma\ge\lambda$.
\end{lemma}

\begin{proof}
If $\gamma<\lambda$, take a cofinal $c:\omega\to\lambda$, a height $\alpha$ with $\cf(\alpha)>\gamma$, and a cover $Y\ni c$ with witness $j$. Then $Y\in V_\alpha$ is fixed, $\nu=|Y|<\crit(j)$, and a bijection $f:\nu\to Y$ in $X$ shows $Y\subseteq X$ and $j``Y=Y$. Let $c_n\in Y$ be the $n$-fold preimage of $c$; then $c(k)=e^n(c_n(k))$ with $e=j\restr V_\lambda$. Choosing $c(k)\ge\crit(j)$ and $n$ with $e^n(\crit(j))>c(k)$ gives a contradiction, since $e^n$ maps ordinals below $\crit(j)$ below $\crit(j)$ and the others to at least $e^n(\crit(j))$.
\end{proof}

\section{From covers to small exact hulls}\label{sec:hulls}

The reports use two routes to obtain a small elementary substructure $\sigma\prec V_\alpha$ that is simultaneously (i) of size at most $\gamma$, (ii) a container for prescribed parameters, and (iii) a predicate with $\REx(\lambda,\sigma)$.

\textbf{Route S} (\R1--\R3, \R7, \R8) imports the stationary characterization, Input~\ref{in:bg}(iii), and intersects the stationary set with the club of elementary hulls containing the parameters.

\textbf{Route H} (\R4, \R6) proves the required hull directly from Definition~\ref{def:cex}:

\begin{lemma}[Persistent small covers]\label{lem:persist}\src{R4, R6}
$\CEx_\gamma(\lambda)$ holds iff every set $y$ belongs to a set $Y$ with $|Y|\le\gamma$ and $\REx(\lambda,Y)$.
\end{lemma}

\begin{proof}
($\Rightarrow$) If $y$ has no such cover, let $b(Y)$ be the least failing height of each candidate $Y$, and let $F(\eta)$ bound $b(Y)+\omega$ over candidates $Y\subseteq V_\eta$; this uses Replacement on sets only. A continuous recursion of length $\gamma^+$ gives $\theta$ with $\cf(\theta)=\gamma^+$ closed under $F$. Cover exactingness at $\theta$ yields $Y\ni y$; since $|Y|<\cf(\theta)$, $Y\in V_\beta$ for some $\beta<\theta$, so $b(Y)<\theta$, and Lemma~\ref{lem:descent}(b) gives a witness at $b(Y)$. ($\Leftarrow$) Use $Z=Y\cap V_\alpha\in\OD_{\{Y\}}$ and Lemma~\ref{lem:transfer}.
\end{proof}

\begin{lemma}[Exact hulls]\label{lem:hull}\src{R4, R6}
Assume $\CEx_\gamma(\lambda)$. For every sufficiently large limit $\alpha$ and finite $A\subseteq V_\alpha$ there is $\sigma$ with $A\subseteq\sigma\prec V_\alpha$, $|\sigma|\le\gamma$, and $\REx(\lambda,\sigma)$.
\end{lemma}

\begin{proof}
Code a Skolem system for $V_\alpha$ (with constants for $A$) as one function $f$, cover $f$ by a persistent $Y$ (Lemma~\ref{lem:persist}), and let $\sigma$ be the closure of $\omega$ under \emph{all} functions $V_\alpha^{<\omega}\to V_\alpha$ that belong to $Y$. Then $|\sigma|\le\gamma$, $\sigma\prec V_\alpha$, and $\sigma$ is definable from $Y$ and $\alpha$ alone, so Lemma~\ref{lem:transfer} applies. The individually chosen $f$ is not a parameter of the definition.
\end{proof}

\begin{lemma}[Trace reconstruction]\label{lem:trace}\src{R1--R4, R6--R8}
Let $U$ be an $\eta$-complete ultrafilter on an ordinal $\tau$, and let $U,\tau\in\sigma\prec V_\alpha$ with $|\sigma|<\eta$ and $\alpha$ large. Then $U\in\OD_{\{\sigma\}}$. Hence for $x\in\CD(\eta)$ with $\eta$ regular uncountable, there is a club $C\subseteq P_\eta(V_\alpha)$ with $x\in\OD_{\{\sigma\}}$ for all $\sigma\in C$.
\end{lemma}

\begin{proof}
Let $\beta=\min\bigcap(U\cap\sigma)$, which exists by completeness. For $A\in\sigma\cap\Pow(\tau)$ we get $A\in U\iff\beta\in A$, using that $\tau\setminus A\in\sigma$. $U$ is the \emph{unique} ultrafilter $W\in\sigma$ on $\tau$ with this trace: two such candidates would be separated by a set in $\sigma$, by elementarity. The definition uses only $\sigma$ and the ordinals $\tau,\beta,\alpha$. The seed $\beta$ need not belong to $\sigma$, $U$ need not be principal, and uniqueness is among candidates \emph{inside} $\sigma$ only.
\end{proof}

The mechanism is the one in the proof of Theorem~3.7 of the Blue--Goldberg notes; the reports' addition is the exact inequality $|\sigma|\le\gamma<\eta$.

\begin{lemma}[The singular endpoint]\label{lem:singular}\src{R6}
If $\gamma$ is singular, every $\gamma$-complete filter is $\gamma^+$-complete; hence $\CD(\gamma)=\CD(\gamma^+)$ and $\HCD(\gamma)=\HCD(\gamma^+)$.
\end{lemma}

\begin{proof}
Split an index set of size $\gamma$ into $\cf(\gamma)$ pieces of size below $\gamma$ and intersect in two stages.
\end{proof}

\section{The completeness barrier}\label{sec:barrier}

\begin{theorem}[Completeness barrier]\label{thm:barrier}\src{R1--R4, R6--R8}
Assume $\CEx_\gamma(\lambda)$ and put $\rho=q(\gamma)$.
\begin{enumerate}
\item \emph{Absorption.} $\REx(\lambda,x)$ for every $x\in\CD(\rho)$.\src{R4, R6}
\item No cofinal $a\subseteq\lambda$ with $\ot(a)<\lambda$ (equivalently $|a|<\lambda$) belongs to $\CD(\rho)$.
\item $\cf^V(\lambda)=\omega$, while $\cf^{\HCD(\eta)}(\lambda)=\lambda$ for every cardinal $\eta\ge\rho$.
\item Every cofinal $b\subseteq\lambda$ in $\HCD(\rho)$, or in any inner model contained in it, has $\ot(b)=\lambda=|b|^V$. Thus a countable cofinal $c\subseteq\lambda$ satisfies $\min\{|a|^V:a\in\HCD(\rho),\,c\subseteq a\subseteq\lambda\}=\lambda$.\src{R2, R4, R6}
\end{enumerate}
In particular $\CEx_\lambda(\lambda)$ implies that $\lambda$ is regular in $\HCD(\lambda)$.\src{R6}
\end{theorem}

\begin{proof}
By Lemma~\ref{lem:singular} we may take $\rho=\gamma^+$. (1) Let $x$ be definable from $\gamma^+$-complete $U_0,\dots,U_{m-1}$ and ordinals. By Lemma~\ref{lem:hull} (Route~H), or by stationarity (Route~S), choose $\sigma\prec V_\alpha$ containing the $U_i$ with $|\sigma|\le\gamma$ and $\REx(\lambda,\sigma)$. Lemma~\ref{lem:trace} gives $x\in\OD_{\{\sigma\}}$, and Lemma~\ref{lem:transfer} gives $\REx(\lambda,x)$. (2) follows from (1) and Lemma~\ref{lem:short}. (3) $\lambda$ remains a cardinal in $N=\HCD(\eta)$; a cofinal $f:\mu\to\lambda$ in $N$ with $\mu<\lambda$ has range of $V$-cardinality below $\lambda$, hence of order type below $\lambda$, and this range lies in $\CD(\rho)$ by \eqref{eq:mono}. (4) The order type of $b$ is at most $\lambda$ and cannot be smaller.
\end{proof}

\begin{remark}[Threshold]
\R1, \R3 and \R8 state the barrier for all $\eta>\gamma$, \R2, \R4 and \R7 at $\gamma^+$, and \R6 at $q(\gamma)$; by \eqref{eq:mono} and Lemma~\ref{lem:singular} the last is the strongest. The successor cannot be dropped for regular $\gamma$ by this proof: $\sigma$ may have size exactly $\gamma$, and a uniform $\gamma$-complete ultrafilter on $\gamma$ contains all $\gamma\setminus\{\xi\}$, whose intersection is empty \R1. No report claims the threshold is optimal.
\end{remark}

\begin{corollary}[What escapes the high-completeness models]\label{cor:escape}\src{R3, R8}
Assume $\CEx_\gamma(\lambda)$ and $\eta\ge q(\gamma)$. No strictly increasing cofinal $c:\omega\to\lambda$ is ordinal definable from an $\eta$-complete ultrafilter on an ordinal, however large its underlying ordinal. No nontrivial elementary $e:V_\lambda\to V_\lambda$ belongs to $\CD(\eta)$.
\end{corollary}

\begin{proof}
$\ran(c)$ is a short cofinal set. From $e$ one defines its critical sequence, which is cofinal by Lemma~\ref{lem:crit}.
\end{proof}

\begin{corollary}[Certificate of failed cover exactingness]\label{cor:cert}\src{R1}
If $\lambda\le\gamma<\eta$ and some cofinal $a\subseteq\lambda$ of size below $\lambda$ lies in $\CD(\eta)$, then for all large $\alpha$ a club of elementary hulls in $P_{\gamma^+}(V_\alpha)$ is disjoint from $\{\sigma:\REx(\lambda,\sigma)\}$; so $\CEx_\gamma(\lambda)$ fails.
\end{corollary}

\begin{corollary}[Lower covering and approximation fail]\label{cor:approx}\src{R1--R4, R6, R8}
Assume $\CEx_\gamma(\lambda)$, let $W\models\ZF$ be an inner model with $W\subseteq\HCD(q(\gamma))$, and let $c\subseteq\lambda$ be cofinal of order type $\omega$. Then $c\notin W$, but $c\cap b$ is finite and belongs to $W$ whenever $b\in W$ is well-orderable in $W$ with $|b|^W<\lambda$. Consequently $(W,V)$ fails the $\theta$-cover and $\theta$-approximation size tests for every uncountable cardinal $\theta\le\lambda$.
\end{corollary}

\begin{proof}
$\lambda$ is regular in $W$, so $b\cap\lambda$ is bounded, and a cofinal $\omega$-sequence has finitely many terms below a bound.
\end{proof}

This is compatible with Input~\ref{in:hcd}: when $W=\HCD(\kappa)$ for a strongly compact $\kappa>\gamma$, approximation and covering hold at $\kappa$, where the test set $b=\lambda$ is allowed and $c\cap\lambda=c\notin W$.

\section{A strongly compact cardinal below \texorpdfstring{$\lambda$}{lambda}}\label{sec:below}

\begin{theorem}[Small-witness separation]\label{thm:sep}\src{R1--R4, R6--R8}
Suppose $\delta<\lambda\le\gamma$, $\SC(\delta)$ and $\CEx_\gamma(\lambda)$; let $\rho=q(\gamma)$. There is a cofinal $a\subseteq\lambda$ with
\[
a\in\HCD(\delta),\qquad|a|^{\HCD(\delta)}<\delta,\qquad\ot(a)<\delta,\qquad a\notin\textstyle\bigcup_{\eta\ge\rho}\CD(\eta).
\]
Hence $\cf^{\HCD(\delta)}(\lambda)<\delta<\lambda=\cf^{\HCD(\eta)}(\lambda)$ for all $\eta\ge\rho$, and $\Pow(\lambda)\cap\HCD(\rho)\subsetneq\Pow(\lambda)\cap\HCD(\delta)$. Moreover:
\begin{enumerate}
\item every $b\in\HCD(\rho)$ with $a\subseteq b$ has $|b|^V\ge\lambda$;\src{R2}
\item no cofinal $b\subseteq a$ belongs to $\HCD(\rho)$;\src{R6}
\item some $\delta$-complete ultrafilter $U$ on an ordinal satisfies $U\notin\CD(\rho)$.\src{R4}
\end{enumerate}
\end{theorem}

\begin{proof}
Cover a countable cofinal $c\subseteq\lambda$ by $b_0\in\HCD(\delta)$ of internal size below $\delta$ (Input~\ref{in:hcd}) and let $a=b_0\cap\lambda$. Its order type is below the cardinal $\delta$ and its enumeration lies in $\HCD(\delta)$. Theorem~\ref{thm:barrier}(2),(4) give the rest; for (3), take $U$ with $a\in\OD_{\{U\}}$ and compose definitions.
\end{proof}

The witness $a$ need not be countable, even in $V$: strong compactness supplies a small \emph{cover} of an $\omega$-sequence, not the sequence. The forcing-size bound of Input~\ref{in:hcd} gives an independent proof of the weaker conclusion $\cf^{\HCD(\delta)}(\lambda)<\lambda$, since a forcing of size below the strong limit $\lambda$ cannot singularize a regular $\lambda$ \R1.

\subsection{The conditional inconsistency, in all its reported forms}

Each report names a slightly different additional principle. For $\delta<\lambda\le\gamma$ and $\rho=q(\gamma)$:

\begin{center}
\small
\renewcommand{\arraystretch}{1.2}
\begin{tabularx}{\textwidth}{@{}lYl@{}}
\toprule
&\textbf{Principle}&\textbf{Source}\\\midrule
$\mathsf P_1$&$\Pow(\lambda)\cap\HCD(\delta)=\Pow(\lambda)\cap\HCD(\rho)$ \ (local stabilization)&\R3, \R4, \R6, \R8\\
$\mathsf P_2$&$[\lambda]^{<\delta}\cap\HCD(\delta)\subseteq\HCD(\rho)$, sizes ambient \ (local promotion $\mathsf{LP}$)&\R1, \R2, \R3\\
$\mathsf P_3$&as $\mathsf P_2$ with size computed in $\HCD(\delta)$ \ ($\mathsf{LS}$)&\R7\\
$\mathsf P_4$&every $x\in[\lambda]^{<\delta}\cap\HCD(\delta)$ is covered by some $b\in\HCD(\rho)$ with $|b|^V<\lambda$ \ ($\mathsf{Cov}$)&\R2\\
$\mathsf P_5$&every cofinal $x\in\HCD(\delta)$ with $\ot(x)<\delta$ has a cofinal subset in $\HCD(\rho)$ \ ($\mathsf{CR}$)&\R6\\
$\mathsf P_6$&every $\delta$-complete ultrafilter on an ordinal lies in $\CD(\rho)$&\R4\\
$\mathsf P_7$&some cofinal $x\subseteq\lambda$ with $|x|<\lambda$ lies in $\CD(\rho)$&\R1, \R4, \R7\\
\bottomrule
\end{tabularx}
\end{center}

\begin{corollary}[Conditional inconsistency below $\lambda$]\label{cor:cond}\src{R1--R4, R6--R8}
$\ZFC$ refutes $\SC(\delta)\wedge\delta<\lambda\le\gamma\wedge\CEx_\gamma(\lambda)\wedge\mathsf P_i$ for each $i\le6$, and refutes $\CEx_\gamma(\lambda)\wedge\mathsf P_7$ outright. The implications are
$\mathsf P_6\Rightarrow\mathsf P_1\Rightarrow\mathsf P_2\Leftrightarrow\mathsf P_3$, $\ \mathsf P_2\Rightarrow\mathsf P_4$, $\ \mathsf P_2\Rightarrow\mathsf P_5$, and under the hypotheses each of $\mathsf P_4,\mathsf P_5$ yields $\mathsf P_7$.
\end{corollary}

\begin{proof}
Apply each principle to the set $a$ of Theorem~\ref{thm:sep}. For $\mathsf P_2\Leftrightarrow\mathsf P_3$: a subset of $\lambda$ in $\HCD(\delta)$ of ambient size below $\delta$ has order type below $\delta$, so its enumeration in $\HCD(\delta)\models\ZFC$ makes it internally small; the converse is trivial. $\mathsf P_6\Rightarrow\mathsf P_1$ holds for sets of ordinals by composing definitions.
\end{proof}

None of the principles is derived from strong compactness. The equivalence $\mathsf P_2\Leftrightarrow\mathsf P_3$ is an observation of this synthesis; it shows that the reports' differing size conventions do not produce different theorems.

\begin{corollary}[The extendible boundary recovered]\label{cor:ext}\src{R1--R4, R6--R8}
If $\delta<\lambda$ is extendible, then $\neg\CEx_\gamma(\lambda)$ for every $\gamma$.
\end{corollary}

\begin{proof}
$\HCD(\delta)=\HCD\subseteq\HCD(\rho)\subseteq\HCD(\delta)$, so $\mathsf P_1$ holds; extendibles are strongly compact. For $\gamma<\lambda$ use Lemma~\ref{lem:bound}.
\end{proof}

This is Blue--Goldberg's Theorem~3.6, not a new result. The factorization shows what each hypothesis contributes: strong compactness supplies the short cofinal set in the low model; extendibility supplies the promotion to high completeness.

\begin{theorem}[First regularization]\label{thm:first}\src{R6}
For an ambient cardinal $\lambda$ let $r(\lambda)$ be the least infinite cardinal $\eta$ with $\cf^{\HCD(\eta)}(\lambda)=\lambda$, if one exists. The map $\eta\mapsto\cf^{\HCD(\eta)}(\lambda)$ is nondecreasing. If $\CEx_\gamma(\lambda)$ then $r(\lambda)\le q(\gamma)$; if also $\SC(\delta)$ for some $\delta<\lambda$ then $r(\lambda)>\delta$. If $\CEx_\lambda(\lambda)$ and $\lambda$ is a limit of strongly compact cardinals, then $r(\lambda)=\lambda$ exactly.
\end{theorem}

The consistency of the last hypothesis is not asserted. The theorem does not show $\HCD(\lambda)=\bigcap_{\eta<\lambda}\HCD(\eta)$, and no single short cofinal set need lie in all the lower models.

\section{A strongly compact cardinal above \texorpdfstring{$\gamma$}{gamma}: grounds}\label{sec:above}

\begin{lemma}[Singularizing forcing is not $\lambda$-cc]\label{lem:cc}\src{R3, R4, R6, R8; size form R1, R2, R7}
Let $W\models\ZFC$ be an inner model in which $\lambda$ is regular uncountable, and let $V=W[G]$ for $\PP\in W$ with $\cf^V(\lambda)=\omega$ and $\lambda$ a cardinal of $V$. Then $\PP$ is not $\lambda$-cc in $W$; it has an antichain of $W$-size $\lambda$, no dense subset of $W$-size below $\lambda$, and $|\PP|^V\ge\lambda$. Without the regularity assumption: if $|\PP|^W\le\nu<\lambda$ for an infinite $\nu$, then $W$ contains a cofinal $a\subseteq\lambda$ with $|a|^W\le\nu$.
\end{lemma}

\begin{proof}
If $\PP$ were $\lambda$-cc, the sets $B_n$ of values of a name $\dot c(n)$ decided by maximal antichains have $W$-size below $\lambda$, so $\bigcup_nB_n$ is bounded in $\lambda$ by regularity in $W$, contradicting forced cofinality. For the last part take all values forced by some $q\le p$.
\end{proof}

\begin{theorem}[Proper canonical ground]\label{thm:ground}\src{R1--R4, R6, R8}
Suppose $\lambda\le\gamma<\kappa$, $\CEx_\gamma(\lambda)$ and $\SC(\kappa)$. Then $W=\HCD(\kappa)$ is a \emph{proper} ground of $V$ with
\[
W\models\text{``$\lambda$ is strongly inaccessible''},\qquad\cf^V(\lambda)=\omega.
\]
Every $\PP\in W$ with $V=W[G]$ satisfies the conclusions of Lemma~\ref{lem:cc}, and $(W,V)$ fails $\theta$-covering and $\theta$-approximation for every uncountable $\theta\le\lambda$ while satisfying both at $\kappa$.
\end{theorem}

\begin{proof}
$\kappa\ge\gamma^+\ge q(\gamma)$, so $W\subseteq\HCD(q(\gamma))$ and $\lambda$ is regular in $W$ by Theorem~\ref{thm:barrier}; $W$ is a ground by Input~\ref{in:hcd}. For the strong-limit clause \R8: an injection $\lambda\to\Pow^W(\mu)$ in $W$ with $\mu<\lambda$ would contradict $2^\mu<\lambda$ in $V$. The rest is Lemma~\ref{lem:cc} and Corollary~\ref{cor:approx}.
\end{proof}

\begin{corollary}[Ground Axiom obstruction]\label{cor:ga}\src{R1--R4, R6, R8}
The theory $\ZFC+\GA+\exists\lambda\le\gamma<\kappa\,[\CEx_\gamma(\lambda)\wedge\SC(\kappa)]$ is inconsistent. Under $\GA$, every strongly compact cardinal is at most every cover bound of every cover-exacting cardinal; and $\ZFC+\GA+{}$``there is a proper class of strongly compact cardinals'' proves that no cardinal is cover exacting for any set-sized bound. No outer model of $\ZFC+\GA$ can preserve both $\CEx_\gamma(\lambda)$ and a strongly compact $\kappa>\gamma$.
\end{corollary}

The inequality $\kappa>\gamma$, not merely $\kappa>\lambda$, is essential: for $\lambda<\kappa\le\gamma$ monotonicity does not place $\HCD(\kappa)$ inside $\HCD(\gamma^+)$ \R4. Nothing here excludes $\GA+\CEx_\gamma(\lambda)$ with a single strongly compact below $\lambda$: then $V=\HCD(\delta)$ and the gap of Section~\ref{sec:below} remains \R3. The theorem is consonant with the known positive constructions, which are Prikry extensions.

\begin{corollary}[General ground obstruction]\label{cor:gengound}\src{R1, R3, R7}
Assume $\CEx_\gamma(\lambda)$. If $W\models\ZFC$ is an inner model with $W\subseteq\HCD(q(\gamma))$ and $V=W[G]$, then the forcing is not $\lambda$-cc in $W$. If instead $V=W[G]$ for an arbitrary ground $W$ with $|\PP|^W\le\nu<\lambda$, then $W$ contains a cofinal $a\subseteq\lambda$ with $|a|^W\le\nu$ and $a\notin\CD(q(\gamma))$; so the assumption that all such small sets of $W$ lie in $\CD(q(\gamma))$ is contradictory.
\end{corollary}

\begin{corollary}[Two strongly compacts]\label{cor:sandwich}\src{R1, R6, R8}
If $\delta<\lambda\le\gamma<\kappa$ with $\SC(\delta)$, $\SC(\kappa)$, $\CEx_\gamma(\lambda)$, then both $\HCD(\kappa)$ and $\HCD(\delta)$ are grounds and
\[
\underbrace{\HCD(\kappa)}_{\lambda\text{ inaccessible}}\subseteq\underbrace{\HCD(q(\gamma))}_{\lambda\text{ regular}}\subsetneq\underbrace{\HCD(\delta)}_{\cf(\lambda)<\delta}\subseteq V,
\]
the strict inclusion being witnessed on $[\lambda]^{<\delta}$. $\HCD(\delta)$ has $\delta$-approximation while $\HCD(\kappa)$ fails it. Nothing prevents $\HCD(\delta)=V$.
\end{corollary}

\emph{Third edition.} Theorem~\ref{thm:sandwich} below shows that the hypotheses of Corollary~\ref{cor:sandwich} are inconsistent. The corollary, and Corollary~\ref{cor:statapps}(b), remain correct as conditional statements and record exactly the structure from which \R{24} extracts the contradiction.


\subsection{A successor stationary-set obstruction (second round)}\label{sec:stat}

Lemma~\ref{lem:cc} bounds a singularizing forcing at $\lambda$. With a strongly compact cardinal \emph{below} $\lambda$ the bound moves to the successor, and a specific stationary set must be lost. The mechanism (non-reflection after Prikry-type singularization versus reflection from strong compactness) is classical; \R{14} frees it from any assumption that $V$ is a forcing extension.

\begin{lemma}[Reflection]\label{lem:screfl}\src{R14; classical}
If $\delta$ is strongly compact, $\theta>\delta$ is regular and $S\subseteq E^\theta_\omega$ is stationary, then $S\cap\alpha$ is stationary in some $\alpha<\theta$ with $\omega<\cf(\alpha)<\delta$.
\end{lemma}

\begin{proof}
Let $j:V\to M$ come from a fine $\delta$-complete ultrafilter on $\Pow_\delta(\theta)$ and $\beta=\sup j``\theta<j(\theta)$; then $\theta\le\cf^M(\beta)<j(\delta)$. Given a club $C\subseteq\beta$ in $M$, interleave $C$ with $j``\theta$ and take $\alpha\in S$ closed under the interleaving function; $\cf(\alpha)=\omega$ gives $j(\alpha)=\sup j``\alpha\in C\cap j(S)$. So $M$ thinks $j(S)$ reflects at $\beta$; apply elementarity.
\end{proof}

\begin{theorem}[Collapse or destroy]\label{thm:successor}\src{R14}
Let $W\subseteq V$ be an inner model of $\ZFC$, $\lambda$ a cardinal of $V$ that is regular in $W$ with $\cf^V(\lambda)=\omega$, and let some $\delta<\lambda$ be strongly compact in $V$. Put $\theta=(\lambda^+)^W$ and $S=(E^\theta_\lambda)^W$. Then either $|\theta|^V=\lambda$, or $\theta=(\lambda^+)^V$ and the $W$-stationary set $S$ is nonstationary in $V$.
\end{theorem}

\begin{proof}
Assume $\theta=(\lambda^+)^V$. $S$ is stationary in $W$ and $S\subseteq(E^\theta_\omega)^V$. If $\cf^V(\alpha)>\omega$ then $\cf^W(\alpha)<\lambda$, and a continuous cofinal sequence of length $\cf^W(\alpha)$ in $W$ is a club of $\alpha$ avoiding $S$; so $S$ reflects nowhere in $V$, and Lemma~\ref{lem:screfl} makes it nonstationary.
\end{proof}

\begin{corollary}[Successor chain-condition bound]\label{cor:succcc}\src{R14}
In the situation of Theorem~\ref{thm:successor}, if $V=W[G]$ for $\PP\in W$ then $\PP$ is not $(\lambda^+)^W$-cc in $W$; every dense subset of $\PP$ in $W$ has $W$-size at least $(\lambda^+)^W$. (A $\theta$-cc forcing preserves $\theta$ and every stationary subset of $\theta$.)
\end{corollary}

\begin{corollary}\label{cor:statapps}\src{R14}
(a) If $\lambda$ is exacting and some $\delta<\lambda$ is strongly compact, then $(\lambda^+)^{\HOD}$ is collapsed or $(E^{(\lambda^+)^{\HOD}}_\lambda)^{\HOD}$ is nonstationary in $V$.
(b) If $\delta<\lambda\le\gamma<\kappa$ with $\SC(\delta)$, $\SC(\kappa)$ and $\CEx_\gamma(\lambda)$, then every presentation of $V$ over $W=\HCD(\kappa)$ has an antichain of $W$-size $(\lambda^+)^W$, and $(\lambda^+)^W$ is collapsed or $S=(E^{(\lambda^+)^W}_\lambda)^W$ is nonstationary in $V$. If $(\lambda^+)^W=(\lambda^+)^V$, then $S$ is already nonstationary in $\HCD(\delta)$: some club $C\in\HCD(\delta)\setminus W$ avoids it.
\end{corollary}

\begin{proof}
(a) $W=\HOD$ and Input~\ref{in:abl}(i). (b) Theorem~\ref{thm:ground} and Corollary~\ref{cor:succcc}. For the last sentence, $V$ is an extension of $U=\HCD(\delta)$ by a forcing of size below $(2^{2^\delta})^+<\lambda$ (Input~\ref{in:hcd}; $\lambda$ is a strong limit), which is $\theta$-cc in $U$ and so preserves stationary subsets of $\theta$; since $S$ is nonstationary in $V$, it is nonstationary in $U$, and a club of $U$ avoiding $S$ cannot lie in $W$.
\end{proof}

Part (b) strengthens the $\lambda$-cc bound of Theorem~\ref{thm:ground} to $\lambda^+$ when there is also a strongly compact cardinal below $\lambda$. It is an obstruction of the same conditional kind: it is incompatible with a $(\lambda^+)$-cc presentation, or with preservation of one named stationary set, not with the configuration itself.

\subsection{No strongly compact sandwich (third round)}\label{sec:sandwich}

For $\cf^V(\alpha)>\omega$ let $\mathcal F^V_{\alpha,\omega}$ be the filter generated by the $\omega$-closed unbounded subsets of $\alpha$. An inner model $N\models\ZFC$ is \emph{$\omega$-club amenable} if $\mathcal F^V_{\alpha,\omega}\cap N\in N$ for all such $\alpha$ (Goldberg \cite[\S2.2]{dichotomy}). The filter is the ambient one; amenability does not make ambient stationarity internal, it only encodes, for sets already in $N$, the test of ambient positivity by a set of $N$.

\begin{lemma}[Amenability]\label{lem:amenable}\src{R24}
If $A$ is ordinal definable and $A\subseteq\HCD(\eta)$ then $A\in\HCD(\eta)$. Hence every $\HCD(\eta)$ that satisfies Choice, and $\HOD$, is $\omega$-club amenable.
\end{lemma}

\begin{theorem}[Covering by an amenable ground]\label{thm:groundcover}\src{R24}
Let $W\models\ZFC$ be a ground of $V$ that is definable from a set parameter and $\omega$-club amenable, and let $\delta$ be strongly compact in $V$. Then $W$ has the $\delta$-cover property in $V$. Consequently, if $\lambda>\delta$ is regular in $W$, then $\cf^V(\lambda)\ge\delta$.
\end{theorem}

\begin{proof}
Let $x\subseteq\mu$ with $|x|<\delta$. Choose $\theta>\max(\mu,\delta)$ a successor cardinal of $W$ above the size of the forcing; in $W$ split $(E^\theta_\omega)^W$ into $\mu$ disjoint stationary sets $A_\xi$; they remain stationary in $V$ because the forcing is $\theta$-cc. Let $j:V\to M$ be the ultrapower by a fine $\delta$-complete ultrafilter on $\Pow_\delta(\theta)$ and $\beta=\sup j``\theta$, so $\omega<\cf^M(\beta)<j(\delta)$. By amenability transferred to $j(W)$, the set $s=\{\xi<j(\mu):j(\vec A)_\xi\cap\beta\text{ is }\mathcal F^M_{\beta,\omega}\text{-positive}\}$ belongs to $j(W)$. It contains $j``\mu$ (each $j(A_\xi\cap E^\theta_\omega)\cap\beta$ meets every club of $M$, as in Lemma~\ref{lem:screfl}) and has $M$-size below $j(\delta)$ (disjoint positive sets meet a club of size $\cf^M(\beta)$ in distinct points). So $s\in j(W\cap\Pow_\delta(\mu))$, the ultrafilter derived from $s$ is fine, $\delta$-complete and concentrates on $W\cap\Pow_\delta(\mu)$, and a member of it covers $x$. For the last sentence cover a cofinal subset of $\lambda$ of size $\cf^V(\lambda)<\delta$.
\end{proof}

This is Goldberg's stationary-seed argument with the compactness cardinal and the inner model decoupled. \R{24} also verifies the consequence from a published statement: by \cite[Lemma~2.7]{dichotomy}, if $N$ is $\omega$-club amenable and $\delta$ is $\omega_1$-strongly compact, then all sufficiently large regular cardinals are measurable in $N$ or $N$ covers countable sets by sets of size below $\delta$; and the first alternative fails for a ground, because arbitrarily large successor cardinals of $N$ remain regular.

\begin{theorem}[No strongly compact sandwich]\label{thm:sandwich}\src{R24}
In $\ZFC$, suppose $\delta<\lambda\le\gamma$, $\SC(\delta)$ and $\CEx_\gamma(\lambda)$. Then no $\HCD(\eta)$ with $\eta\ge q(\gamma)$ is a ground of $V$. In particular there is no strongly compact cardinal above $\gamma$.
\end{theorem}

\begin{proof}
A ground satisfies Choice, so by Lemma~\ref{lem:amenable} and Theorem~\ref{thm:groundcover} it would have the $\delta$-cover property; but $\lambda$ is regular in $\HCD(\eta)$ (Theorem~\ref{thm:barrier}) while $\cf^V(\lambda)=\omega<\delta$. If $\kappa>\gamma$ were strongly compact, $\HCD(\kappa)$ would be a ground (Input~\ref{in:hcd}).
\end{proof}

\begin{corollary}\label{cor:sandwichcons}\src{R24}
(a) If a cover-exacting $\lambda$ lies above a strongly compact cardinal, the strongly compact cardinals are bounded by every cover bound of $\lambda$; with a proper class of strongly compact cardinals, every cover-exacting cardinal lies below the least of them. (b) If $\lambda$ is exacting and some $\delta<\lambda$ is strongly compact, then $\HOD$ is not a ground of $V$. (c) The same holds with ``$\delta$ strongly compact'' weakened to: for all large regular $\theta$ there is a countably complete fine ultrafilter on $\Pow_\delta(\theta)$. (d) Without any upper strongly compact: if $\SC(\delta)$, $\delta<\lambda\le\gamma$, $\CEx_\gamma(\lambda)$ and $N=\HCD(\eta)\models\ZFC$ for some $\eta\ge q(\gamma)$, then for every $V$-regular $\theta>\lambda$ the filter $\mathcal F^V_{\theta,\omega}\cap N$ admits no $\lambda$ disjoint positive sets in $N$ (so fewer than $\lambda$ cells of any $N$-partition of $\theta$ remain stationary in $V$), and every $V$-regular $\theta>(2^\lambda)^N$ is measurable in $N$.
\end{corollary}

Compare Corollary~\ref{cor:ga}: there the Ground Axiom was needed to turn the proper ground $\HCD(\kappa)$ into a contradiction; a strongly compact cardinal \emph{below} $\lambda$ does the same. This is a partial negative answer to Problem~5.2. It is consistent with the known models: after Prikry forcing at $\kappa$ there is no strongly compact cardinal below $\kappa$.

\begin{theorem}[Cofinality cascade]\label{thm:cascade}\src{R20}
Let $V=W[G]$ for $\PP\in W$, let $\delta<\lambda$ be strongly compact in $V$, and let $\lambda$ be a cardinal of $V$, regular in $W$, with $\cf^V(\lambda)<\delta$. Then there are $\delta$ many $W$-regular cardinals $\mu_i>\lambda$, increasing, whose $V$-cofinalities $\tau_i<\delta$ are strictly increasing and cofinal in $\delta$; they can be found below any $W$-cardinal bounding the density of $\PP$. Hence $\PP$ has in $W$ an antichain of size $\Theta=(\lambda^{+\delta})^W$ and no dense subset of smaller $W$-size. In particular $\lambda$ cannot be singularized while all larger ground-regular cardinals stay regular, nor with a bound below $\delta$ on the new cofinalities.
\end{theorem}

\begin{proof}
If $\PP$ were $\chi$-cc for a regular $\chi\le\Theta$ (for singular $\Theta$ use Erd\H os--Tarski), all strata $(E^\theta_\mu)^W$ at a suitable $\theta$ would stay stationary. $(E^\theta_\lambda)^W\subseteq(E^\theta_{<\delta})^V$ reflects (Lemma~\ref{lem:screfl}, for fewer than $\delta$ sets simultaneously) at some $\alpha$ with $\cf^V(\alpha)<\delta$, and then $\mu_0=\cf^W(\alpha)>\lambda$ with $\cf^V(\mu_0)>\cf^V(\lambda)$; iterate, reflecting all earlier strata simultaneously.
\end{proof}

Theorem~\ref{thm:cascade} is a statement about arbitrary grounds and is not affected by Theorem~\ref{thm:sandwich}; its applications to $W=\HOD$ and $W=\HCD(\kappa)$ in \R{20} are, by Theorem~\ref{thm:sandwich} and Corollary~\ref{cor:sandwichcons}(b), statements about an impossible situation.

\section{Exacting cardinals: definable families of cofinal maps}\label{sec:width}

This section uses only Input~\ref{in:abl}(i). Write $N=\HOD_{V_\lambda}$. An $\OD_{V_\lambda}$ set need not belong to $N$, but an $\OD_{V_\lambda}$ set of ordinals, or graph between ordinals, does.

\begin{theorem}[Projection width]\label{thm:width}\src{R9}
Let $\lambda$ be exacting, $0<\mu<\lambda$, and let $F\in\OD_{V_\lambda}$ be a nonempty set of cofinal maps $\mu\to\lambda$. Then for some $\xi<\mu$ the projection $P_\xi=\{f(\xi):f\in F\}$ is unbounded in $\lambda$, and for every such $\xi$, $\ot(P_\xi)=\lambda=|P_\xi|^V$; hence $|F|^V\ge\lambda$. If all $f\in F$ are increasing this holds from some coordinate onward.
\end{theorem}

\begin{proof}
All $P_\xi$ are defined uniformly from one definition of $F$, so $\xi\mapsto\sup\{\beta+1:\beta\in P_\xi\}$ lies in $N$. If every $P_\xi$ were bounded, regularity of $\lambda$ in $N$ would bound this function below $\lambda$, and then every $f\in F$ would be bounded. An unbounded $P_\xi\in N$ has order type $\lambda$, again by regularity.
\end{proof}

No member of $F$ is assumed definable and no choice is made inside $N$.

\begin{corollary}\label{cor:odcover}\src{R9}
If $\lambda$ is exacting, there is no $\OD_{V_\lambda}$ family $F$ with $0<|F|<\lambda$ of cofinal maps from a fixed $\mu<\lambda$ into $\lambda$. If $c:\omega\to\lambda$ is cofinal, every $A\in\OD_{V_\lambda}$ with $c\in A$ has $|A|\ge\lambda$; thus the principle ``every $x\in V_{\lambda+1}$ lies in an $\OD_{V_\lambda}$ set of size below $\lambda$'' fails.
\end{corollary}

\begin{proposition}[Fresh cofinal set]\label{prop:fresh}\src{R9; HOD form R2, R3, R8}
If $\lambda$ is exacting and $c\subseteq\lambda$ is cofinal of order type $\omega$, then $c\notin\HOD_{V_\lambda}$, but $c\cap a$ is finite and in $\HOD_{V_\lambda}$ for every $a\in\HOD_{V_\lambda}$ that injects there into an ordinal below $\lambda$.
\end{proposition}

The cloud of Section~\ref{sec:i3} does not contradict Corollary~\ref{cor:odcover}: it may have size exactly $\lambda$, and its definition uses the generic sequence and a name, not parameters in $V_\lambda$.

\subsection{Small families of short cofinal sets, bounded separation, and laws (second round)}\label{sec:smallfam}

Write $\Clam$ for the set of cofinal $a\subseteq\lambda$ with $|a|<\lambda$ (equivalently $\ot(a)<\lambda$); so $D_\lambda\subseteq\Clam$ in the notation of Section~\ref{sec:uex}.

\begin{lemma}[Fixed small sets]\label{lem:fixedsmall}\src{R10, R12, R14--R18}
Let $j:X\to V_\zeta$ be an exacting witness at $\lambda$ with $\kappa=\crit(j)$.
(a) If $S\in X$, $S\subseteq\lambda$, $j(S)=S$ and $|S|<\lambda$, then $S\subseteq\kappa$.
(b) There is no $\FF\in X$ with $j(\FF)=\FF$, $\varnothing\ne\FF\subseteq\Clam$ and $|\FF|<\lambda$. The members of $\FF$ may have different, unbounded, order types and need not be fixed.
\end{lemma}

\begin{proof}
(a) $\mu=|S|$ is a fixed ordinal below $\lambda$, so $\mu<\kappa$; a bijection $f:\mu\to S$ in $X$ gives $j``S=\ran j(f)=S$, and an order-preserving bijection of a set of ordinals onto itself is the identity. So $S$ consists of fixed ordinals, which lie below $\kappa$ by Lemma~\ref{lem:crit}. (b) Let $\tau$ be the least order type of a member of $\FF$ and $u=\bigcup\{b\in\FF:\ot(b)=\tau\}$. Then $u$ is cofinal, $|u|\le|\FF|\cdot|\tau|<\lambda$, and $u$ is definable from $\FF$, hence fixed, contradicting (a).
\end{proof}

The compression to the minimal order type \R{11} is what removes any uniform bound: a plain union could have size $\lambda$ when the sizes of the members are unbounded. For subfamilies of $D_\lambda$ the plain union \R{13}, \R{16} or the coordinate projections \R{12}, \R{14} serve equally.

\begin{theorem}[No small definable family of short cofinal sets]\label{thm:smallC}\src{R10; parameter device R18}
If $\lambda$ is exacting, every nonempty $A\in\OD_{V_\lambda}$ with $A\subseteq\Clam$ has $|A|\ge\lambda$.
\end{theorem}

\begin{proof}
Otherwise let $\rho$ be the least rank of a parameter $p\in V_\lambda$ for which some counterexample lies in $\OD_p$. At a sufficiently correct height, $\rho$ is definable from $\lambda$, hence fixed by any witness $j$, hence $\rho<\crit(j)$; so such a $p$ lies in $V_{\crit(j)}$ and is fixed. The least counterexample in the canonical well-order of $\OD_p$ is then fixed, contradicting Lemma~\ref{lem:fixedsmall}(b).
\end{proof}

The \emph{minimal-rank-parameter device} of \R{18} is worth isolating: no ordinal or low-rank parameter of the original definition is assumed fixed, and no witness with a prescribed high critical point is needed. \R{10} and \R{14}--\R{17} instead use Input~\ref{in:abl}(ii) to raise the critical point above the rank of a given parameter; both routes are correct. Theorem~\ref{thm:smallC} contains Input~\ref{in:abl}(i) (take $A=\{a\}$), Corollary~\ref{cor:odcover}, and the version of the latter for maps with varying domains below $\lambda$ \R{18}.

\begin{theorem}[Bounded-rank separation]\label{thm:pinning}\src{R10, R14, R17}
Let $\lambda$ be an uncountable cardinal with $|V_\alpha|<\lambda$ for all $\alpha<\lambda$. Then $\lambda$ is regular in $\HOD_{V_\lambda}$ iff every $A\in\OD_{V_\lambda}$ with $A\subseteq V_{\lambda+1}$ and $|A|<\lambda$ is separated at a bounded rank, i.e.\ $x\mapsto x\cap V_\alpha$ is injective on $A$ for some $\alpha<\lambda$. In that case $A\in\HOD_{V_\lambda}$, $A\subseteq\HOD_{V_\lambda}$, and $\HOD_{V_\lambda}$ computes $|A|$ correctly. In particular this holds at every exacting $\lambda$, where for $\OD_{V_\lambda}$ sets $A\subseteq V_{\lambda+1}$: $|A|<\lambda$ iff $A$ is separated at a bounded rank.
\end{theorem}

\begin{proof}
($\Rightarrow$) The set of first-difference ranks $\{d(x,y):x\ne y\in A\}$ is an $\OD_{V_\lambda}$ set of fewer than $\lambda$ ordinals, hence bounded. Each $x\in A$ is then definable from $A$ and $x\cap V_\alpha\in V_\lambda$. ($\Leftarrow$) If $b\in\HOD_{V_\lambda}$ is short cofinal, $\{b\cap\beta:\beta\in b\}$ is not separated at any bounded rank.
\end{proof}

The threshold is sharp ($A=\lambda$). The argument does not reach selectors on $[Q_\lambda]^n$, whose codes are subsets, not elements, of $V_{\lambda+1}$ \R{10}.

\begin{theorem}[Laws on cofinal sequences]\label{thm:laws}\src{R10}
Let $\cf(\lambda)=\omega$ and let $m$ be a countably additive probability measure on the $\sigma$-algebra on $S\subseteq D_\lambda$ generated by the sets $\{a\in S:a(n)<\alpha\}$. Then the medians $\min\{\alpha:m(\{a:a(n)<\alpha\})>\tfrac12\}$, $n<\omega$, form a cofinal set of order type $\omega$, definable from $S$ and $m$. Consequently, if $\lambda$ is exacting there is no nonempty $\OD_{V_\lambda}$ family of fewer than $\lambda$ such laws on $D_\lambda$.
\end{theorem}

\begin{proof}
The medians exist by continuity from below and are nondecreasing. If bounded by $\beta$, the sets $\{a:a(n)<\beta\}$ decrease to $\varnothing$ but all have measure above $\tfrac12$, contradicting continuity from above. The family of median sets contradicts Theorem~\ref{thm:smallC}.
\end{proof}

Countable additivity is essential: a nonprincipal ultrafilter on the tails of one sequence gives a finitely additive law with no medians \R{10}.

\section{Ultraexacting cardinals: failure of definable finite choice}\label{sec:uex}

For $\cf(\lambda)=\omega$ let $D_\lambda=\{a\subseteq\lambda:\ot(a)=\omega,\ \sup a=\lambda\}$, let $a=^*b$ mean $|a\triangle b|<\omega$, and $Q_\lambda=D_\lambda/{=^*}$. For $1\le r<n<\omega$ let $\Sel_{n,r}(Q)=\{f:[Q]^n\to[Q]^r:f(B)\subseteq B\}$. The even/odd splitting of a critical sequence is in the Blue--Goldberg notes (Remark~2.4); the reports extend it in two complementary ways, through two different imported interfaces.

\subsection{Finite cycles and finite families (interface: Input~\ref{in:abl}(ii))}

\begin{lemma}[Finite algebra]\label{lem:algebra}\src{R5}
Let $\tau$ be an $N$-cycle on $A$, $\sigma$ a permutation of a finite set $\FF$ with $|\FF|=m$, and $v_f:[A]^n\to[A]^r$ with $v_f(B)\subseteq B$ and $v_{\sigma(f)}(\tau``B)=\tau``v_f(B)$. (a) If $N=n$, every $\sigma$-orbit has length divisible by $n/\gcd(n,r)$, and this bound is sharp. (b) If $N=n\cdot\lcm(1,\dots,m)$, no such data exist.
\end{lemma}

\begin{proof}
(b) With $L=\lcm(1,\dots,m)$, $\sigma^L=\id$, so $v_f$ commutes with $\tau^L$, which acts on $B_0=\{q_{kL}:k<n\}$ as a single $n$-cycle; a transitive cycle has no invariant subset of size $0<r<n$.
\end{proof}

Part (a) alone does \emph{not} exclude finite families (two binary selectors can be swapped on a 2-cycle); the amplification (b) is the essential step.

\begin{lemma}[Internal cycles]\label{lem:cycles}\src{R5}
Let $j:X\to V_\zeta$ be an ultraexacting witness with $\zeta>\lambda+\omega$. Then the critical sequence $\vec\kappa$ of $e=j\restr V_\lambda$ belongs to $X$, and for every $N\ge1$ the classes $q_i=[\{\kappa_{Nk+i}:k<\omega\}]$, $i<N$, lie in $X$ and satisfy $j(q_i)=q_{i+1\bmod N}$.
\end{lemma}

\begin{proof}
$\vec\kappa$ is definable from $e\in X$. Since enumerations with domain $\omega$ are in $X$, $j(a_i)=j``a_i$, which is $a_{i+1}$, or $a_0\setminus\{\kappa_0\}$ for $i=N-1$; the quotient removes exactly this defect.
\end{proof}

\begin{lemma}[Fixed canonical witnesses]\label{lem:canonical}\src{R5}
If $\lambda$ is ultraexacting, $p\in V_\lambda$, and some $z\in\OD_p$ satisfies a first-order property $\Phi(z,p,\lambda,\bar k)$, then there are such a $z_*$ and an ultraexacting witness $j$ with $z_*\in X$ and $j(z_*)=z_*$.
\end{lemma}

\begin{proof}
Take the $\OD_p$-least witness, reflect its definition, and use Input~\ref{in:abl}(ii) to choose $\crit(j)>\rank(p)$.
\end{proof}

\begin{theorem}[No finite definable family of selectors]\label{thm:fsel}\src{R5}
If $\lambda$ is ultraexacting and $1\le r<n<\omega$, there is no nonempty finite $\FF\in\OD_{V_\lambda}$ with $\FF\subseteq\Sel_{n,r}(Q_\lambda)$. More generally, no ultraexacting witness $j$ has a nonempty finite $\FF\in X$, $\FF\subseteq\Sel_{n,r}(Q_\lambda)$, with $j(\FF)=\FF$.
\end{theorem}

\begin{proof}
For the second statement: $\FF\subseteq X$ and $j\restr\FF$ is a permutation $\sigma$; restrict the selectors to the $n\lcm(1,\dots,|\FF|)$-cycle of Lemma~\ref{lem:cycles} and apply Lemma~\ref{lem:algebra}(b). The first statement follows by Lemma~\ref{lem:canonical}. Only permutations of finite sets are iterated; no iterability of $j$ is used.
\end{proof}

The family may be definable with no definable member. Note that $\Sel_{n,r}(Q_\lambda)$ itself is ordinal definable and nonempty, yet has no nonempty finite $\OD_{V_\lambda}$ subset. \R7 proves the singleton, ordinal-parameter case by a different route (Corollary~\ref{cor:label-sel}).

\begin{corollary}\label{cor:orders}\src{R5, R7}
At an ultraexacting $\lambda$ there is no nonempty finite $\OD_{V_\lambda}$ family of linear orders of $Q_\lambda$, of tournaments on $Q_\lambda$, or of choice functions on its finite subsets; and no ordinal-definable injection of $Q_\lambda$ into the ordinals.
\end{corollary}

\begin{theorem}[Selector-coded Icarus enrichment]\label{thm:icarus}\src{R5}
Let $N\models\ZF$ be a transitive inner model with $V_{\lambda+1}\subseteq N$ and $j:N\to N$ elementary with $\crit(j)<\lambda=\sup_nj^n(\crit j)$. Then no nonempty finite $\FF\in N$, $\FF\subseteq\Sel_{n,r}(Q_\lambda)$, has $j(\FF)=\FF$. For $f\in\Sel_{n,r}(Q_\lambda)$ let $E_f\subseteq V_{\lambda+1}$ code $f$ on representatives by flat codes $\{\langle0,i\rangle\}\cup\{\langle k+1,\alpha\rangle:\alpha\in a_k\}$. Then there is no elementary $j:L(V_{\lambda+1},E_f)\to L(V_{\lambda+1},E_f)$ with $\crit(j)<\lambda$ and $j(E_f)=E_f$.
\end{theorem}

This is a forbidden \emph{information condition}, not a refutation of $\Izero^\#$ or of unspecified Icarus sets: nothing shows that a canonical enrichment defines such a selector. The quotient classes have rank $\lambda+1$, which is why representatives must be coded.

\subsection{Finite labels and transversals (interface: Input~\ref{in:abl}(iii))}

Let $D_m(\lambda)$ be the set of $m$-tuples from $D_\lambda$ with pairwise distinct classes. A rule $L:D_m(\lambda)\to k$ is coded by the predicate $A_L=\{\{\langle i,\xi\rangle:\xi\in a_i\}\cup\{\langle m,t\rangle\}:t=L(\vec a)\}\subseteq V_{\lambda+1}$, and every elementary $e$ of $(V_{\lambda+1},A_L)$ with $\crit(e)<\lambda$ satisfies $L(e(\vec a))=L(\vec a)$, using $e(a)=e``a$ for countable $a$ \R7.

\begin{lemma}[Finite-permutation realization]\label{lem:perm}\src{R7}
For $e:V_{\lambda+1}\to V_{\lambda+1}$ with $\crit(e)<\lambda$, every $m$, and every $\pi\in S_m$, there is a pairwise disjoint $\vec a\in D_m(\lambda)$ with $e(a_i)=^*a_{\pi(i)}$: on a cycle $(i_0\dots i_{\ell-1})$ with offset $t$, put $a_{i_s}=\{\kappa_{\ell n+s}+t:n<\omega\}$.
\end{lemma}

\begin{theorem}[Finite-label obstruction]\label{thm:label}\src{R7}
Let $\lambda$ be ultraexacting, $F$ a finite $S_m$-set, and $L:D_m(\lambda)\to F$ ordinal definable, invariant under finite modification of each coordinate, and $S_m$-equivariant. Then every $\pi\in S_m$ fixes some element of $F$.
\end{theorem}

\begin{proof}
With $e$ preserving $A_L$ and $\vec a$ from Lemma~\ref{lem:perm}: $L(\vec a)=L(e\vec a)=L(\pi^{-1}\cdot\vec a)=\pi^{-1}\cdot L(\vec a)$.
\end{proof}

\begin{corollary}\label{cor:label-sel}\src{R7}
At an ultraexacting $\lambda$ there is no ordinal-definable $s\in\Sel_{m,r}(Q_\lambda)$ for $1\le r<m$ (an $m$-cycle fixes no proper nonempty subset of $m$).
\end{corollary}

The test does not ask for a label fixed by the whole group: $S_3$ acting on the disjoint union of its natural 3-point action and the 2-point sign action has a fixed point for each permutation but no common fixed point. \R7 left open whether a rule $L$ exists for such a target. \R{13} settles this: the elementwise test is also \emph{sufficient} (Theorem~\ref{thm:classification}).

\begin{theorem}[No finite-valued transversal]\label{thm:transversal}\src{R7}
If $\lambda$ is ultraexacting there is no ordinal-definable $T\subseteq D_\lambda$ with $0<|T\cap[a]|<\omega$ for every $a\in D_\lambda$.
\end{theorem}

\begin{proof}
Take $e$ preserving $T$ and $a=\{\kappa_n\}$, so $e(a)=a\setminus\{\kappa_0\}$. Then $e$ injects the finite set $B=T\cap[a]$ into itself, so $e^p(b)=b$ for some $b\in B$, $p\ge1$. Then $e^p\restr b$ is an increasing bijection of a set of order type $\omega$, hence the identity, while $e^p$ moves every ordinal in $[\crit e,\lambda)$ and $b$ is cofinal.
\end{proof}

\begin{remark}[Parameters in $V_\lambda$; observation of the synthesis]
Theorem~\ref{thm:transversal} extends to $T\in\OD_{V_\lambda}$ by the method of \R5: Lemma~\ref{lem:canonical} gives a fixed $T_*$ and a witness $j$ with $a=\ran(\vec\kappa)\in X$ (Lemma~\ref{lem:cycles}); $B=T_*\cap[a]$ is a finite member of $X$ with $j(B)=T_*\cap[j(a)]=B$, so $j$ permutes $B\subseteq X$, and the same periodicity argument applies since $j(b)=j``b$ for $b\in B$.
\end{remark}

\begin{remark}[Controls]\label{rem:controls}
None of this contradicts Choice: every excluded object exists in $V$ and merely fails to be definable. Modification invariance is essential, since $D_\lambda$ itself carries an ordinal-definable lexicographic linear order \R7. The argument does not apply to ordinary exacting cardinals, whose witnesses need not contain the critical sequence \R5. And under $V_\lambda\subseteq\HOD$ the quotient $Q^{\mathrm{bd}}_\lambda$ of \emph{bounded} countable subsets of $\lambda$ modulo finite difference lies in $\HOD$ and has an ordinal-definable well-order \R5.
\end{remark}

\section{Consistency results}\label{sec:cons}

\subsection{Cover exactingness from \texorpdfstring{$\Iwf$}{I3wf(0)}}\label{sec:i3}

\begin{theorem}\label{thm:i3}\src{R9}
Let $j:V_\Lambda\to V_\Lambda$ be an $\Iwf$-embedding, $\kappa=\crit(j)$, $U_0=\{A\subseteq\kappa:\kappa\in j(A)\}$. Then $\mathbf 1_{\PP_{U_0}}\Vdash\CEx_{\check\kappa}(\check\kappa)$ for Prikry forcing $\PP_{U_0}$. Consequently there is a countable transitive model of $\ZFC+\exists\lambda\,\CEx_\lambda(\lambda)$, and
\[
\Con(\ZFC+\Iwf)\Longrightarrow\Con(\ZFC+\exists\lambda\,\CEx_\lambda(\lambda)).
\]
\end{theorem}

The proof combines three lemmas, the first two of which are reusable.

\begin{lemma}[Enumerated-domain internalization]\label{lem:internal}\src{R9}
Let $W$ be a transitive set model of enough of $\ZFC$, $e:W\to W$ an \emph{external} elementary map, $\lambda$ a cardinal of $W$, and $c=\langle\nu_n\rangle\in W$ increasing cofinal in $\lambda$ with $e(\lambda)=\lambda$, $e(\nu_n)=\nu_{n+1}$, $\omega<\nu_0$. Assume $V_\lambda\subseteq W$, $W\models|V_\lambda|=\lambda$, and $\alpha,Y\in W$ with $Y\subseteq V^W_\alpha$, $e(\alpha)=\alpha$, $e(Y)=Y$. Then $W$ contains a relative $\lambda$-exacting witness for $(V_\alpha^W,Y)$. It is not assumed that $e$ or $e\restr V_\lambda$ is in $W$.
\end{lemma}

\begin{proof}
In $W$ take $X\prec(V^W_\alpha,\in,Y)$ of size $\lambda$ containing $V_\lambda\cup\{\lambda\}$, a bijection $f:\lambda\to X$, and put $Z=e(X)$, $g=e(f)$. The maps $e\restr f``\nu_n$ equal $g\circ(e\restr\nu_n)\circ f^{-1}$ and belong to $W$ because $e\restr\nu_n\in V_\lambda$. In $W$ form the tree of coherent finite lists of partial elementary maps $f``\nu_m\to g``\nu_{m+1}$ sending $\lambda\mapsto\lambda$, $\nu_0\mapsto\nu_1$. It has an external branch, so by absoluteness of well-foundedness it has a branch in $W$, whose union is the required embedding $X\to Z\prec V^W_\alpha$.
\end{proof}

\begin{lemma}[Finite-change cloud]\label{lem:cloud}\src{R9}
Let $c$ be Prikry generic over a transitive $N$ for a normal measure $D\in N$ on $\lambda$, $W=N[G_c]$, and $\tau\in N$ a name forced into $V_\alpha$. Let $Y_{\tau,c}=\{\tau^{G_h}:h\text{ increasing},\ |\ran h\triangle\ran c|<\omega\}$. Then $\tau^{G_c}\in Y_{\tau,c}\subseteq V^W_\alpha$, $Y_{\tau,c}\in W$, $|Y_{\tau,c}|^W\le\lambda$, and any elementary $e:W\to W$ fixing $\lambda,\PP^N_D,\tau,\alpha$ with $e(c)=c^-$ (the tail) satisfies $e(Y_{\tau,c})=Y_{\tau,c}$.
\end{lemma}

The value $\tau^{G_c}$ itself may move; cover exactingness tolerates exactly this motion inside an invariant small family.

\begin{lemma}[Localization]\label{lem:local}\src{R9, after ABGL Thm.~5.4}
Let $N_0$ be transitive of size $\kappa$, $V_\kappa\cup\{\kappa,U_0\cap N_0\}\subseteq N_0$, modelling enough of $\ZFC$. With $J=j_{0,\omega}$: $N=J(N_0)$ is well-founded, there is an elementary $i:N\to N$ with $i\restr\Lambda=j\restr\Lambda$ and $i(J(a))=J(a)$ for \emph{every} $a\in N_0$, the critical sequence $c$ is Prikry generic over $N$ for $J(U_0\cap N_0)$, and $i$ lifts to $e:N[G_c]\to N[G_c]$ with $e(c)=c^-$.
\end{lemma}

\begin{proof}[Proof of Theorem~\ref{thm:i3}]
Otherwise, by cone homogeneity of Prikry forcing, some $\alpha$ and name $\tau$ have $\mathbf1\Vdash$``no $Y\ni\tau$ of size $\le\kappa$ admits a relative witness at $\alpha$''. Collapse a hull $H\prec V_\zeta$ of size $\kappa$ containing $V_\kappa\cup\{U_0,\alpha,\tau\}$ to $N_0$ and apply Lemma~\ref{lem:local}; the images $a=J(\bar\alpha)$, $t=J(\bar\tau)$ are fixed by $e$. In $W=N[G_c]$ the cloud $Y_{t,c}$ is a set of size at most $\Lambda$ containing $t^{G_c}$ with $e(Y_{t,c})=Y_{t,c}$ (Lemma~\ref{lem:cloud}), so Lemma~\ref{lem:internal} yields a relative witness in $W$, contradicting the transferred forcing statement. For the transitive model, note $V_\Lambda[G]=V[G]_\Lambda$ and pass to a countable elementary submodel of $V_\Lambda$.
\end{proof}

\begin{corollary}[Strength interval]\label{cor:strength}\src{R9}
In consistency strength, $\Itwo>\Iwf\ge\exists\lambda\,\CEx_\lambda(\lambda)\ge\exists\lambda\,\Ex(\lambda)>\Ithree$; the strict comparisons are those of Input~\ref{in:iter}. No equiconsistency is claimed, and the ultraexacting clause cannot be read into the construction.
\end{corollary}

This lowers the $\Itwo$ upper bound of Input~\ref{in:bg}(iv). \R9 notes that the Blue--Goldberg notes anticipate cover variants of related constructions, so no priority is claimed. The construction does not preserve a prescribed strongly compact below $\kappa$ and so does not touch Problem~5.2.

\begin{corollary}[Calibration of the barrier]\label{cor:calib}\src{R1; with R6, R9}
$\Con(\ZFC+\Itwo)$, and by Theorem~\ref{thm:i3} already $\Con(\ZFC+\Iwf)$, implies
\[
\Con\bigl(\ZFC+\exists\lambda\,[\CEx_\lambda(\lambda)\wedge\cf^V(\lambda)=\omega\wedge\cf^{\HCD(\lambda)}(\lambda)=\lambda]\bigr).
\]
\end{corollary}

This adds no strength: the extra clauses are theorems about every such model (Theorem~\ref{thm:barrier}). Its role is to confirm that the barrier is a non-contradictory structural feature. The combination of \R1, \R6 and \R9 in the displayed form is made here.

\subsection{One \texorpdfstring{$\Izero$}{I0}-equiconsistent boundary theory}

Five first-round reports, all nine second-round reports, and eight third-round reports each derive an ``equiconsistency package'' from Input~\ref{in:abl}(iv),(v). Since every clause is a consequence of $\UEx(\lambda)\wedge V_\lambda\subseteq\HOD$, they merge into one statement.

\begin{theorem}[Merged boundary package]\label{thm:i0}\src{R2, R3, R5, R7, R8; R10--R18}
Let $T_{\partial}$ be $\ZFC$ plus the existence of a cardinal $\lambda$ with $\UEx(\lambda)$ and $V_\lambda\subseteq\HOD$. Then $\Con(T_\partial)\Leftrightarrow\Con(\ZFC+\Izero)$, and $T_\partial$ proves all of the following for that $\lambda$ and every cofinal $c\subseteq\lambda$ of order type $\omega$:
\begin{enumerate}
\item $\lambda$ is the least exacting cardinal.\src{known: ABGL Cor.~3.10}
\item $V^{\HOD}_\lambda=V_\lambda$; $\Pow(\alpha)^{\HOD}=\Pow(\alpha)$ and $\cf^{\HOD}(\alpha)=\cf(\alpha)$ for all $\alpha<\lambda$.\src{R3, R8}
\item $\cf(\lambda)=\omega$ while $\HOD\models$ ``$\lambda$ is strongly inaccessible''; $c\in V_{\lambda+1}\setminus\HOD$, so $\lambda$ is the first ordinal of cofinality or powerset disagreement and $\lambda+1$ the first rank of disagreement.\src{R2, R3, R8}
\item $c\cap a$ is finite and in $\HOD$ whenever $a\in\HOD$, $|a|^{\HOD}<\lambda$; every $b\in\HOD$ with $c\subseteq b$ has $|b|^V\ge\lambda$; so $(\HOD,V)$ fails $\theta$-approximation and $\theta$-covering for all uncountable $\theta\le\lambda$.\src{R2, R3, R8}
\item $Q^{\mathrm{bd}}_\lambda\in\HOD$ has an ordinal-definable well-order, whereas the conclusions of Theorems~\ref{thm:fsel}, \ref{thm:label}, \ref{thm:transversal} and Corollary~\ref{cor:orders} hold for $Q_\lambda$.\src{R5, R7}
\item Every conclusion of Section~\ref{sec:quot} holds, with $\OD$ in place of $\OD_{V_\lambda}$ since the two coincide: at least $\lambda$ rigid classes $q$, for which moreover $H_q=\HOD_{\{q\}}\models\ZFC$ agrees with $V$ below $\lambda$, regards $\lambda$ as inaccessible, and contains no member of $q$; full-width fibres for every $\OD$ complete section; the exact finite-label classification, including an $\OD$ choice-or-orientation rule on triples although neither pure rule is $\OD$; no $\OD$ probability kernel. The bounded quotient has an $\OD$ transversal.\src{R10--R18}
\item There are $\lambda$ classes $q$ with disjoint representatives such that $H_q=\HOD_{\{q\}}$ contains $V_\lambda$ and regards $\lambda$ as measurable, by the tail measure; for the critical-sequence class the measure is normal and concentrates on measurables, every representative is Prikry generic over $H_q$, and $V_\lambda^{H_q[a]}=V_\lambda$. $(\lambda^+)^V$ is measurable in every $\HOD_{\vec q_F}$, in particular in $\HOD$, with one club of simultaneous inaccessibles. No $\OD$ ultrafilter-valued kernel exists, although an $\OD$ charge-valued kernel does.\src{R19, R21--R23, R26, R27}
\end{enumerate}
\end{theorem}

\begin{proof}
Equiconsistency is Input~\ref{in:abl}(iv),(v). For (1): a smaller exacting $\mu$ would have a cofinal $\omega$-sequence of rank $\mu<\lambda$, hence in $\HOD$, contradicting Input~\ref{in:abl}(i). For (2): the relevant subsets and cofinal maps have rank below $\lambda$. (3),(4): $\HOD$-regularity is Input~\ref{in:abl}(i); strong limit as in Theorem~\ref{thm:ground}; the rest as in Corollary~\ref{cor:approx}. (5): Remark~\ref{rem:controls} and Section~\ref{sec:uex}. (6): Section~\ref{sec:quot}; under $V_\lambda\subseteq\HOD$ every parameter in $V_\lambda$, in particular a well-order of the reals, is ordinal definable. (7): Section~\ref{sec:measures}; under $V_\lambda\subseteq\HOD$ the bijection $b$ and the mean $M$ are ordinal definable and can be dropped from the parameters.
\end{proof}

The equiconsistency compares consistency statements; it does not say that $\Izero$ holds at the same ordinal, and it is not a new proof of the ABGL comparison. Its use is as a \emph{control}: total agreement with $\HOD$ below $\lambda$ does not supply the one countable sequence that would be fatal, so no negative argument may import covering at $\lambda$ from rank agreement below $\lambda$.

\section{Class embeddings: the internal unbounded-range obstruction}\label{sec:range}

Let $j:V\to M$ be a nontrivial elementary embedding into a transitive inner model (in a class theory where the needed restrictions exist), $\kappa=\crit(j)$, and let $\theta>\kappa$ be regular uncountable with $j(\theta)=\theta$. For $\lambda=\sup_nj^n(\kappa)$, $\theta=\lambda^+$ is such a cardinal \emph{without} any cardinal-correctness assumption: $j(\theta)=(\lambda^+)^M\le\lambda^+\le j(\theta)$ \R3. Let $C_j=\{\beta<\theta:j``\beta\subseteq\beta\}$, a club whose points of cofinality $\omega$ are fixed by $j$; fix a partition $\vec S$ of $E^\theta_\omega$ into $\theta$ stationary sets and let $\vec T=j(\vec S)$.

\begin{lemma}[Stationarity spectrum; known]\label{lem:spectrum}\src{HKP; proof in R3}
$\{\xi<\theta:T_\xi\text{ is stationary in }V\}=j``\theta$, and the single club $C_j$ is disjoint from every $T_\xi$ with $\xi\notin j``\theta$.
\end{lemma}

\R2, \R6 and \R7 each examined this reconstruction as a candidate result and rejected it as already contained in Hamkins--Kirmayer--Perlmutter \cite[Thm.~11--12]{hkp}; \R3 attributes it the same way and adds:

\begin{theorem}[No internal unbounded part of the range]\label{thm:range}\src{R3}
No $A\in M$ with $A\subseteq j``\theta$ is unbounded in $\theta$.
\end{theorem}

\begin{proof}
Let $f\in M$ enumerate such an $A$ increasingly. Since $A=j``B$ with $B$ unbounded, $f(\xi)=j(b_\xi)\ge j(\xi)$. So the $M$-club $D_f=\{\beta:f``\beta\subseteq\beta\}$ is contained in $C_j$. But $T_\kappa$ is stationary in $M$ and $\kappa\notin j``\theta$, so $D_f\cap T_\kappa\ne\varnothing=C_j\cap T_\kappa$.
\end{proof}

\begin{corollary}\label{cor:losses}\src{R3}
Every unbounded $B\subseteq\theta$ in $M$ contains $\theta$ many labels $\xi$ with $T_\xi$ stationary in $M$ but not in $V$. The ambient nonstationary ideal on $\theta$ is not amenable to $M$. (The set of surviving labels is nevertheless stationary in $V$: abundance in $V$ and availability in $M$ are different.)
\end{corollary}

\begin{remark}[A sanity control added in the synthesis]
The proof shows that $M$ contains no club subset of $C_j$, so no function in $M$ dominates $j\restr\theta$. One may check this on the ultrapower by a normal measure on $\kappa$ with $\theta>\kappa$ inaccessible. The $M$-definable set of $M$-strong-limit cardinals of cofinality $\omega$ is \emph{not} contained in $C_j$: if $\alpha>\kappa$ is a $V$-strong limit of cofinality $\kappa$, then $\sup j``\alpha=\alpha<j(\alpha)$ and $j(\alpha)$ is an $M$-strong limit, so the least $M$-strong limit $\beta'$ above $\alpha$ satisfies $\alpha<\beta'\le j(\alpha)$ and is not closed under $j$. The familiar fixed points (the $V$-strong limits of cofinality $\omega$) do not form a set in $M$.
\end{remark}

The result is a conditional inconsistency for a strengthened embedding theory. It uses no cardinal correctness, so the global cardinal-preserving problem of the research plan still lacks the implication from $\mathrm{Card}^M=\mathrm{Card}^V$ to any internal selection of externally stationary cells.

\section{Sections of the cofinal quotient}\label{sec:quot}

This section merges the second round of reports. Throughout, $\lambda$ has cofinality $\omega$, $D_\lambda$ is the set of cofinal subsets of $\lambda$ of order type $\omega$, $a=^*b$ means $|a\triangle b|<\omega$, $[a]$ is the class of $a$ and $Q_\lambda=D_\lambda/{=^*}$, as in Section~\ref{sec:uex}; $\Clam\supseteq D_\lambda$ is the set of short cofinal subsets of $\lambda$ (Section~\ref{sec:smallfam}). A \emph{complete section} is a set $T\subseteq D_\lambda$ meeting every class; it is \emph{thin} if $|T\cap q|<\lambda$ for every $q$, and a \emph{transversal} if $|T\cap q|=1$. No bound below $\lambda$ uniform in $q$ is ever assumed. $\OD_{V_\lambda\cup\{q\}}$ allows ordinals, finitely many parameters in $V_\lambda$, and the class $q$ as \emph{one} set parameter --- not its members. All cardinalities are computed in $V$.

Every class has exactly $\lambda$ members ($d\mapsto a\triangle d$ is a bijection from $[\lambda]^{<\omega}$), so $\lambda$ is the largest possible fibre; and $|Q_\lambda|=2^\lambda$ for a strong limit $\lambda$ of cofinality $\omega$ \R{15}.

\subsection{Fixed classes}

\begin{lemma}[Orbits and offsets]\label{lem:orbits}\src{R12, R15; R13, R14, R17; R10, R16, R18}
Let $j:X\to V_\zeta$ be an ultraexacting witness with $\zeta>\lambda+\omega$, $e=j\restr V_\lambda\in X$, $\kappa=\crit(j)$, $\kappa_n=e^n(\kappa)$.
\begin{enumerate}
\item The set of \emph{roots} $R_e=[\kappa,\lambda)\setminus e``[\kappa,\lambda)$ has size $\lambda$, and the forward orbits $a_r=\{e^n(r):n<\omega\}$, $r\in R_e$, partition $[\kappa,\lambda)$ into members of $D_\lambda$.
\item Each $a_r$ and $q_r=[a_r]$ belong to $X$, and $j(a_r)=a_r\setminus\{r\}$, so $j(q_r)=q_r$. The classes $q_r$ are pairwise distinct.
\item In particular this holds for the critical sequence ($r=\kappa$) and for the offsets $a_t=\{\kappa_n+t:n<\omega\}$, $t<\kappa$.
\end{enumerate}
\end{lemma}

\begin{proof}
$e$ maps $[\kappa_{n-1},\kappa_n)$ into $[\kappa_n,\kappa_{n+1})$, which has size $\kappa_{n+1}$, so $\kappa_{n+1}$ roots lie there. Every ordinal in $[\kappa,\lambda)$ has only finitely many $e$-predecessors (Lemma~\ref{lem:crit}), hence lies on a unique orbit. Orbits are definable from $e\in X$ and $r\in V_\lambda\subseteq X$; enumerations with domain $\omega$ are mapped pointwise.
\end{proof}

This is the only place where ultraexactingness, rather than exactingness, is used in this section: an exacting witness need not contain any fixed class. All nine reports record the question whether exactingness suffices as open.

\subsection{Rigid classes}

\begin{definition}[Rigidity]\label{def:rigid}\src{R18}
A class $q\in Q_\lambda$ is \emph{rigid}, $\Rig(\lambda,q)$, if every nonempty $A\in\OD_{V_\lambda\cup\{q\}}$ with $A\subseteq\Clam$ has $|A|\ge\lambda$.
\end{definition}

\begin{theorem}[Rigid classes exist]\label{thm:rigid}\src{R18; with R12, R15, R17}
Let $\lambda$ be ultraexacting and let $j:X\to V_\zeta$ be an ultraexacting witness at a sufficiently correct height. Then every $q\in X\cap Q_\lambda$ with $j(q)=q$ is rigid. Consequently at least $\lambda$ classes, with pairwise disjoint representatives, are rigid; and the bound in Definition~\ref{def:rigid} is attained, by $A=q$.
\end{theorem}

\begin{proof}
Exactly as for Theorem~\ref{thm:smallC}, with $q$ as an additional fixed parameter: minimize the rank of the $V_\lambda$-parameter (so that it falls below $\crit(j)$ and is fixed), take the least counterexample in the canonical well-order, and contradict Lemma~\ref{lem:fixedsmall}(b). Lemma~\ref{lem:orbits} supplies $\lambda$ fixed classes.
\end{proof}

\R{18} proves rigidity of the critical-sequence class; \R{17} obtains $\lambda$ classes that are saturated for all $\OD_{V_\lambda}$ sets simultaneously, using instead a bijection $b:\lambda\to V_\lambda$ in $X$ with $j(b)=b$ (so that $\OD_{V_\lambda}\subseteq\OD_b$) and a uniform bounded-rank correctness lemma. Applying the device of \R{18} to each orbit class of \R{12}, \R{15} gives both conclusions at once; this combination is made here.

\begin{corollary}[What rigidity gives]\label{cor:rigid}\src{R11, R15, R17, R18}
Let $q$ be rigid.
\begin{enumerate}
\item \emph{Zero-or-full traces.} $|T\cap q|\in\{0,\lambda\}$ for every $T\subseteq D_\lambda$ in $\OD_{V_\lambda\cup\{q\}}$.
\item \emph{Small families.} The same holds for every member $T$ of any family $\mathcal T\subseteq\Pow(D_\lambda)$ in $\OD_{V_\lambda\cup\{q\}}$ with $|\mathcal T|<\lambda$, and for all values $F(x,\alpha)$ of a definable $F:V_\lambda\times\theta\to\Pow(D_\lambda)$. The members need not be definable.
\item \emph{Cofinal summaries.} If $R\in\OD_{V_\lambda\cup\{q\}}$ assigns to each class a subset of $\Clam$, then $R(q)=\varnothing$ or $|R(q)|\ge\lambda$.
\item \emph{The model $N_q=\HOD_{V_\lambda\cup\{q\}}$.} $\cf^{N_q}(\lambda)=\lambda$; $q\notin N_q$ and $q\cap N_q=\varnothing$; every $\OD_{V_\lambda\cup\{q\}}$ set containing a member of $q$ has size at least $\lambda$; and the projection-width theorem~\ref{thm:width} holds relative to $q$.
\end{enumerate}
\end{corollary}

\begin{proof}
(1) $T\cap q$ is a definable subfamily of $D_\lambda\subseteq\Clam$. (2) For $T\in\mathcal T$ with $0<|T\cap q|<\lambda$ the union $u_T=\bigcup(T\cap q)$ is cofinal of size below $\lambda$; the set of all such $u_T$ is a definable subfamily of $\Clam$ of size at most $|\mathcal T|<\lambda$, hence empty. Each value $F(x,\alpha)$ is itself definable from $x$ and $\alpha$, so (1) applies to it. (3), (4) are instances of the definition.
\end{proof}

\begin{corollary}[Sharp width of definable sections]\label{cor:sharpwidth}\src{R10--R18}
Let $\lambda$ be ultraexacting.
\begin{enumerate}
\item Every complete section $T\in\OD_{V_\lambda}$ has at least $\lambda$ classes with $|T\cap q|=\lambda$, and one set of $\lambda$ classes serves all such $T$ \R{17}. For arbitrary $T\subseteq D_\lambda$ in $\OD_{V_\lambda}$, at least $\lambda$ fibres are empty or full \R{13}.
\item There is no thin complete section in $\OD_{V_\lambda}$: the \emph{definable width} $\min\{\sup_q|T\cap q|:T\in\OD_{V_\lambda}\text{ complete}\}$ is exactly $\lambda$. This contains Theorem~\ref{thm:transversal}.
\item No nonempty $\OD_{V_\lambda}$ family of fewer than $\lambda$ thin complete sections exists (in particular no countable definable family of transversals), and no such family of thin \emph{partial} sections has complete union \R{15}. Every $\OD_{V_\lambda}$ set having a thin complete section as a member has size at least $\lambda$ \R{12}, \R{15}.
\item Even the local demand fails: it is not the case that every class $q$ has some $T\in\OD_{V_\lambda}$ with $0<|T\cap q|<\lambda$ \R{17}.
\end{enumerate}
\end{corollary}

\begin{remark}[Sharpness]\label{rem:sharpquot}
(a) $T=D_\lambda$ is ordinal definable with all fibres of size $\lambda$; thin transversals exist by Choice. (b) The parameter rank cannot be raised: no class is saturated for all $\OD_{V_{\lambda+1}}$ sets (use $T=\{a\}$, $a\in q$) \R{17}. (c) The family bound cannot be raised in the local lemma: a fixed class is itself a fixed family of size $\lambda$ \R{10}. (d) Not every class is rigid: for each increasing cofinal $c$ there is a redundantly coded class $q_c$ from which $c$ is ordinal definable, so $\cf^{N_{q_c}}(\lambda)=\omega$ \R{11}. (e) For every class $q$ there are $\lambda$ distinct classes $q_t$ mutually ordinal definable with $q$, hence with $N_{q_t}=N_q$; so $\lambda$ distinct rigid classes can share one model \R{11}.
\end{remark}

\subsection{The structure of a full fibre}

For $a=^*b$ let $\ch_a(b)=|b\setminus a|-|a\setminus b|\in\mathbb Z$; then $\ch_a(c)=\ch_a(b)+\ch_b(c)$ and $\ch_{a\setminus\{\min a\}}=\ch_a+1$. For $B\subseteq[a]$ put $\Spec_a(B)=\{\ch_a(b):b\in B\}$ and $P_n(B)=\{b(n):b\in B\}$.

\begin{theorem}[Coordinate width and full phase]\label{thm:phase}\src{R12}
Let $j$ be an ultraexacting witness, $a\in D_\lambda\cap X$ with $j(a)=a\setminus\{\min a\}$, and $\varnothing\ne B\subseteq[a]$ with $B\in X$, $j(B)=B$. Then there is $n_0$ such that for every $r\in\mathbb Z$ and $n\ge n_0$ the phase $B_r=\{b\in B:\ch_a(b)=r\}$ has $P_n(B_r)$ cofinal in $\lambda$ of size $\lambda$. In particular $\Spec_a(B)=\mathbb Z$ and every phase has size $\lambda$. Hence every complete section $T\in\OD_{V_\lambda}$ has these properties on at least $\lambda$ classes.
\end{theorem}

\begin{proof}
$\Spec_a(B)=\Spec_{j(a)}(j(B))=\Spec_a(B)+1$, so the nonempty spectrum is $\mathbb Z$. The cardinals $\theta_r=|B_r|$ satisfy $j(\theta_r)=\theta_{r-1}$; a two-sided orbit of ordinals $\le\lambda$ is constant, with value $\lambda$ or below $\kappa$, and the latter would make $B$ a fixed family of size below $\lambda$. A fixed small $P_n(B)$ would lie below $\kappa$ by Lemma~\ref{lem:fixedsmall}(a).
\end{proof}

\begin{corollary}\label{cor:phasecons}\src{R12}
At an ultraexacting $\lambda$ no complete section in $\OD_{V_\lambda}$ (i) misses an integer phase in every class, or a residue class of phases modulo $m$; (ii) has in every class a pointwise bound $g_q:\omega\to\lambda$ on its members; or (iii) has in every class a finite bound on the number of points inserted, or deleted, relative to some reference representative. The two conditions are independent: there is a complete section with all fibres of size $\lambda$ and a single phase, and one with countable fibres and full spectrum.
\end{corollary}

\begin{corollary}[Colourings, coarser relations, homomorphisms]\label{cor:colour}\src{R14--R18}
Let $\lambda$ be ultraexacting.
(a) For every rigid $q$, every $f:D_\lambda\to\mathrm{Ord}$ in $\OD_{V_\lambda}$ and every $\xi$, $|q\cap f^{-1}\{\xi\}|\in\{0,\lambda\}$; so there is no $\OD_{V_\lambda}$ assignment of a well-order to each class that is constant on classes.
(b) If $E\in\OD_{V_\lambda}$ is an equivalence relation on $D_\lambda$ with $a\mathrel E(a\setminus\{\min a\})$ for all $a$, every $\OD_{V_\lambda}$ set meeting every $E$-class meets some $E$-class in at least $\lambda$ points.
(c) If $E\in\OD_{V_\lambda}$ has all classes of size below $\lambda$, there is no nonempty $\OD_{V_\lambda}$ family of fewer than $\lambda$ maps $f:D_\lambda\to D_\lambda$ with $a=^*b\Rightarrow f(a)\mathrel E f(b)$; for instance when $E$ is the orbit relation of a definable action of a group of size below $\lambda$.
\end{corollary}

\subsection{Rules with cofinal output and randomized selection}

\begin{theorem}\label{thm:conversion}\src{R10, R11, R16, R18}
If $\lambda$ is ultraexacting, there is no nonempty $\OD_{V_\lambda}$ family of fewer than $\lambda$ rules $R$ on $Q_\lambda$ with $\varnothing\ne R(q)\subseteq\Clam$ and $|R(q)|<\lambda$ for every $q$. The outputs need not represent, or be related to, the input class. Equivalently there is no such tail-invariant rule on $D_\lambda$; in particular no $\OD_{V_\lambda}$ map $Q_\lambda\to\Clam$.
\end{theorem}

\begin{proof}
Let $q$ be rigid and let $u(B)$ denote the least-order-type compression of Lemma~\ref{lem:fixedsmall}. For a family $\mathcal F$ as described, $\{u(R(q)):R\in\mathcal F\}$ is a nonempty subfamily of $\Clam$ of size below $\lambda$, definable from $\mathcal F$ and $q$, contradicting rigidity.
\end{proof}

\begin{corollary}[No definable randomized selection]\label{cor:kernel}\src{R10}
If $\lambda$ is ultraexacting there is no nonempty $\OD_{V_\lambda}$ family of fewer than $\lambda$ \emph{probability kernels}, i.e.\ assignments to each class $q$ of a countably additive probability law on $q$ (on the $\sigma$-algebra of Theorem~\ref{thm:laws}). Dirac kernels exist by Choice.
\end{corollary}

\begin{proof}
Theorem~\ref{thm:laws} converts a kernel into the rule $q\mapsto$ median set of $K(q)$.
\end{proof}

\subsection{The exact finite-label classification}\label{sec:class}

In the setting of Theorem~\ref{thm:label}, let $\Gamma\le S_m$ act on a finite nonempty set $F$, call a rule $L:D_m(\lambda)\to F$ \emph{admissible} if it is invariant under finite modification of the coordinates and $\Gamma$-equivariant, and let $\CF(\Gamma,F)$ say that every single $g\in\Gamma$ fixes some element of $F$.

\begin{lemma}[Cyclic stabilizers]\label{lem:cyclic}\src{R13}
To $\vec a\in D_m(\lambda)$ assign its \emph{incidence word} $w_{\vec a}(n)=\{i<m:\delta_n\in a_i\}$, where $\langle\delta_n\rangle$ enumerates $\bigcup_ia_i$; finite modifications change $w$ only up to \emph{shift-tail equivalence} ($w(r+n)=v(s+n)$ for all $n$), and $w_{g\cdot\vec a}=g\cdot w_{\vec a}$. For every word $w$ in which each label occurs infinitely often and any two labels are separated infinitely often, the stabilizer in $\Gamma$ of the shift-tail class of $w$ is cyclic.
\end{lemma}

\begin{proof}
$G_w=\{(d,g)\in\mathbb Z\times\Gamma:w(n+d)=g\cdot w(n)\text{ for all large }n\}$ is a subgroup whose projection to $\mathbb Z$ is injective (if $(0,g)\in G_w$ and $g$ moves $i$, the labels $i$ and $g^{-1}i$ are not separated infinitely often). So $G_w$ is cyclic, and the stabilizer is its image in $\Gamma$.
\end{proof}

\begin{theorem}[Classification]\label{thm:classification}\src{R13; necessity R7}
Let $\lambda$ be ultraexacting. An admissible rule $L:D_m(\lambda)\to F$ in $\OD_{V_\lambda}$ exists if and only if $\CF(\Gamma,F)$. For sufficiency ultraexactingness is not needed: for every uncountable $\lambda$ of cofinality $\omega$ the rule is definable from a well-order of $\Pow(\omega)$, a parameter in $V_{\omega+10}$; it is $\OD$ if the reals have an $\OD$ well-order, in particular if $V_\lambda\subseteq\HOD$.
\end{theorem}

\begin{proof}
Necessity is Theorem~\ref{thm:label}. For sufficiency, the rule factors through shift-tail classes of words, which are coded by reals. In each $\Gamma$-orbit of classes take the class of the least coded word $v$; its stabilizer is cyclic (Lemma~\ref{lem:cyclic}), a generator fixes some label by $\CF$, and so the whole stabilizer fixes the least such label $f$; put $L(g\cdot v)=g\cdot f$.
\end{proof}

\begin{corollary}\label{cor:five}\src{R13}
(a) At an ultraexacting $\lambda$ there is an $\OD_{V_\lambda}$ rule assigning to every unordered triple of classes \emph{either} a distinguished member \emph{or} a cyclic orientation, although no $\OD_{V_\lambda}$ rule always does the first and none always does the second; every such mixed rule uses each output type on at least $\lambda$ triples. (b) The five-element $S_3$-set of Section~\ref{sec:uex} is, up to isomorphism, the smallest $S_3$-set satisfying $\CF$ without a global fixed point. (c) A finite group admits a finite set with $\CF$ and no global fixed point iff it is not cyclic. (d) If strictly increasing maps $f$ satisfy $f``a_i=^*a_{\pi_f(i)}$ on one common tuple, the $\pi_f$ generate a cyclic group: several embeddings cannot realize a noncyclic relabelling group on one tuple, which is why the elementwise test cannot be upgraded to a whole-group test.
\end{corollary}

This settles the question left open in the first edition after Theorem~\ref{thm:label}, and shows that the obstruction of Section~\ref{sec:uex} is exactly the cyclic one.

\subsection{Preserved-predicate versions}

\begin{theorem}\label{thm:preserved}\src{R10--R13, R15--R18}
(a) Let $e:(V_{\lambda+1},\in,T)\to(V_{\lambda+1},\in,T)$ be elementary with $\crit(e)<\lambda$ and $T\subseteq D_\lambda$. Then $|T\cap[a]|\in\{0,\lambda\}$ for the critical-sequence set $a$ and for each offset $a_t$, $t<\crit(e)$; so $T$ is not a thin complete section.
(b) Let $N\models\ZF$ be a transitive inner model with $V_{\lambda+1}\subseteq N$ and $j:N\to N$ elementary with $\crit(j)<\lambda=\sup_nj^n(\crit j)$ and set-sized restrictions. If $\mathcal T\in N$, $j(\mathcal T)=\mathcal T$, $\mathcal T\subseteq\Pow(D_\lambda)$ and $|\mathcal T|^V<\lambda$, then $|T\cap q|^V\in\{0,\lambda\}$ for all $T\in\mathcal T$ and all $\lambda$ orbit classes $q$; likewise Theorem~\ref{thm:phase} holds for a fixed complete section, and no $f:Q_\lambda\to D_\lambda$ in $N$ is fixed. No Choice in $N$ is used, and all sizes are ambient.
(c) Hence for a thin (or phase-proper) complete section $T$, or for the code $A_f$ of a map $f:Q_\lambda\to D_\lambda$, there is no elementary $j:L(V_{\lambda+1},T)\to L(V_{\lambda+1},T)$ with $\crit(j)<\lambda$ and $j(T)=T$.
\end{theorem}

\begin{proof}
(a) The fibre is a subset, not an element, of $V_{\lambda+1}$; one applies $e$ to the elements $\bigcup(T\cap[a])$ \R{13} or $P_n(T\cap[a])$ \R{16} of $V_{\lambda+1}$, which are fixed and would be small. (b), (c) as in Theorems~\ref{thm:icarus} and~\ref{thm:rigid}, the orbits now being available because $V_{\lambda+1}\subseteq N$.
\end{proof}

These extend Theorem~\ref{thm:icarus} from finite families of selectors to thin sections. As there, they are forbidden \emph{information conditions}: no report shows that a canonical Icarus enrichment defines a thin section.

\subsection{Periodic classes of one embedding}

\begin{proposition}\label{prop:periodic}\src{R15}
For $e:V_\lambda\to V_\lambda$ let $\bar e([a])=[e``a]$. A class is fixed by $\bar e^{\,r}$ iff it has a representative that is a finite nonempty union of forward $e^r$-orbits in $[\kappa,\lambda)$. There are exactly $\lambda$ classes of each least period $r\ge1$, hence exactly $\lambda$ periodic classes, out of $2^\lambda$; every periodic class lies in $X$, and Corollary~\ref{cor:rigid}(1),(2) holds at each of them for families fixed by $j$.
\end{proposition}

So a single witness cannot give more than $\lambda$ test classes; whether a definable complete section must have more than $\lambda$ full fibres is open.

\section{Measures hidden in rigid classes}\label{sec:measures}

This section merges the third round on the quotient $Q_\lambda$. Notation is that of Section~\ref{sec:quot}. For $q\in Q_\lambda$ and $A\subseteq\lambda$ write $\operatorname{Tail}(q,A)$ if some (equivalently every) representative of $q$ is almost contained in $A$. For a parameter $d$ put $N(d,q)=\HOD_{\{d,q\}}$, an inner model of $\ZFC$; $H_q=N(\varnothing,q)$; and $N_q=\HOD_{V_\lambda\cup\{q\}}$, an inner model of $\ZF$ containing $V_\lambda$. The \emph{tail filter} is $U(d,q)=\{A\in\Pow(\lambda)\cap N(d,q):\operatorname{Tail}(q,A)\}$, and similarly $U_q^{N_q}$.

\subsection{Rigidity is regularity}

\begin{theorem}\label{thm:rigreg}\src{R25}
For every uncountable $\lambda$ of cofinality $\omega$ and $q\in Q_\lambda$: $\Rig(\lambda,q)$ iff $\lambda$ is regular in $N_q$, iff no short cofinal subset of $\lambda$ is in $\OD_{V_\lambda\cup\{q\}}$. For every class $q$ there are $\lambda$ classes with pairwise disjoint representatives that are mutually ordinal definable with $q$; so one rigid class yields $\lambda$ rigid classes with the same $N_q$, and Corollaries~\ref{cor:rigid} and~\ref{cor:sharpwidth} hold at all of them.
\end{theorem}

\begin{proof}
A small definable family of short cofinal sets compresses to a single definable short cofinal set of ordinals (Lemma~\ref{lem:fixedsmall}(b)), which lies in $N_q$; the converse is the case of a singleton. The copies are those of Remark~\ref{rem:sharpquot}(e).
\end{proof}

This answers the first half of Question (iv) of the second edition: rigidity has an intrinsic description with no reference to an embedding.

\subsection{Tail measures and Prikry cores}

\begin{theorem}[Tail decisions]\label{thm:taildec}\src{R22, R23, R26, R27; R21}
Let $j:X\to V_\zeta$ be an ultraexacting witness at a sufficiently correct height, $d\in X$ with $j(d)=d$ (for instance $d=\varnothing$, or the fixed bijection $b:\lambda\to V_\lambda$ of \R{17}), $r$ a root, $a_r$ its orbit and $q_r=[a_r]$.
\begin{enumerate}
\item Every $A\subseteq\lambda$ in $\OD_{\{d,q_r\}}$, and every $A\in\OD_{V_\lambda\cup\{q_r\}}$, contains almost all of $a_r$ or almost none. Every $f:\lambda\to\tau$, $\tau<\lambda$, in these classes is constant on a tail of $a_r$.
\item For the critical-sequence class $q=q_\kappa$, every such $f:\lambda\to\lambda$ that is regressive on a tail of $a_\kappa$ is constant on a tail of $a_\kappa$.
\end{enumerate}
\end{theorem}

\begin{proof}
Take the canonical least counterexample (minimal-rank parameter, Theorem~\ref{thm:smallC}); it is fixed. For fixed $A$: $e^n(r)\in A\iff e^{n+1}(r)\in j(A)=A$, so membership is constant along the orbit. For fixed $f$ with fixed $\tau<\lambda$, hence $\tau<\kappa$: $f(e^{n+1}(r))=j(f(e^n(r)))=f(e^n(r))$. (2) For regressive $f$, $f(\kappa_0)<\kappa_0=\crit(j)$ is fixed, and $f(\kappa_{n+1})=j(f(\kappa_n))$. This uses $r=\kappa$: for another root, $f(r)<r$ need not be fixed.
\end{proof}

\begin{theorem}[Internal normal measure and Prikry core]\label{thm:core}\src{R22, R23, R26, R27}
In the situation of Theorem~\ref{thm:taildec}:
\begin{enumerate}
\item For every root $r$, $U(d,q_r)\in N(d,q_r)$ is there a nonprincipal $\lambda$-complete ultrafilter on $\lambda$; so $\lambda$ is measurable in $N(d,q_r)$, for $\lambda$ classes with disjoint representatives. Likewise $U^{N_{q_r}}_{q_r}\in N_{q_r}$ is an ultrafilter closed under intersections of $N_{q_r}$-sequences of length below $\lambda$ (in $\ZF$).
\item For $q=q_\kappa$ these ultrafilters are \emph{normal}, and $U(d,q)$ concentrates on the cardinals that are measurable in $N(d,q)$ \R{27}. $U(d,q)$ is the unique ultrafilter of $N(d,q)$ all of whose members contain a tail of a representative of $q$.
\item Every representative of $q$ is Prikry generic over $N(d,q)$ for $U(d,q)$, and all give the same model $K(d,q)=N(d,q)[a]$, with $N(d,q)\subsetneq K(d,q)\subseteq V$, $q\in K(d,q)$, $q\cap N(d,q)=\varnothing$, the same cardinals, and $\cf^{K(d,q)}(\lambda)=\omega$. One class $q$ serves all $p\in V_\lambda$ simultaneously, with $U(p,q)=U(r,q)\cap N(p,q)$ whenever $\OD_{\{p,q\}}\subseteq\OD_{\{r,q\}}$ \R{23}.
\item With $d=b$, or with $d=\varnothing$ under $V_\lambda\subseteq\HOD$, $V_\lambda^{N(d,q)}=V_\lambda^{K(d,q)}=V_\lambda$: the ambient singularity of $\lambda$ already occurs in a Prikry extension of a model that contains all of $V_\lambda$ and regards $\lambda$ as measurable.
\end{enumerate}
\end{theorem}

\begin{proof}
(1) The filter is ordinal definable from $d,q_r$ and a subset of the model. Ultrafilter and completeness are Theorem~\ref{thm:taildec}(1): for a sequence $\langle A_i:i<\mu\rangle$ in the model, colour $\alpha$ by the least $i$ with $\alpha\notin A_i$. (2) Normality is Theorem~\ref{thm:taildec}(2); each $\kappa_n$ is measurable in $V$ with a measure definable enough to survive, and membership in the definable set of internally measurable cardinals is decided on a tail. (3) Mathias' criterion: a cofinal $\omega$-sequence is Prikry generic for a normal measure iff it is almost contained in every measure-one set; finite changes preserve this and generate the same extension.
\end{proof}

\begin{remark}[Limits]\label{rem:corelimits}
(a) Completeness is internal: $\lambda\setminus(\kappa_n+1)\in U(d,q)$ for all $n$, with empty intersection in $V$. (b) Normality is a property of the class, not of the model: $q^+=[\{\kappa_n+1\}]$ is mutually definable with $q$, so $N(d,q^+)=N(d,q)$, but its tail ultrafilter is not normal; and the class of $\{\kappa_n\}\cup\{\kappa_n+1\}$ gives a $\lambda$-complete filter that is not an ultrafilter \R{22}, \R{23}, \R{26}. (c) $K(d,q)$ is not claimed to equal $V$, and $V$ is not a $\lambda$-cc extension of $N(d,q)$. (d) \emph{Intrinsic criterion} \R{22}: for an arbitrary class $q$, the tail filter is a normal measure in $H_q$ iff some representative is Prikry generic over $H_q$ for some normal measure; then the measure is the tail filter.
\end{remark}

\begin{corollary}[Definable splitting]\label{cor:split}\src{R22, R23, R26}
If $\lambda$ is ultraexacting, it is not the case that every class $q$ has a set $A\in\OD_{V_\lambda\cup\{q\}}$ splitting it ($A$ and $\lambda\setminus A$ both meeting a representative infinitely often); in particular there is no $\OD_{V_\lambda}$ function $F$ with $F(q)$ splitting $q$ for all $q$. Nor can $\lambda$ be non-measurable in every $H_q$, nor can every $\HOD_{\{b,q\}}$ have internal Mitchell order $o(\lambda)\le1$ \R{27}.
\end{corollary}

\subsection{The successor of $\lambda$}

\begin{theorem}\label{thm:succmeas}\src{R19}
Let $\lambda$ be ultraexacting, $\theta=(\lambda^+)^V$, $S=(E^\theta_\omega)^V$. There are an inaccessible $\kappa<\lambda$, a bijection $b:\lambda\to V_\lambda$ and $\lambda$ orbit classes $q_i$ with disjoint representatives such that for every finite $F\subseteq\lambda$ and $N_F=\HOD_{b,\vec q_F}$:
\begin{enumerate}
\item $V_\lambda^{N_F}=V_\lambda$, $\lambda$ is inaccessible in $N_F$, and no member of any $q_i$, nor any $q_i$, belongs to $N_F$.
\item $I_F=\{A\in\Pow(S)\cap N_F:A\text{ nonstationary in }V\}\in N_F$ is a $\theta$-complete ideal whose quotient is atomic with fewer than $\kappa$ atoms; each atom gives a normal measure on $\theta$ in $N_F$. So $\theta$ is measurable in $N_F$ and $(\lambda^+)^{N_F}<\theta$.
\item One club $C\subseteq\theta$ of $V$ has all its points of cofinality $\omega$ inaccessible in every $N_F$; $(E^\theta_\omega)^{N_F}$ is stationary in $N_F$ and nonstationary in $V$; and $V$ is not an extension of $N_F$ by a forcing that is $\theta$-cc in $N_F$.
\end{enumerate}
\end{theorem}

\begin{proof}
A surjection $V_{\lambda+1}\to\theta$ definable from $\lambda,\theta$ and the canonical extension $e^+\in X$ yield $t\in X$, $t:\theta\to\theta$, agreeing with $j$ on $X\cap\theta$; its fixed points form an $\omega$-closed unbounded set meeting every stationary subset of $S$. As in Kunen's argument there is then no $\OD_p$ partition of $S$ into $\kappa$ sets stationary in $V$, for a fixed parameter $p$. So $\Pow(S)^{N_F}/I_F$ is $\kappa$-cc with $2^\kappa<\theta$, hence atomic with fewer than $\kappa$ atoms.
\end{proof}

This relativizes to the models $N_F$ the theorem of Aguilera--Bagaria--L\"ucke that $\lambda^+$ is $\omega$-strongly measurable in $\HOD$ at an ultraexacting $\lambda$ \cite[Cor.~3.15]{abl}; the additions are the parameters $b,\vec q_F$, the uniform club, and the explicit measures. With a strongly compact $\delta<\lambda$, Theorem~\ref{thm:successor} applies to $W=K(d,q)$ or $N(d,q)$: its $\lambda^+$ is collapsed or its $(E_\lambda)$ is destroyed \R{23}.

\subsection{Exact consistency strength}

\begin{theorem}\label{thm:strength}\src{R25, R27}
(a) The following are equiconsistent over $\ZFC$: a measurable cardinal; a strong limit singular $\lambda$ of cofinality $\omega$ that is regular in $\HOD_{V_\lambda}$; the existence of a rigid class at such a $\lambda$; the existence of $\lambda$ rigid classes with disjoint representatives and a common $N_q$; and the normal-trace principle $\mathsf{NT}$: for some $\lambda,b,q$, the tail filter is a normal measure on $\lambda$ in $\HOD_{\{b,q\}}$.
(b) $\mathsf{NT}^+$, which adds that the measure concentrates on the internally measurable cardinals, is equiconsistent with a cardinal of Mitchell order $o(\kappa)\ge2$.
\end{theorem}

\begin{proof}
Upper bounds: Prikry forcing $\PP_U$ at a measurable $\kappa$. Any two conditions have direct extensions with isomorphic cones, by $(s^\frown u,D)\mapsto(t^\frown u,D)$, and this isomorphism preserves the canonical name of the \emph{class} $[c]$ of the generic sequence (though not of $c$). So every set of ordinals definable in $W[G]$ from $[c]$ and ground parameters lies in $W$; thus $\kappa$ is regular in $N_{[c]}$ and, by Theorem~\ref{thm:rigreg}, $[c]$ is rigid; the tail filter is the restriction of $U$. For (b) start with $U$ concentrating on measurables, equivalently $o(\kappa)\ge2$. Lower bounds: a singular cardinal that is regular in an inner model gives an inner model with a measurable (Dodd--Jensen \cite{dj}); an internal normal measure concentrating on internal measurables gives $o\ge2$ in that inner model.
\end{proof}

\begin{assessment}{What this calibration means}
Every conclusion of Section~\ref{sec:quot} that follows from rigidity alone --- zero-or-full traces, the sharp definable width $\lambda$, the absence of thin sections, small families, colourings, cofinal-output rules, probability kernels --- already holds in the Prikry extension of a model with one measurable cardinal. These are theorems about definability at a singular cardinal that is regular in a $\HOD$-like model; ultraexactingness is one sufficient condition for them and is not detected by them. What is specific to the embedding is: the simultaneity over finite sets of orbit classes and the measurability of $\lambda^+$ in $N_F$ (Theorem~\ref{thm:succmeas}), and agreement $V_\lambda^{N}=V_\lambda$ at the $\Izero$ level. (The third edition listed here also the necessity half of Theorem~\ref{thm:classification} and Theorem~\ref{thm:fsel}, as lacking a Prikry-model counterpart; Section~\ref{sec:prikry} supplies it.)
\end{assessment}

\begin{proposition}[The Vop\v enka algebra of a class]\label{prop:vopenka}\src{R25}
For a rigid $q$ at a strong limit $\lambda$ there is a complete Boolean algebra $\mathbb B_q\in H_q$, atomless and weakly homogeneous under the action of the finite-change group, such that every $b\in q$ induces an $H_q$-generic $G_b$ with $H_q[G_b]=\HOD_{\{q,b\}}$, the same model $M_q$ for all $b\in q$. $\mathbb B_q$ singularizes $\lambda$ to cofinality $\omega$, is not $\lambda$-cc in $H_q$, and adds no bounded subsets of $\lambda$ if $V_\lambda\subseteq\HOD$. The $\lambda$ filters $G_b$ are pairwise not mutually generic.
\end{proposition}

\subsection{Finitely additive kernels}

A \emph{charge} on a set is a finitely additive probability measure on its power set; a (charge-valued) \emph{kernel} assigns to each $q\in Q_\lambda$ a charge on $q$. Corollary~\ref{cor:kernel} excluded definable countably additive kernels.

\begin{theorem}\label{thm:charges}\src{R21}
(a) \emph{A definable kernel exists.} Let $M$ be a shift-invariant mean on $\omega$ (a set in $V_{\omega+20}$). Then $K_M(q)(A)=M\bigl(\langle\mathbf 1_A(a\setminus\{\text{first }n\text{ points}\}):n<\omega\rangle\bigr)$ does not depend on $a\in q$ and defines a kernel in $\OD_{\{M\}}\subseteq\OD_{V_\lambda}$, for every uncountable $\lambda$ of cofinality $\omega$. It gives mass $1$ to every countable tail ray and mass $0$ to every point and to every event $\{b:b(i)<\alpha\}$.
(b) \emph{Null partitions.} If no short cofinal subset of $\lambda$ is in $\OD_z$, every nonempty $\OD_z$ family of fewer than $\lambda$ charges on a set $S\subseteq D_\lambda$ has a common countable partition of $S$ into null sets; so its members are purely finitely additive. This holds at exacting $\lambda$ for $z\in V_\lambda$, and at rigid $q$ for $z=q$.
(c) \emph{Phase balance.} At each of the $\lambda$ orbit classes $q$ of an ultraexacting witness, every charge $\mu$ on $q$ in $\OD_{V_\lambda\cup\{q\}}$ satisfies $\mu(\{b:\ch_a(b)\equiv r\bmod m\})=1/m$ for all $m\ge1$, $r$, and $a\in q$; every phase is null.
(d) \emph{No definable ultrafilters.} Hence no nonempty finite $\OD_{V_\lambda\cup\{q\}}$ family of finite-range charges on such $q$ exists, in particular no definable ultrafilter on $q$; and an ultraexacting $\lambda$ admits no nonempty finite $\OD_{V_\lambda}$ family of ultrafilter-valued kernels. The definable family $q\mapsto\{\text{charges on }q\}$ of compact convex sets has a definable point selection but no definable extreme-point selection.
\end{theorem}

\begin{proof}
(a) Two representatives have eventually synchronized tails, and $M$ is shift invariant. (b) For each coordinate $n$ the distribution functions $\alpha\mapsto\mu(\{a(n)<\alpha\})$ take, over the small definable family, fewer than $\lambda$ jump points, a definable set of ordinals, hence bounded by some $\beta$; the sets $\{a:a(n)<\beta\le a(n+1)\}$ and their analogues give the partition. (c) For a fixed charge $\mu$ with $j(a)=a\setminus\{\min a\}$ the phase distribution of $\mu$ is invariant under $r\mapsto r+1$, as in Theorem~\ref{thm:phase}. (d) A finite-range charge takes only finitely many values, contradicting (c); extreme points of the set of charges are exactly the ultrafilters.
\end{proof}

So randomized selection is excluded exactly at countable additivity: between the $\{0,1\}$-valued and the countably additive there is a definable finitely additive kernel, and every plateau value in $[0,1]$ occurs for suitable mixtures \R{21}. Whether a countably infinite definable family of ultrafilter-valued kernels can exist at an ultraexacting cardinal is open; in the Prikry model one exists, with a ground-model parameter (Theorem~\ref{thm:gap}).

\subsection{Vop\v enka reflection in the core (fourth round)}\label{sec:vopenka}

For $A\subseteq V_\lambda$ let $E_A$ be the set of $\delta<\lambda$ that are \emph{$A$-extendible below $\lambda$}: for every $\beta\in(\delta,\lambda)$ there is an elementary $i:(V_\beta,A\cap V_\beta)\to(V_\gamma,A\cap V_\gamma)$ with $\gamma<\lambda$, $\crit(i)=\delta$ and $i(\delta)>\beta$. A cardinal $\lambda$ is Vop\v enka iff it is inaccessible and every $E_A$, $A\subseteq V_\lambda$, is nonempty; the sets $E_A$ generate the Vop\v enka filter.

\begin{theorem}[Predicate reflection]\label{thm:vopenka}\src{R29}
In the situation of Theorem~\ref{thm:core} with $d=b$, let $N=\HOD_{\{b,q\}}$ and $U=U(b,q)$. For every $A\in\Pow(V_\lambda)^N$, $E^N_A=E^V_A\in U$; so a tail of the critical sequence consists of $A$-extendible cardinals. Hence $N\models$ ``$\lambda$ is a measurable Vop\v enka cardinal and $U$ extends the Vop\v enka filter'', with diagonal intersections over sequences of predicates in $N$; one embedding can preserve all predicates $A_\xi$, $\xi<\delta$, at once. Also $\HOD_{\{b\}}\models$ ``$\lambda$ is Vop\v enka''; and if $\lambda$ is merely exacting and $V_\lambda\subseteq\HOD$, then $\HOD\models$ ``$\lambda$ is Vop\v enka''.
\end{theorem}

\begin{proof}
The embeddings in question have rank below $\lambda$, so they belong to $N$ and $E^N_A=E^V_A$. Take the least $A$, in the canonical well-order of $N$, with $E_A\notin U$; it is fixed by $j$, so $e=j\restr V_\lambda$ is elementary on $(V_\lambda,A)$, and $e^n\restr V_\beta$ witnesses that $\kappa$ is $A$-extendible below $\lambda$ for $\beta<\kappa_n$. So $\kappa\in E_A$, and $E_A$ is a fixed set, hence contains the whole critical sequence (Theorem~\ref{thm:taildec}) and belongs to $U$ --- a contradiction. For ordinary exactingness use Input~\ref{in:abl} in the same way, without the measure.
\end{proof}

The quantifier over predicates is internal to $N$, and must be: a representative $c\in q$, used as a predicate, has $E^V_c=\varnothing$, because the $n$th element $\alpha\ge\delta$ of $c$ is definable without parameters in $(V_\beta,c\cap V_\beta)$ and in the target, so $i(\alpha)=\alpha$, while $i(\alpha)\ge i(\delta)>\beta>\alpha$ \R{29}. The corresponding principles (a normal measure containing all $E_A$, and its trace version after Prikry forcing) are equiconsistent with each other; their exact strength is not determined.

\section{The Prikry model: finite symmetry at measurable strength}\label{sec:prikry}

This section merges the fourth round (except \R{29}). Throughout, $W\models\ZFC$, $U\in W$ is a normal measure on $\kappa$, $c$ is Prikry generic for $U$ over $W$, $V=W[c]$, $\lambda=\kappa$, $q=[\ran c]\in Q_\kappa$. In $V$, $\kappa$ is a strong limit cardinal of cofinality $\omega$, $V_\kappa=V^W_\kappa$ and all cardinals are preserved. The class $\mathcal D(W,q)$ consists of the sets that are definable in $V$ from ordinals, $q$, and finitely many individual elements of $W$; no predicate for $W$ is allowed. It contains $\OD_{V_\kappa\cup\{q\}}$. \R{32} and \R{35} allow in addition any parameter $z(c)$ that is \emph{tail invariant} (unchanged when $c$ is changed finitely), such as a finite tuple or an indexed family of classes built from $c$.

\subsection{Stem surgery}

\begin{lemma}[Insertion]\label{lem:insertion}\src{R28, R30--R36}
(a) For stems $s,t$ and $B\in U$ above both, $(s^\frown u,C)\mapsto(t^\frown u,C)$ is an isomorphism of the cones below $(s,B)$ and $(t,B)$ fixing all check names; corresponding generics $a,d$ satisfy $W[a]=W[d]$ and $[a]=[d]$. (b) Below any condition $(s,A)$ and for any $\alpha\in A$ there are generics $a=s^\frown h$ and $d=s^\frown\langle\alpha\rangle^\frown h$, both below $(s,A)$, with $W[a]=W[d]$, $[a]=[d]$ and $d_{n+1}=a_n$ for $n\ge|s|$. (c) \emph{Decide, then insert.} Let $T(c)$ be a name for an element of a finite set $Y\in W$, defined from $c$ and parameters in $\mathcal D(W,q)$, and let $\rho$ be a permutation of $Y$ such that $T(d)=\rho(T(a))$ whenever $a,d$ are as in (b). Then $T(c)\in\operatorname{Fix}(\rho)$.
\end{lemma}

\begin{proof}
(c) Some condition decides $T(c)=y$. Take $a,d$ below it as in (b). Every parameter has the same interpretation for $a$ and for $d$, because the extension and the class are literally the same; only the enumeration has shifted. So $y=T(d)=\rho(T(a))=\rho(y)$.
\end{proof}

\R{28} deletes a point instead of inserting one, \R{32} calls the operation splicing and \R{33} tail recoding; these are the same lemma. It replaces the embedding in every argument of Sections~\ref{sec:uex} and~\ref{sec:quot} that used only ``$j(a)=a$ minus a point and $j(q)=q$''.

\begin{lemma}[Ground capture]\label{lem:capture}\src{R30--R33, R35; R25}
Every set of ordinals in $\mathcal D(W,q)$, also with tail-invariant parameters, belongs to $W$. Hence $\HOD_{\{p,q\}}\subseteq W$ for $p\in W$, and with a bijection $b:\kappa\to V_\kappa$ in $W$, $V_\kappa\subseteq N=\HOD_{\{b,q\}}\subseteq W$.
\end{lemma}

\subsection{Rigidity with ground parameters}

\begin{theorem}\label{thm:prikryrigid}\src{R30, R33, R36; $2^\kappa$ classes R31, R32, R35}
\begin{enumerate}
\item Every nonempty $\mathcal A\in\mathcal D(W,q)$ with $\mathcal A\subseteq\mathcal C_\kappa$ has $|\mathcal A|\ge\kappa$; so $q$ is rigid, $|T\cap q|\in\{0,\kappa\}$ for every $T\subseteq D_\kappa$ in $\mathcal D(W,q)$, and $\kappa$ is regular in $N_q$.
\item \emph{Local constancy.} Every map $h:q\to E$ in $\mathcal D(W,q)$ with $E\in W$ is locally constant in the topology on $q$ generated by the sets $\{b\in q:s\sqsubseteq b,\ b\setminus s\subseteq A\}$; every such $T\subseteq q$ is clopen. If $T\cap q\ne\varnothing$ then $|\{b\in T\cap q:\ch_a(b)=r\}|=\kappa$ for every $a\in q$ and $r\in\mathbb Z$ \R{33}.
\item \emph{$2^\kappa$ rigid classes.} For $A\in\Pow(\kappa)^W$ let $q_A=[\{e(c_n,A\cap c_n):n<\omega\}]$, where $e\in W$ is an injection of $\{(\alpha,x):\alpha<\kappa,\ x\subseteq\alpha\}$ into $\kappa$ with $e(\alpha,x)\ge\alpha$ (\R{31}; \R{32} and \R{35} use other codings). The $q_A$ are $2^\kappa$ distinct classes with pairwise almost disjoint representatives, $\OD_{\{e,q_A\}}=\OD_{\{e,A,q\}}$, every $\HOD_{\{p,q_A\}}\subseteq W$ regards $\kappa$ as measurable by its tail filter, and every $T\subseteq D_\kappa$ definable from ground sets and finitely many $q_A$ has $|T\cap q_A|\in\{0,\kappa\}$ for all $A$. So every such complete section has $2^\kappa=|Q_\kappa|$ full fibres. The whole indexed family is one tail-invariant parameter $z$, and $\HOD_{\{z\}}$ still carries a normal measure \R{35}.
\item \emph{Rank capture} \R{32}. With a parameter coding $\Pow(\kappa)^W$, the core $N$ has $U_N=U$, $(\kappa^+)^N=(\kappa^+)^V$, $V^N_{\kappa+1}=V^W_{\kappa+1}$ and $V^{N[c]}_{\kappa+1}=V_{\kappa+1}$.
\end{enumerate}
\end{theorem}

Item (3) is the maximal possible number; at an ultraexacting cardinal only $\lambda$ rigid classes are known (Theorem~\ref{thm:rigid}).

\subsection{Finite labels and finite families of selectors}

\begin{theorem}[Classification in the Prikry model]\label{thm:prikryclass}\src{R28, R30--R36}
For a finite $\Gamma\le S_m$ and a finite $\Gamma$-set $F\in W$, the following are equivalent: $\CF(\Gamma,F)$; there is an admissible rule $D_m(\kappa)\to F$ in $\OD_{V_\kappa}$; there is one in $\mathcal D(W,q)$, also with tail-invariant parameters; there is a nonempty \emph{finite family} of admissible rules in $\mathcal D(W,q)$ \R{33}, \R{36}.
\end{theorem}

\begin{proof}
Sufficiency is Theorem~\ref{thm:classification}, which needs no large cardinals. Necessity: for $\pi\in\Gamma$ and an increasing cofinal $a$ put $a^\pi_i(a)=\{\omega\cdot a_n+\pi^{-n}(i):n<\omega\}$ (Lemma~\ref{lem:perm} with $a$ in place of the critical sequence). Inserting one point into $a$ gives $[a^\pi_i(d)]=[a^\pi_{\pi^{-1}(i)}(a)]$, i.e.\ $\vec q^{\,\pi}(d)=\pi\cdot\vec q^{\,\pi}(a)$. For a rule $L$, the name $T(c)=L(\vec q^{\,\pi}(c))$ satisfies $T(d)=\pi\cdot T(a)$, so by Lemma~\ref{lem:insertion}(c) its value is fixed by $\pi$. For a finite family of rules, apply the same to the finite \emph{set} of values, with $\rho$ the induced permutation of $\Pow(F)$, and then argue as in Lemma~\ref{lem:algebra}.
\end{proof}

\begin{theorem}[No finite definable family of selectors]\label{thm:prikrysel}\src{R28, R30--R36}
For $1\le r<n<\omega$ there is no nonempty finite $\mathcal F\subseteq\Sel_{n,r}(Q_\kappa)$ in $\mathcal D(W,q)$, also with tail-invariant parameters. Hence there is no nonempty finite definable family of linear orders of $Q_\kappa$, of tournaments on $Q_\kappa$, or of choice functions on its finite subsets, and no definable injection of $Q_\kappa$ into the ordinals.
\end{theorem}

\begin{proof}
Let $|\mathcal F|\le m$, $L=\operatorname{lcm}(1,\dots,m)$, $N=nL$ and code the $N$-cycle $\tau$ into an $N$-tuple of classes as above. The finite observation is the set of \emph{restriction tables} $v_f:[N]^n\to[N]^r$, $f\in\mathcal F$; insertion acts on it by $(\tau\cdot v)(B)=\tau[v(\tau^{-1}[B])]$. By Lemma~\ref{lem:insertion}(c) the set of tables is $\tau$-invariant, which contradicts Lemma~\ref{lem:algebra}.
\end{proof}

Theorems~\ref{thm:prikryclass} and~\ref{thm:prikrysel} answer Question~(ii) of the third edition positively for the Prikry model. They do not show that rigidity alone implies these conclusions: the proof uses the forcing presentation, not only the regularity of $\kappa$ in $N_q$.

\subsection{Ultrafilters, kernels and charges on the generic class}

For $a\in q$ let $t_a:\mathbb Z\to q$ (defined on a final segment) send $k$ to $a$ with its first $k$ points removed, extended coherently; for $a,b\in q$, $t_b(n)=t_a(n-\ch_a(b))$ for large $n$.

\begin{theorem}[The local threshold is $\aleph_0$]\label{thm:ufthreshold}\src{R28, R30, R33--R36}
\begin{enumerate}
\item \emph{In $\ZFC$, for every uncountable $\lambda$ of cofinality $\omega$.} Fix a nonprincipal ultrafilter $D$ on $\omega$. Then $\mathcal U_D(q)=\{(t_a)_*D:a\in q\}$ is a countably infinite family of nonprincipal ultrafilters on $q$, each containing every tail ray; $q\mapsto\mathcal U_D(q)$ is definable from $D$. The family is a $\mathbb Z$-torsor: it carries a canonical successor with one infinite orbit and no designated origin. (No shift $S_j$, $j\ne0$, fixes an ultrafilter on $\mathbb Z$.)
\item \emph{In the Prikry model} there is no nonempty finite family of ultrafilters on $q$ in $\mathcal D(W,q)$, and no nonempty finite family of total ultrafilter-valued kernels on $Q_\kappa$. So the least size of a nonempty $\OD_{V_\kappa\cup\{q\}}$ family of ultrafilters on $q$ is exactly $\aleph_0$, while the least size of a nonempty definable family of \emph{members} of $q$ is $\kappa$.
\item \emph{Profinite density} \R{36}. For any nonempty $\mathcal W\in\mathcal D(W,q)$ of ultrafilters on $q$, the set $H$ of phases of its members in $\widehat{\mathbb Z}$ satisfies $H+1=H$; so $H$ is dense, $\mathcal W$ is infinite, and if $\mathcal W$ is compact then $H=\widehat{\mathbb Z}$ and $|\mathcal W|\ge2^{2^{\aleph_0}}$. The closure of $\mathcal U_D(q)$ attains this bound.
\item \emph{Charges} \R{36}. If $\mathcal M\in\mathcal D(W,q)$ is a family of $h<\omega$ charges on $q$, every member gives mass $1/\ell$ to every residue class of the index modulo $\ell$, for every prime $\ell>h$ (for all $\ell$ if $h=1$); the bound $\ell>h$ is sharp, by $d$-element families of periodic ray averages.
\end{enumerate}
\end{theorem}

\begin{proof}
(2),(3): the phase of an ultrafilter is a finite observation modulo each $n$; insertion adds $1$ to it. A finite nonempty subset of $\widehat{\mathbb Z}$ is not invariant under $+1$. (4): insertion permutes $\mathcal M$ with orbits of length $t\le h$, so each member is invariant under the shift by $t$ of the residues modulo $\ell$, and $t$ is invertible modulo a prime $\ell>h$; compare Theorem~\ref{thm:charges}(c).
\end{proof}

\begin{theorem}[The finite--countable gap for total objects]\label{thm:gap}\src{R34}
Let $d\in W$ be a well-ordered presentation of the $\PP_U$-names for subsets of $\kappa$, and $\nu\in W$ a nonprincipal ultrafilter on $\omega$. For all $1\le r<n$ there is a countably infinite family of \emph{total} members of $\Sel_{n,r}(Q_\kappa)$ definable from $d,\nu,q$, and a countably infinite definable family of total ultrafilter-valued kernels; no member of either family is in $\mathcal D(W,q)$. Moreover $q$ is definable from $U$ (as the $\subseteq^*$-largest class of generic sequences). After coding $d,\nu,U$ into the continuum function high above $\kappa$, one obtains a model in which these families are ordinal definable, $V_\kappa\subseteq\HOD$, $\HOD\models$ ``$\kappa$ is measurable'', $q$ is ordinal definable and rigid, and the least size of a nonempty $\OD_{V_\kappa}$ family of $(n,r)$-selectors, and of kernels, is exactly $\aleph_0$.
\end{theorem}

\begin{proof}
Each $b\in q$ is a generic sequence with $W[b]=V$, so $d$ and $b$ define a well-order of $V_{\kappa+1}$ and hence a total selector $f_b$ and a total kernel $K_b$. The values lie in finite sets (respectively, are decided set by set), so $f_{a,k}(x)=\nu\text{-}\lim_j f_{t_a(j+k)}(x)$ is defined; by synchronization of tails, replacing $a$ by $b\in q$ shifts $k$ by $\ch_a(b)$, uniformly in the input $x$. So $\{f_{a,k}:k\in\mathbb Z\}$ depends only on $q$. It is infinite by Theorem~\ref{thm:prikrysel}.
\end{proof}

This is the first positive result of its kind: between ``no finite definable family'' and ``a definable family of size $\kappa$'' the truth, in the Prikry model and with a ground parameter, is $\aleph_0$. At an ultraexacting cardinal the question is open, and \R{28}, \R{30}, \R{33}, \R{35}, \R{36}, which allow only low-rank parameters, correctly state that their local construction does not settle it.

\subsection{Strength and separation}

\begin{theorem}\label{thm:prikrystrength}\src{R28, R30--R36}
Each of the following packages is equiconsistent with one measurable cardinal: the conclusions of Theorems~\ref{thm:prikryrigid}--\ref{thm:ufthreshold} together with an internal normal tail measure in $\HOD_{\{b,q\}}$ (every report); the same with $2^\lambda$ jointly rigid classes and rank capture (\R{31}, \R{32}, \R{35}); the package of Theorem~\ref{thm:gap} (\R{34}). In every case the lower bound comes from the clause asserting an internal normal measure (or measurability in $\HOD$); no lower bound is known for the finite-symmetry clauses alone.
If $\kappa$ is the least measurable cardinal of $W$, then in $W[c]$ all of this holds at $\lambda=\kappa$ while there is no nontrivial elementary $e:V_\lambda\to V_\lambda$, no exacting or ultraexacting cardinal $\le\lambda$, and the core regards $\lambda$ as its least measurable cardinal, so its normal tail measure does not concentrate on measurables (compare Theorem~\ref{thm:strength}(b)) \R{30}, \R{32}, \R{36}.
\end{theorem}

\begin{assessment}{What the fourth round changes}
After this round, no conclusion of Sections~\ref{sec:uex}, \ref{sec:quot} and~\ref{sec:measures} about definable choice on the quotient $Q_\lambda$ is known to need more than one measurable cardinal, except those that mention the embedding or $\lambda^+$ explicitly (Theorem~\ref{thm:succmeas}, Theorem~\ref{thm:vopenka}, the $\Izero$ package). Conversely the Prikry model has more: $2^\kappa$ rigid classes instead of $\lambda$, and ground parameters of arbitrary rank. The results of groups C and F are therefore best described as theorems about Prikry-type singularization, of which an ultraexacting embedding provides an instance inside $V$ without forcing (Theorem~\ref{thm:core}(3)). Whether they follow from rigidity alone remains open.
\end{assessment}

\section{Countable analogues}\label{sec:countable}

This section is not drawn from the reports. It records what the results of Sections~\ref{sec:width}--\ref{sec:prikry} become when $\lambda$ is replaced by a countable ordinal of cofinality $\omega$ and ``ordinal definable'' by a notion from descriptive set theory or higher recursion theory. The barrier, the ground theorems and everything about $\HCD(\eta)$ have no such analogue: they are statements about covering in $V$. Nor has the threshold $\aleph_0$ of Theorem~\ref{thm:ufthreshold}, which needs an ultrafilter on $\omega$ as a parameter. Proof-theoretic ordinals are not touched by anything here.

\subsection{Admissibility as rigidity}

Code a cofinal $\omega$-sequence in $\omega_1^{\mathrm{CK}}$ by a real $x=\langle e_n:n<\omega\rangle$ with $e_n\in\mathcal O$ and $\sup_n|e_n|=\omega_1^{\mathrm{CK}}$; let $C$ be the set of such codes.

\begin{proposition}\label{prop:ck}
Every $\Sigma^1_1$ subset of $C$ is empty or contains a perfect set. The same holds for $\Sigma^1_1(z)$ sets whenever $\omega_1^z=\omega_1^{\mathrm{CK}}$.
\end{proposition}

\begin{proof}
A $\Sigma^1_1$ set without a perfect subset is countable and all its members are hyperarithmetic \cite[Ch.~III]{sacks}. If $x=\langle e_n\rangle\in C$ were hyperarithmetic, $\{e_n:n<\omega\}$ would be a $\Delta^1_1$ subset of $\mathcal O$, hence bounded below $\omega_1^{\mathrm{CK}}$ by Spector's $\Sigma^1_1$-boundedness theorem.
\end{proof}

This has exactly the shape of Theorem~\ref{thm:smallC} and Definition~\ref{def:rigid}: a definable family of short cofinal sets is empty or as large as possible, with ``countable'' in place of ``of size below $\lambda$'', lightface $\Sigma^1_1$ in place of $\OD_{V_\lambda}$, and admissibility of $\omega_1^{\mathrm{CK}}$ --- the absence of a cofinal $\omega$-sequence that is $\Sigma_1$ over $L_{\omega_1^{\mathrm{CK}}}$ --- in place of the regularity of $\lambda$ in $\HOD_{V_\lambda}$ (Theorem~\ref{thm:rigreg}). It is a routine consequence of classical theorems and presumably folklore; its interest is that it identifies the rigidity phenomenon of Section~\ref{sec:quot} with $\Sigma^1_1$-boundedness one universe up. The class-by-class statement of Corollary~\ref{cor:rigid}(1) degenerates, because a class modulo finite difference is now countable; a nonempty $\Sigma^1_1$ subset of $C$ instead meets uncountably many classes.

\subsection{The classical counterpart of the Prikry core}

Theorem~\ref{thm:core} should be compared with the theorem of Bukovsk\'y and Dehornoy \cite{bukovsky,dehornoy}: if $M_n$ is the $n$th iterate of $M_0$ by a normal measure $U$ with critical points $\kappa_n$, then $\langle\kappa_n:n<\omega\rangle$ is Prikry generic over the direct limit $M_\omega$ for the image of $U$, and $M_\omega[\langle\kappa_n\rangle]=\bigcap_nM_n$. Theorem~\ref{thm:core} is the form this takes when the iteration is replaced by the single embedding of an ultraexacting witness and $M_\omega$ by $\HOD_{\{d,q\}}$; the synthesis should have cited it in the third edition. When $M_0$ is a countable iterable structure (for instance the mouse $0^\dagger$), all ordinals involved are countable, and this is where the countable counterpart of Section~\ref{sec:measures} lives. One level lower, the tail filter of the Silver indiscernibles below $i_\omega$ is a normal $L$-ultrafilter on $i_\omega$ that does not belong to $L$; there is a measure but no genericity.

\begin{question}
In the situation of Theorem~\ref{thm:core}, is the Prikry core $K(d,q)=N(d,q)[a]$ an intersection of models, as in the Bukovsk\'y--Dehornoy theorem --- for instance of the models $\HOD_{\{d,\,e^n``a\}}$, or of the images under the iterates of $e$ of a suitable model?
\end{question}

\subsection{$\lambda=\omega$: the Baire-measurable classification}

For $\lambda=\omega$ the quotient $Q_\lambda$ is the set of infinite subsets of $\omega$ modulo finite difference, a form of $E_0$. Give $X_m=\Pow(\omega)^m$ the product topology and let $D_m\subseteq X_m$ be the set of $m$-tuples of infinite sets that are pairwise different modulo finite sets; $D_m$ is comeagre and invariant under finite changes and under permutations of the coordinates. For $\Gamma\le S_m$ and a finite $\Gamma$-set $F$, call $L:D_m\to F$ \emph{admissible} if it is invariant under finite changes and $L(\pi\cdot\vec x)=\pi\cdot L(\vec x)$ for $\pi\in\Gamma$, as before.

\begin{proposition}\label{prop:baire}
A Baire-measurable admissible $L:D_m\to F$ exists if and only if some label in $F$ is fixed by every element of $\Gamma$. Consequently, in $\ZF+\mathsf{DC}+{}$``every set of reals has the Baire property'' this is the criterion for the existence of any admissible rule.
\end{proposition}

\begin{proof}
A constant map to a globally fixed label is admissible. Conversely, finite changes act on $X_m$ by homeomorphisms and topologically transitively, so by the topological zero-one law \cite[Section~8]{kechris} each fibre of $L$ is meagre or comeagre; as $F$ is finite, $L$ has a constant value $\ell_0$ on a comeagre set. A permutation of coordinates is a homeomorphism, so for comeagre many $\vec x$ both $L(\vec x)=\ell_0$ and $L(\pi\cdot\vec x)=\ell_0$, whence $\pi\cdot\ell_0=\ell_0$.
\end{proof}

Compare Theorems~\ref{thm:classification} and~\ref{thm:prikryclass}, where the criterion is the strictly weaker $\CF(\Gamma,F)$: every single $\pi\in\Gamma$ fixes a label. The two differ already for the Klein four-group acting on three blocks of two labels, each non-identity element fixing one block pointwise. The difference is located in the sufficiency half of Theorem~\ref{thm:classification}: it chooses, for each $\Gamma$-orbit of tail-classes of incidence words, one class on which to start, and does so from a well-order of $\Pow(\omega)$ --- a parameter in $V_\lambda$. That is a selection from a quotient of $E_0$ type and cannot be Baire measurable. The necessity halves of Sections~\ref{sec:uex}, \ref{sec:quot} and~\ref{sec:prikry} only ever test one $\pi$ at a time (Lemma~\ref{lem:insertion}(c)), which is why they yield $\CF$ and no more. So the classification measures how much choice the definability class contains, and the two known answers are the extremes.

\begin{question}
(i) Let $\lambda$ be ultraexacting, or let $V=W[c]$ be a Prikry extension, and allow as parameters only ordinals, reals and the class $q$, in a universe where no well-order of the reals is ordinal definable from a real (for instance after adding $\omega_1$ Cohen reals; large cardinals alone do not ensure this). Is the criterion for a definable admissible rule $\CF(\Gamma,F)$, a global fixed point, or something in between? (ii) What is the criterion in a Prikry extension of a model of $\mathsf{AD}+\mathsf{DC}$?
\end{question}

Proposition~\ref{prop:baire} is formalized in Lean, in the generality of a finite-change-invariant labelling of $\{0,1\}^\iota$ that is generically equivariant for a group acting on the index set and on the labels, together with the zero-one law it rests on and the Klein four-group example (\nolinkurl{Cardinals/Countable/GenericErgodicity.lean}; nothing is admitted there).

A further connection, through the virtual large cardinals of Gitman and Schindler, is worked out in the separate research note \nolinkurl{docs/cardinals/research-notes/Exacting_Embeddings_of_Countable_Structures}: exacting witnesses that exist only in a forcing extension, equivalently between countable structures $L_\zeta$, are compatible with $V=L$; they separate the consequences of a single witness from those of the Kunen inconsistency, and they define a family of large countable ordinals below the least stable ordinal.

\section{Comparison of the reports}\label{sec:compare}

\subsection{Agreement}

The seven group-A reports agree on every mathematical conclusion they share. All nine agree on status: no unconditional inconsistency, no resolution of Problem~5.2, no priority claim, no formal verification. No contradiction between any two reports was found.

\subsection{Differences that matter}

\begingroup\small
\renewcommand{\arraystretch}{1.18}
\begin{longtable}{@{}P{.2\textwidth}P{.75\textwidth}@{}}
\toprule
\textbf{Topic}&\textbf{Finding}\\\midrule\endfirsthead
\toprule
\textbf{Topic}&\textbf{Finding}\\\midrule\endhead
Dependence on lecture notes&\R1--\R3, \R8 import the stationary characterization (Prop.~3.4) and prove a focused transfer lemma; \R7 imports both Lemma~3.3 and Prop.~3.4; \R4 and \R6 prove persistent covers, exact hulls and general transfer from the definition. Route~H therefore removes the barrier's dependence on unrefereed results other than definitions, and should be preferred.\\
Threshold&$\eta>\gamma$ (\R1, \R3, \R8), $\gamma^+$ (\R2, \R4, \R7), $q(\gamma)$ (\R6). Only \R6 notices the singular endpoint, which gives the clean statement ``$\CEx_\lambda(\lambda)\Rightarrow\lambda$ regular in $\HCD(\lambda)$''.\\
Extra hypothesis below $\lambda$&Seven formulations $\mathsf P_1$--$\mathsf P_7$; all are instances of one argument, and the two size conventions coincide (Corollary~\ref{cor:cond}). The weakest uniform principles are $\mathsf P_4$ (\R2) and $\mathsf P_5$ (\R6), which are not comparable on their face.\\
Forcing lower bound&\R1, \R2 bound $|\PP|$; \R3, \R4, \R6, \R8 prove failure of the $\lambda$-chain condition, which is stronger and also yields density $\ge\lambda$. \R4 stresses the bound is $\lambda$, not $\lambda^+$.\\
Ground profile&Only \R8 proves that $\lambda$ is strongly inaccessible (not merely regular) in the ground; the same argument is reused for $\HOD$ in Theorem~\ref{thm:i0}(3).\\
Finite choice&\R5 and \R7 use different interfaces. \R5 (witnesses with high critical point and $j\restr V_\lambda\in X$) gets finite families and parameters in $V_\lambda$; \R7 (one preserved $\OD$ predicate on $V_{\lambda+1}$) gets the general label theorem and the transversal theorem, with ordinal parameters only. \R7's selector corollary is the singleton case of \R5's theorem; \R7's other results are not covered by \R5.\\
$\Izero$ packages&Five overlapping packages ($\mathsf{Bdry}$, $T_{\mathrm{first}}$, $\mathrm{Profile}$, $T_1$, $T_{\mathrm{bd}}$); all are corollaries of $\UEx(\lambda)\wedge V_\lambda\subseteq\HOD$ and merge into Theorem~\ref{thm:i0}.\\
Stationary track&Consistent: three reports discard the HKP reconstruction as known; \R3 also credits it to HKP and contributes only the strengthening, Theorem~\ref{thm:range}.\\
Citations&The $\UEx$/$\Izero$ equiconsistency is cited as ``Theorem~A'' (\R3, \R7) and ``Theorem~4.5'' (\R2, \R5, \R8) of \cite{abgl}; the $\delta$-cover theorem as Thm.~4.14, with Cor.~4.7 added in \R7; the trace argument as ``proof of Thm.~3.7'' or ``p.~10'' of \cite{cover}. These appear to be the same items under introduction and body numbering, but were not re-checked against the sources here.\\
\bottomrule
\end{longtable}
\endgroup

\subsection{The second round}

The nine reports \R{10}--\R{18} were written independently from the first edition of this synthesis, and all nine prove Corollary~\ref{cor:sharpwidth}(2). No contradiction among them, or with the first round, was found. They differ as follows.

\begingroup\small
\renewcommand{\arraystretch}{1.18}
\begin{longtable}{@{}P{.2\textwidth}P{.75\textwidth}@{}}
\toprule
\textbf{Topic}&\textbf{Finding}\\\midrule\endfirsthead
\toprule
\textbf{Topic}&\textbf{Finding}\\\midrule\endhead
Local lemma&Three equivalent devices turn a small fixed family into a small fixed set of ordinals: the union (\R{13}, \R{16}; only for families of countable sets), coordinate projections (\R{12}, \R{14}), and compression to the least order type (\R{10}, \R{11}, \R{18}; works for arbitrary short cofinal sets). Only the last gives Lemma~\ref{lem:fixedsmall}(b) in full.\\
Fixing the definable object&\R{10}, \R{12}--\R{17} raise the critical point above the rank of the parameter (Input~\ref{in:abl}(ii)); \R{11}, \R{13} use the predicate form (Input~\ref{in:abl}(iii)); \R{18} minimizes the rank of the parameter, which needs neither and is uniform in the parameter. \R{17} achieves uniformity differently, through a fixed bijection $\lambda\to V_\lambda$.\\
How many classes&One (\R{10}, \R{11}, \R{16}, \R{18}); $\kappa$ offsets per witness, hence $\lambda$ (\R{13}, \R{14}, \R{17}); $\lambda$ orbits of one witness (\R{12}, \R{15}). \R{15} shows $\lambda$ is the most one witness can give.\\
Strength of the statement&Single section (\R{11}, \R{13}, \R{16}); small families of complete sections (\R{10}, \R{12} with a uniform bound, \R{14}, \R{17}); small families of partial sections (\R{15}); all $\OD_{V_\lambda}$ sets at once (\R{17}); definability from $q$ itself (\R{11} for single sets, \R{18} for families). Theorem~\ref{thm:rigid} and Corollary~\ref{cor:rigid} are the common upper bound. \R{12}'s Proposition on small clouds is the weakest form and is subsumed.\\
Beyond cardinality&Only \R{12} analyses the inside of a full fibre (coordinates, phases). Only \R{13} treats finite labels, and it is the only second-round report that answers a question posed in the first edition. Only \R{14} returns to cover exactingness. Only \R{10} treats measures.\\
$\Izero$ packages&Nine more, all corollaries of $\UEx(\lambda)\wedge V_\lambda\subseteq\HOD$; merged into Theorem~\ref{thm:i0}(6).\\
Relation to the Lean development&The abstract theorem \nolinkurl{no_finite_valued_transversal} is the finite case of Lemma~\ref{lem:fixedsmall}(b); the second-round results are not yet formalized.\\
\bottomrule
\end{longtable}
\endgroup

\subsection{The third round}

\begingroup\small
\renewcommand{\arraystretch}{1.18}
\begin{longtable}{@{}P{.2\textwidth}P{.75\textwidth}@{}}
\toprule
\textbf{Topic}&\textbf{Finding}\\\midrule\endfirsthead
\toprule
\textbf{Topic}&\textbf{Finding}\\\midrule\endhead
Consistency with earlier rounds&No report contradicts an earlier one, but \R{24} changes the status of several earlier statements: Corollary~\ref{cor:sandwich}, Corollary~\ref{cor:statapps}(b) (\R{14}), the $\HCD$ and $\HOD$ applications of \R{20}, and the ``two-sided'' clauses of several first-round reports describe a situation that Theorem~\ref{thm:sandwich} shows impossible. \R{14}, \R{20} and \R{23} each came within one step of it: they derived collapse-or-destroy alternatives and chain-condition bounds from the same reflection argument, but did not use amenability of the $\omega$-club filter, which is what converts reflection into covering.\\
The normal-measure theorem&Proved four times (\R{22}, \R{23}, \R{26}, \R{27}) with the same two lemmas (tail decisions; regressive functions). They differ in the ground model: $\HOD_{\{q\}}$ and $\HOD_{V_\lambda\cup\{q\}}$ (\R{22}, \R{26}), $\HOD_{\{p,q\}}$ uniformly in $p\in V_\lambda$ with coherence (\R{23}), $\HOD_{\{b,q\}}$ with the fixed bijection (\R{26}, \R{27}). Only \R{27} proves concentration on measurables; only \R{22} gives the intrinsic criterion and the three-class example; \R{23} and \R{26} give the successor-shift example.\\
Strength&\R{25} and \R{27} independently use the same cone-homogeneity observation (Prikry cone isomorphisms preserve the class of the generic sequence). \R{25} calibrates rigidity, \R{27} the trace principles; together they give Theorem~\ref{thm:strength}. No other report in any round determines an exact consistency strength that is not inherited from $\Izero$.\\
Overlap with the literature&Theorem~\ref{thm:succmeas} refines \cite[Cor.~3.15]{abl}, which \R{19} cites; Theorem~\ref{thm:groundcover} is a decoupled form of Goldberg's seed argument \cite{dichotomy}, and \R{24} cross-checks its use against the published lemma. The Prikry facts (Mathias criterion, strong Prikry lemma) are reproved in appendices of four reports and are classical.\\
Charges&Only \R{21}. It complements \R{10} (countably additive: impossible) and \R{12} (phases), and its phase-balance theorem is the measure-theoretic form of Theorem~\ref{thm:phase}.\\
$\Izero$ packages&Eight more; merged into Theorem~\ref{thm:i0}(7).\\
\bottomrule
\end{longtable}
\endgroup

\subsection{The fourth round}

\begingroup\small
\renewcommand{\arraystretch}{1.18}
\begin{longtable}{@{}P{.2\textwidth}P{.75\textwidth}@{}}
\toprule
\textbf{Topic}&\textbf{Finding}\\\midrule\endfirsthead
\toprule
\textbf{Topic}&\textbf{Finding}\\\midrule\endhead
Convergence&Eight reports prove Theorems~\ref{thm:prikryclass} and~\ref{thm:prikrysel} with the same two ingredients: one-point surgery on the generic sequence (insertion in \R{30}, \R{31}, \R{34}--\R{36}; deletion in \R{28}; splicing in \R{32}; tail recoding in \R{33}) and the finite algebra of Lemma~\ref{lem:algebra} with the $\operatorname{lcm}$ amplification. All credit the sufficiency half and the amplification to the synthesis. No report contradicts another.\\
Parameters&All allow arbitrary ground-model sets, which is more than $\OD_{V_\kappa\cup\{q\}}$; \R{32} and \R{35} add tail-invariant parameters; \R{33} and \R{36} treat finite \emph{families} of rules. None allows a predicate for $W$ (\R{30}, \R{34} say so explicitly); since the cone isomorphisms fix $W$, and a ground is definable from a parameter, this is bookkeeping rather than a restriction.\\
Ultrafilters&Six reports construct the same countable family $\mathcal U_D(q)$ under six names. \R{36} adds the profinite and compact thresholds and the residue laws; \R{35} the torsor structure. Only \R{34} constructs \emph{total} selectors and kernels, using a ground parameter of high rank; the other five note that their construction is local. These statements are compatible.\\
Rigid classes&\R{31}, \R{32}, \R{35} obtain $2^\kappa$ rigid classes by three different codings of ground subsets of $\kappa$ into recodings of $c$; the conclusions agree.\\
Strength&Every report states the equiconsistency with one measurable for a package that \emph{includes} an internal normal measure, and every report disclaims a lower bound for finite symmetry alone. This discipline was absent from some first-round packages.\\
\R{29}&Independent of the others: it strengthens Theorem~\ref{thm:core} at an ultraexacting cardinal. Its restriction to predicates in the core is essential and is proved to be (Proposition~9.1 there).\\
\bottomrule
\end{longtable}
\endgroup

\subsection{Relative reliability}

The group-A, B and E arguments are short and were re-derived during the merge without difficulty; the only external mathematics they use are Inputs~\ref{in:kunen}, \ref{in:hcd} and \ref{in:abl}(i). The group-C arguments are elementary once the two imported interfaces are granted; \R5 and \R7 include Python scripts that exhaustively check only the finite algebra (Lemma~\ref{lem:algebra}, Lemma~\ref{lem:perm}, the $S_3$ example) and are not evidence for the set-theoretic steps. Theorem~\ref{thm:i3} is the most delicate result: it is the only new forcing construction, it depends on the precise form of the ABGL well-founded-part localization (Input~\ref{in:iter}, Lemma~\ref{lem:local}), and it appears in a single report. Its two new lemmas (\ref{lem:internal}, \ref{lem:cloud}) were re-derived here and appear sound; the imported interface was not checked. It should be the first target of specialist review. The second-round arguments are short and uniform: everything rests on Lemma~\ref{lem:fixedsmall}, Lemma~\ref{lem:orbits} and the canonical-counterexample device, all re-derived here. The two places that needed care are the choice of a height correct for the definability class \emph{before} the witness (and hence the fixed class) is chosen, which \R{17} and \R{18} treat explicitly, and the rank bookkeeping in Theorem~\ref{thm:preserved}(a). Lemma~\ref{lem:cyclic} is elementary and was checked in full. \R{12} and \R{13} include Python checks of finite combinatorics only. In the third round, Theorem~\ref{thm:sandwich} deserves the closest scrutiny because it is the only outright non-coexistence theorem: its proof was re-derived here in both forms (the seed argument, and the two-line deduction from \cite[Lemma~2.7]{dichotomy}); the points to check are that the ground is definable from a set parameter and satisfies Choice, and that amenability refers to the ambient filter; the statement of the cited lemma was checked against arXiv:2107.00513v1. The Lean development derives the theorem from that lemma, with Lemma~\ref{lem:amenable} as a hypothesis (\nolinkurl{Cardinals/Sandwich.lean}). Since the fourth edition the Lean development also works with the actual embedding of a witness rather than with its combinatorial shadow: from the definition of an (ultra)exacting witness it proves Lemma~\ref{lem:crit} (given the cofinality of the critical sequence, which is admitted with a reference), Lemma~\ref{lem:fixedsmall}, the membership $a_r\in X$ and $j(a_r)=a_r\setminus\{r\}$ of Lemma~\ref{lem:orbits}, the zero-or-full theorem for sets fixed by $j$ (Theorem~\ref{thm:rigid}, Corollary~\ref{cor:rigid}(1)), and Theorem~\ref{thm:taildec} (\nolinkurl{Cardinals/FixedSets.lean}, \nolinkurl{Cardinals/Orbits.lean}). The passage from ordinal definability to fixedness by $j$ is not formalized. It inherits the dependence of Theorem~\ref{thm:barrier} on the Blue--Goldberg notes. Theorems~\ref{thm:taildec} and~\ref{thm:core} are short and were re-derived in full; Theorem~\ref{thm:succmeas} and Proposition~\ref{prop:vopenka} were checked only in outline. The lower bounds in Theorem~\ref{thm:strength} rest on classical core-model theorems quoted, not reproved, by \R{25} and \R{27}. In the fourth round, Lemma~\ref{lem:insertion} and the two main theorems were re-derived in full and are short; their finite algebra is the part of the first round that is machine-checked (\nolinkurl{Cardinals/Combinatorics/FiniteCycles.lean}, \nolinkurl{FiniteLabel.lean}). Theorem~\ref{thm:gap} was checked in outline: the point to verify is that a name presentation $d\in W$ and a generic $b$ define a well-order of $V_{\kappa+1}$ uniformly in $b$, and that the $\nu$-limit is taken of finitely-valued data. The coding construction behind its ordinal-definable version, Theorem~\ref{thm:prikryrigid}(4), and the compactness count in Theorem~\ref{thm:ufthreshold}(3) were not re-derived. Classical facts about Prikry forcing (the Prikry property, Mathias' criterion, preservation) are quoted by all reports from \cite{prikry,mathias}.

\section{What remains open}\label{sec:open}

\begin{question}[Promotion; all group-A reports]
Assume $\delta<\lambda\le\gamma$, $\SC(\delta)$, $\CEx_\gamma(\lambda)$, and (by Theorem~\ref{thm:sandwich}, necessarily) no strongly compact cardinal above $\gamma$. Do these hypotheses (or a named strengthening of strong compactness short of extendibility) imply $\mathsf P_4$, $\mathsf P_5$, or $\mathsf P_7$ for one of the sets $a$ of Theorem~\ref{thm:sep}? A positive answer refutes the configuration of Problem~5.2. The entire unproved step is raising the completeness of an ultrafilter code from $\delta$ to $q(\gamma)$; strong compactness extends filters at their given completeness and does not do this.
\end{question}

\begin{question}[Positive side]
Can a cover-exacting construction preserve a prescribed strongly compact $\delta$ below the cover-exacting cardinal? Any such model must exhibit the separation of Theorem~\ref{thm:sep}, must fail $\mathsf P_1$--$\mathsf P_6$, and, if it has a strongly compact above $\gamma$, must have the proper ground of Theorem~\ref{thm:ground}; it cannot satisfy $\GA$ in that case.
\end{question}

\begin{question}[\R9]
Can the $\Iwf$ upper bound for $\exists\lambda\,\CEx_\lambda(\lambda)$ be lowered, or does cover exactingness carry an iteration lower bound beyond that of exactingness? Under a smaller strongly compact, is there an $\OD_{V_\lambda}$ family of fewer than $\lambda$ candidates containing a short cofinal map (which would contradict Corollary~\ref{cor:odcover})?
\end{question}

\begin{question}[\R5, \R7]
Does ordinary exactingness already exclude an $\OD_{V_\lambda}$ binary selector on $Q_\lambda$? What is the least size of a nonempty $\OD_{V_\lambda}$ subfamily of $\Sel_{n,r}(Q_\lambda)$ at an ultraexacting $\lambda$; can it be countable? Does a canonical Icarus enrichment define a preserved selector, thin section, or map $Q_\lambda\to D_\lambda$ (Theorem~\ref{thm:preserved})? (The first edition also asked whether a modification-invariant equivariant rule can exist for a finite $S_m$-set with elementwise but no global fixed points; Theorem~\ref{thm:classification} answers yes.)
\end{question}

\begin{question}[\R{10}--\R{18}]
(i) Does ordinary exactingness exclude an $\OD_{V_\lambda}$ thin complete section? The local lemma needs only exactingness; the missing step is a fixed class in the domain. (ii) Can the thin complete sections, the kernels, or the maps $Q_\lambda\to D_\lambda$ have a nonempty $\OD_{V_\lambda}$ subfamily of size exactly $\lambda$? (iii) Must an $\OD_{V_\lambda}$ complete section have more than $\lambda$ full fibres, perhaps $2^\lambda$? (iv) Is there an intrinsic combinatorial characterization of the rigid classes, and are there more than $\lambda$ of them, or more than one model $N_q$? (v) In Corollary~\ref{cor:statapps}(b), must $(\lambda^+)^{\HCD(\kappa)}$ be collapsed?
\end{question}

\begin{question}[\R3]
Does $\mathrm{Card}^M=\mathrm{Card}^V$ for $j:V\to M$ yield an unbounded subset of $j``\theta$ in $M$, or amenability of the ambient nonstationary ideal? Either would refute cardinal-preserving embeddings by Theorem~\ref{thm:range}.
\end{question}

\begin{question}[\R{19}--\R{27}]
(i) Can Theorem~\ref{thm:sandwich} be proved without the upper strongly compact cardinal, for instance by showing that some $\HCD(\eta)$, $\eta\ge q(\gamma)$, is a ground, or by using Corollary~\ref{cor:sandwichcons}(d) against the exacting witnesses? (ii) Do the necessity half of Theorem~\ref{thm:classification} and Theorem~\ref{thm:fsel} hold in the Prikry extension of a measurable, i.e.\ do they follow from rigidity alone? (iii) Must an ultraexacting $\lambda$ have rigid classes with different models $N_q$, or different normal tail measures? (iv) Can a normal measure on $\lambda$ be found in $\HOD$ itself, without the parameter $q$? (v) Can ordinary exactingness produce a rigid class? (The equiconsistency with a measurable shows this cannot be decided by strength alone.) (vi) Can a countably infinite definable family of ultrafilter-valued kernels exist?
\end{question}

\begin{question}[\R{28}--\R{36}]
Question (ii) above is answered positively for the Prikry model (Section~\ref{sec:prikry}). Open: (i) Does rigidity of $q$ alone, or rigidity together with a normal tail measure in $\HOD_{\{b,q\}}$, imply finite-label necessity and the absence of finite definable families of selectors? (ii) What is the consistency strength of those two conclusions by themselves; can they be forced without a measurable? (iii) Is there a countably infinite $\OD_{V_\lambda}$ family of total selectors, or of total ultrafilter-valued kernels, at an ultraexacting $\lambda$; and in an arbitrary Prikry extension without the ground parameter $d$? (iv) Can an ultraexacting cardinal have $2^\lambda$ rigid classes? (v) What is the strength of a normal measure containing all predicate-extendibility sets $E_A$ (Theorem~\ref{thm:vopenka})?
\end{question}

Not addressed by any report: ordinary exacting cardinals above an extendible, the HOD Conjecture fork, HOD-Berkeley cardinals, and the choiceless hypotheses of the research plan.

\appendix
\section{Concordance}\label{app:concordance}

\yes\ = stated and proved in the report; \prt\ = used, imported, or proved only in a weaker or special form.

\begingroup\small
\renewcommand{\arraystretch}{1.12}
\begin{longtable}{@{}P{.47\textwidth}CCCCCCCCC@{}}
\toprule
\textbf{Result}&\R1&\R2&\R3&\R4&\R5&\R6&\R7&\R8&\R9\\\midrule\endfirsthead
\toprule
\textbf{Result}&\R1&\R2&\R3&\R4&\R5&\R6&\R7&\R8&\R9\\\midrule\endhead
Lemma~\ref{lem:crit} critical-sequence normalization&\yes&\yes&\yes&\yes&\prt&\yes&\yes&\yes&\\
Lemma~\ref{lem:short} short cofinal predicate&\yes&\yes&\prt&\yes&&\yes&\yes&\prt&\\
Lemma~\ref{lem:transfer} transfer (general form)&\prt&\prt&\prt&\yes&&\yes&\prt&\prt&\\
Lemma~\ref{lem:bound} $\gamma\ge\lambda$&&&&&&\yes&&&\\
Lemmas~\ref{lem:persist}--\ref{lem:hull} persistent covers, exact hulls&&&&\yes&&\yes&&&\\
Lemma~\ref{lem:trace} trace reconstruction&\yes&\yes&\yes&\yes&&\yes&\yes&\yes&\\
Lemma~\ref{lem:singular} singular endpoint&&&&&&\yes&&&\\
Theorem~\ref{thm:barrier}(2),(3) barrier&\yes&\yes&\yes&\yes&&\yes&\yes&\yes&\\
Theorem~\ref{thm:barrier}(1) absorption&&&&\yes&&\yes&&&\\
Theorem~\ref{thm:barrier}(4) cover gap&&\yes&&\yes&&\yes&&&\\
Corollary~\ref{cor:escape} sequences/embeddings escape&&&\prt&&&&&\yes&\\
Corollary~\ref{cor:cert} failed-stationarity certificate&\yes&&&&&&&&\\
Corollary~\ref{cor:approx} cover/approximation failure&\yes&\prt&\yes&\prt&&\prt&&\yes&\\
Theorem~\ref{thm:sep} separation below $\lambda$&\yes&\yes&\yes&\yes&&\yes&\yes&\yes&\\
\quad item (1) no small cover / (2) no refinement / (3) ultrafilter&&(1)&&(3)&&(2)&&&\\
Corollary~\ref{cor:cond} conditional inconsistency&\yes&\yes&\yes&\yes&&\yes&\yes&\yes&\\
Corollary~\ref{cor:ext} extendible recovery&\yes&\yes&\yes&\yes&&\yes&\yes&\yes&\\
Theorem~\ref{thm:first} first regularization&&&&&&\yes&&&\\
Lemma~\ref{lem:cc} ($\lambda$-cc form)&\prt&\prt&\yes&\yes&&\yes&\prt&\yes&\\
Theorem~\ref{thm:ground} proper ground&\yes&\yes&\yes&\yes&&\yes&&\yes&\\
\quad $\lambda$ strongly inaccessible in the ground&&&&&&&&\yes&\\
Corollary~\ref{cor:ga} Ground Axiom obstruction&\yes&\yes&\yes&\yes&&\yes&&\yes&\\
Corollary~\ref{cor:gengound} general / small-ground form&\yes&&\yes&&&&\yes&&\\
Corollary~\ref{cor:sandwich} two strongly compacts&\yes&&&&&\yes&&\yes&\\
Theorem~\ref{thm:width}, Corollary~\ref{cor:odcover} projection width&&&&&&&&&\yes\\
Proposition~\ref{prop:fresh} fresh cofinal set&&\prt&\prt&&&&&\prt&\yes\\
Theorem~\ref{thm:fsel} finite selector families&&&&&\yes&&\prt&&\\
Theorem~\ref{thm:icarus} Icarus enrichment&&&&&\yes&&&&\\
Theorem~\ref{thm:label} finite labels&&&&&&&\yes&&\\
Theorem~\ref{thm:transversal} finite-valued transversal&&&&&&&\yes&&\\
Theorem~\ref{thm:i3} $\Iwf\Rightarrow\CEx$&&&&&&&&&\yes\\
Corollary~\ref{cor:calib} calibration&\yes&&&&&\prt&&&\prt\\
Theorem~\ref{thm:i0} $\Izero$ package&&\yes&\yes&&\yes&&\yes&\yes&\\
Theorem~\ref{thm:range} internal unbounded range&&&\yes&&&&&&\\
HKP reconstruction audited as known&&\yes&\prt&&&\yes&\yes&&\\
\bottomrule
\end{longtable}
\endgroup

\subsection*{Second round}

\begingroup\small
\renewcommand{\arraystretch}{1.12}
\begin{longtable}{@{}P{.47\textwidth}CCCCCCCCC@{}}
\toprule
\textbf{Result}&\R{10}&\R{11}&\R{12}&\R{13}&\R{14}&\R{15}&\R{16}&\R{17}&\R{18}\\\midrule\endfirsthead
\toprule
\textbf{Result}&\R{10}&\R{11}&\R{12}&\R{13}&\R{14}&\R{15}&\R{16}&\R{17}&\R{18}\\\midrule\endhead
Lemma~\ref{lem:fixedsmall}(a) fixed small sets of ordinals&\yes&\yes&\yes&\yes&\yes&\yes&\yes&\yes&\yes\\
Lemma~\ref{lem:fixedsmall}(b) for arbitrary short cofinal sets&\yes&\yes&&&&&&&\yes\\
\quad for families in $D_\lambda$&\yes&\yes&\yes&\prt&\yes&\yes&\yes&\yes&\yes\\
Theorem~\ref{thm:smallC} (exacting)&\yes&&&&\prt&&&\prt&\yes\\
\quad minimal-rank-parameter device&&&&&&&&&\yes\\
Theorem~\ref{thm:pinning} bounded separation&\yes&&&&\yes&&&\yes&\\
Theorem~\ref{thm:laws}, Cor.~\ref{cor:kernel} laws and kernels&\yes&&&&&&&&\\
Theorem~\ref{thm:successor}, Cor.~\ref{cor:succcc}, \ref{cor:statapps}&&&&&\yes&&&&\\
Lemma~\ref{lem:orbits} roots and orbits&&&\yes&&&\yes&&&\\
\quad offsets $\kappa_n+t$&&&&\yes&\yes&&&\yes&\\
\quad critical-sequence class only&\yes&\yes&&&&&\yes&&\yes\\
Theorem~\ref{thm:rigid} rigidity from $q$ as parameter&&\prt&&&&&&&\yes\\
Cor.~\ref{cor:rigid}(1) for all $\OD_{V_\lambda}$ sets at once&&&&&&&&\yes&\yes\\
Cor.~\ref{cor:rigid}(2) small families, partial sections&&&&&\prt&\yes&&\prt&\yes\\
Cor.~\ref{cor:rigid}(4) the model $N_q$&&\yes&&&&&&&\prt\\
Cor.~\ref{cor:sharpwidth}(2) no thin complete section&\yes&\yes&\yes&\yes&\yes&\yes&\yes&\yes&\yes\\
\quad $\lambda$ full fibres&&&\yes&\yes&\yes&\yes&&\yes&\\
\quad small families of thin sections&\yes&&\prt&\prt&\yes&\yes&\prt&\yes&\yes\\
Remark~\ref{rem:sharpquot}(b) parameter-rank cutoff&&&&&&&&\yes&\\
Remark~\ref{rem:sharpquot}(d),(e) recoverable classes, recodings&&\yes&&&&&&&\\
Theorem~\ref{thm:phase}, Cor.~\ref{cor:phasecons} width and phase&&&\yes&&&&&&\\
Cor.~\ref{cor:colour}(a) colourings&&&&&&\yes&\yes&\yes&\\
Cor.~\ref{cor:colour}(b) coarser relations&&&&&\yes&&\yes&&\\
Cor.~\ref{cor:colour}(c) homomorphisms&&&&&&&&&\yes\\
Theorem~\ref{thm:conversion} cofinal-output rules&\yes&\yes&&&&&\yes&&\yes\\
Lemma~\ref{lem:cyclic}, Theorem~\ref{thm:classification}, Cor.~\ref{cor:five}&&&&\yes&&&&&\\
Theorem~\ref{thm:preserved}(a) $V_{\lambda+1}$ with predicate&&&&\yes&&&\yes&&\\
Theorem~\ref{thm:preserved}(b),(c) inner models, Icarus&\yes&\yes&\yes&\prt&&\yes&\yes&\yes&\yes\\
Proposition~\ref{prop:periodic} periodic classes; $|Q_\lambda|=2^\lambda$&&&&&&\yes&&&\\
Theorem~\ref{thm:i0}(6) $\Izero$ package&\yes&\yes&\yes&\yes&\yes&\yes&\yes&\yes&\yes\\
\bottomrule
\end{longtable}
\endgroup

\subsection*{Third round}

\begingroup\small
\renewcommand{\arraystretch}{1.12}
\begin{longtable}{@{}P{.47\textwidth}CCCCCCCCC@{}}
\toprule
\textbf{Result}&\R{19}&\R{20}&\R{21}&\R{22}&\R{23}&\R{24}&\R{25}&\R{26}&\R{27}\\\midrule\endfirsthead
\toprule
\textbf{Result}&\R{19}&\R{20}&\R{21}&\R{22}&\R{23}&\R{24}&\R{25}&\R{26}&\R{27}\\\midrule\endhead
Lemma~\ref{lem:amenable}, Theorem~\ref{thm:groundcover} amenable grounds cover&&&&&&\yes&&&\\
Theorem~\ref{thm:sandwich}, Cor.~\ref{cor:sandwichcons} no sandwich&&&&&&\yes&&&\\
Theorem~\ref{thm:cascade} cofinality cascade&&\yes&&&&&&&\\
Theorem~\ref{thm:successor} reused (collapse or destroy)&\prt&\prt&&&\yes&&&&\\
Lemma~\ref{lem:orbits} orbits; fixed bijection $b$&\yes&&\yes&\yes&\prt&&&\yes&\yes\\
Theorem~\ref{thm:rigreg} rigidity $\Leftrightarrow$ regularity in $N_q$; copies&&&\prt&\prt&&&\yes&&\\
Theorem~\ref{thm:taildec}(1) tail decisions&&&\prt&\yes&\yes&&&\yes&\yes\\
Theorem~\ref{thm:taildec}(2) regressive functions&&&&\yes&\yes&&&\yes&\yes\\
Theorem~\ref{thm:core}(1) $\lambda$ many measurable cores&&&&\yes&&&&\yes&\\
Theorem~\ref{thm:core}(2) normality; uniqueness&&&&\yes&\yes&&&\yes&\yes\\
\quad concentration on measurables&&&&&&&&&\yes\\
Theorem~\ref{thm:core}(3) Prikry genericity, common extension&&&&\yes&\yes&&&\yes&\yes\\
\quad uniform in $p\in V_\lambda$, coherence&&&&&\yes&&&&\\
Remark~\ref{rem:corelimits}(b) shift examples&&&&\yes&\yes&&&\yes&\\
Remark~\ref{rem:corelimits}(d) intrinsic criterion&&&&\yes&&&&&\\
Cor.~\ref{cor:split} definable splitting&&&&\yes&\yes&&&\yes&\prt\\
Theorem~\ref{thm:succmeas} $\lambda^+$ measurable in $N_F$&\yes&&&&&&&&\\
Theorem~\ref{thm:strength}(a) strength: a measurable&&&&&&&\yes&&\yes\\
Theorem~\ref{thm:strength}(b) strength: $o(\kappa)\ge2$&&&&&&&&&\yes\\
Proposition~\ref{prop:vopenka} Vop\v enka algebra&&&&&&&\yes&&\\
Theorem~\ref{thm:charges} charges&&&\yes&&&&&&\\
Theorem~\ref{thm:i0}(7) $\Izero$ package&\yes&&\yes&\yes&\yes&&&\yes&\yes\\
\bottomrule
\end{longtable}
\endgroup

\subsection*{Fourth round}

\begingroup\small
\renewcommand{\arraystretch}{1.12}
\begin{longtable}{@{}P{.47\textwidth}CCCCCCCCC@{}}
\toprule
\textbf{Result}&\R{28}&\R{29}&\R{30}&\R{31}&\R{32}&\R{33}&\R{34}&\R{35}&\R{36}\\\midrule\endfirsthead
\toprule
\textbf{Result}&\R{28}&\R{29}&\R{30}&\R{31}&\R{32}&\R{33}&\R{34}&\R{35}&\R{36}\\\midrule\endhead
Theorem~\ref{thm:vopenka} Vop\v enka reflection&&\yes&&&&&&&\\
Theorem~\ref{thm:core} reproved as prerequisite&&\yes&&&&&&&\\
Lemma~\ref{lem:insertion} stem surgery&\yes&&\yes&\yes&\yes&\yes&\yes&\yes&\yes\\
Lemma~\ref{lem:capture} ground capture&\prt&&\yes&\yes&\yes&\yes&\yes&\yes&\yes\\
\quad tail-invariant parameters&&&&&\yes&&&\yes&\\
Theorem~\ref{thm:prikryrigid}(1) rigidity, zero-or-full&\yes&&\yes&\yes&\yes&\yes&\yes&\yes&\yes\\
Theorem~\ref{thm:prikryrigid}(2) local constancy, phases&&&&&&\yes&&&\\
Theorem~\ref{thm:prikryrigid}(3) $2^\kappa$ rigid classes&&&\prt&\yes&\yes&&&\yes&\\
Theorem~\ref{thm:prikryrigid}(4) rank capture&&&&&\yes&&&&\\
Theorem~\ref{thm:prikryclass} classification&\yes&&\yes&\yes&\yes&\yes&\yes&\yes&\yes\\
\quad finite families of rules&&&&&&\yes&&&\yes\\
Theorem~\ref{thm:prikrysel} no finite selector family&\yes&&\yes&\yes&\yes&\yes&\yes&\yes&\yes\\
Theorem~\ref{thm:ufthreshold}(1),(2) threshold $\aleph_0$&\yes&&\yes&&&\yes&\yes&\yes&\yes\\
\quad torsor structure&\prt&&&&&&&\yes&\prt\\
Theorem~\ref{thm:ufthreshold}(3) profinite density, compact families&&&&&&&&&\yes\\
Theorem~\ref{thm:ufthreshold}(4) residue laws for charges&&&&&&&&&\yes\\
Theorem~\ref{thm:gap} countable total families&&&&&&&\yes&&\\
Internal normal measure, common Prikry layer&\yes&&\yes&\yes&\yes&&\prt&\yes&\yes\\
Theorem~\ref{thm:prikrystrength} equiconsistency package&\yes&\prt&\yes&\yes&\yes&\prt&\yes&\yes&\yes\\
\quad separation at the least measurable&&&\yes&&\yes&&&&\yes\\
Finite checks (Python)&\yes&&\yes&\yes&\yes&\yes&&\yes&\yes\\
\bottomrule
\end{longtable}
\endgroup

\section{Directory map}\label{app:map}

\begin{center}
\small
\renewcommand{\arraystretch}{1.12}
\begin{tabularx}{\textwidth}{@{}lY@{}}
\toprule
\textbf{New}&\textbf{Original location under \nolinkurl{docs/cardinals/research-reports/} (folder, or archive : folder)}\\\midrule
\nolinkurl{01/}&\nolinkurl{Large_Cardinals_Continuation (1)/Large_Cardinals_Continuation/}\\
\nolinkurl{02/}&\nolinkurl{Large_Cardinals_Research_Followup/}\\
\nolinkurl{03/}&\nolinkurl{Large_Cardinals_Research_Continuation/}\\
\nolinkurl{04/}&\nolinkurl{Large_Cardinals_Research_Continuation(1) (1)/Large_Cardinals_Research_Continuation/}\\
\nolinkurl{05/}&\nolinkurl{Large_Cardinals_Research_Continuation (1)/Large_Cardinals_Research_Continuation/}\\
\nolinkurl{06/}&\nolinkurl{Large_Cardinals_Research_Continuation (2)/}\\
\nolinkurl{07/}&\nolinkurl{Large_Cardinals_Continuation/}\\
\nolinkurl{08/}&\nolinkurl{Large_Cardinals_Continuation(1)/Large_Cardinals_Continuation/}\\
\nolinkurl{09/}&\nolinkurl{Large_Cardinals_Research_Continuation(1)/Large_Cardinals_Research_Continuation/}\\
\nolinkurl{10/}&\nolinkurl{Cardinals_Research_Continuation.zip} : \nolinkurl{Cardinals_Continuation/}\\
\nolinkurl{11/}&\nolinkurl{Cardinals3_Research_Continuation.zip} : \nolinkurl{Cardinals3/}\\
\nolinkurl{12/}&\nolinkurl{Large_Cardinals_Maximal_Width.zip} : \nolinkurl{Large_Cardinals_Maximal_Width/}\\
\nolinkurl{13/}&\nolinkurl{Cardinals3_Continuation.zip} : \nolinkurl{Cardinals3_Continuation/}\\
\nolinkurl{14/}&\nolinkurl{Cardinals_Continuation.zip} : \nolinkurl{Cardinals_Continuation/}\\
\nolinkurl{15/}&\nolinkurl{Large_Cardinals_Orbit_Saturation.zip} : \nolinkurl{Orbit_Saturation_Report/}\\
\nolinkurl{16/}&\nolinkurl{Cardinals_Continuation_Small_Fibres.zip} : \nolinkurl{Cardinals_Continuation_Small_Fibres/}\\
\nolinkurl{17/}&\nolinkurl{Large_Cardinals_Thin_Sections.zip} : \nolinkurl{Large_Cardinals_Thin_Sections/}\\
\nolinkurl{18/}&\nolinkurl{Cardinals3_Cofinal_Orbit_Rigidity.zip} : \nolinkurl{Cardinals3_Continuation/}\\
\nolinkurl{19/}&\nolinkurl{Cardinals_Continuation_Fixed_Tails.zip} : \nolinkurl{Fixed_Tail_Parameters/}\\
\nolinkurl{20/}&\nolinkurl{Cardinals4_Cofinality_Cascades.zip} : \nolinkurl{Cardinals4_Cofinality_Cascades/}\\
\nolinkurl{21/}&\nolinkurl{Cardinals4_Finitely_Additive_Kernels.zip} : \nolinkurl{Cardinals4_Continuation/}\\
\nolinkurl{22/}&\nolinkurl{Cardinals4_Tail_Measures.zip} : \nolinkurl{Cardinals4_Tail_Measures/}\\
\nolinkurl{23/}&\nolinkurl{Cardinals4_Normal_Measures.zip} : \nolinkurl{Cardinals4_Normal_Measures/}\\
\nolinkurl{24/}&\nolinkurl{Cardinals4_Stationary_Seeds.zip} : \nolinkurl{Cardinals4_Stationary_Seeds/}\\
\nolinkurl{25/}&\nolinkurl{Cardinals4_Research_Continuation.zip} : \nolinkurl{Cardinals4_Rigid_Classes/}\\
\nolinkurl{26/}&\nolinkurl{Ultraexacting_Prikry_Cores.zip} : \nolinkurl{Ultraexacting_Prikry_Cores/}\\
\nolinkurl{27/}&\nolinkurl{Cardinals4_Normal_Traces.zip} : \nolinkurl{Cardinals4_Normal_Traces/}\\
\nolinkurl{28/}&\nolinkurl{Cardinals5_Finite_Symmetry.zip} : \nolinkurl{Cardinals5_Finite_Symmetry/}\\
\nolinkurl{29/}&\nolinkurl{Cardinals5_Vopenka_Reflection.zip} : \nolinkurl{Cardinals5_Vopenka_Reflection/}\\
\nolinkurl{30/}&\nolinkurl{Cardinals5_Prikry_Symmetry(1).zip} : \nolinkurl{Cardinals5_Prikry_Symmetry/}\\
\nolinkurl{31/}&\nolinkurl{Cardinals5_Prikry_Symmetry.zip} : \nolinkurl{Cardinals5_Prikry_Symmetry/}\\
\nolinkurl{32/}&\nolinkurl{Cardinals5_Prikry_Separation.zip} : \nolinkurl{Cardinals5_Prikry_Separation/}\\
\nolinkurl{33/}&\nolinkurl{Prikry_Finite_Symmetry.zip} : \nolinkurl{Prikry_Finite_Symmetry/}\\
\nolinkurl{34/}&\nolinkurl{Prikry_Choice_Gap.zip} : \nolinkurl{Prikry_Choice_Gap/}\\
\nolinkurl{35/}&\nolinkurl{Cardinals5_Prikry_Choice.zip} : \nolinkurl{Cardinals5_Prikry_Choice/}\\
\nolinkurl{36/}&\nolinkurl{Cardinals5_Prikry_Finite_Choice.zip} : \nolinkurl{Cardinals5_Prikry_Finite_Choice/}\\
\bottomrule
\end{tabularx}
\end{center}

\phantomsection
\addcontentsline{toc}{section}{References}
\begin{thebibliography}{99}
\small
\setlength{\itemsep}{6pt}

\bibitem{plan}
\emph{Large cardinals at the inconsistency frontier: a unified assessment of ZFC-compatible hypotheses, their obstructions, and the evidence on both sides}. Research plan, 18 September 2026; \nolinkurl{docs/cardinals/research-plan/Large_Cardinals_Unified_Report.tex}.

\bibitem{cover}
G.~Goldberg, reporting joint work with D.~Blue. \emph{Consistency beyond the Kunen inconsistency}. Lecture notes, Warsaw Inner Model Theory Meeting, 1 July 2026. \url{https://math.berkeley.edu/~goldberg/Slides/CoverExact.pdf}. Definitions~2.6, 3.1--3.2; Remark~2.4; Lemma~3.3; Proposition~3.4; Theorems~3.5--3.8; Problem~5.2.

\bibitem{unique}
G.~Goldberg. The uniqueness of elementary embeddings. \emph{J.\ Symbolic Logic} \textbf{89}(4) (2024), 1430--1454. \doi{10.1017/jsl.2024.60}. Author version used for numbering; earlier preprint \arxiv{2103.13961} (numbering shifted).

\bibitem{abl}
J.~P.~Aguilera, J.~Bagaria, P.~L\"ucke. \emph{Large cardinals, structural reflection, and the HOD Conjecture}. \arxiv{2411.11568v4}, 2025.

\bibitem{abgl}
J.~P.~Aguilera, J.~Bagaria, G.~Goldberg, P.~L\"ucke. \emph{Large cardinals beyond HOD}. \arxiv{2509.10254v1}, 2025.

\bibitem{ad}
A.~Andretta, V.~Dimonte. \emph{The iterability hierarchy above $\mathrm{I3}$}. \arxiv{1712.03877v2}.

\bibitem{kunen}
K.~Kunen. Elementary embeddings and infinitary combinatorics. \emph{J.\ Symbolic Logic} \textbf{36}(3) (1971), 407--413. \doi{10.2307/2269948}.

\bibitem{hkp}
J.~D.~Hamkins, G.~Kirmayer, N.~L.~Perlmutter. Generalizations of the Kunen inconsistency. \emph{Ann.\ Pure Appl.\ Logic} \textbf{163} (2012), 1872--1890. \arxiv{1106.1951}.

\bibitem{dichotomy}
G.~Goldberg. Strongly compact cardinals and ordinal definability. \emph{J.\ Math.\ Log.} \textbf{24}(1) (2024), 2250010. \arxiv{2107.00513}. \S2.2 ($\omega$-club amenability), Proposition~2.3, Lemma~2.7.

\bibitem{dj}
A.~J.~Dodd, R.~B.~Jensen. The covering lemma for $K$. \emph{Ann.\ Math.\ Logic} \textbf{22} (1982), 1--30; The covering lemma for $L[U]$, ibid., 127--135. See also W.~J.~Mitchell, Definable singularity, \emph{Trans.\ Amer.\ Math.\ Soc.} \textbf{327} (1991), 407--426.

\bibitem{prikry}
K.~Prikry. Changing measurable into accessible cardinals. \emph{Dissertationes Math.} \textbf{68} (1970). See also M.~Gitik, Prikry-type forcings, in \emph{Handbook of Set Theory}, Springer 2010, \S1.

\bibitem{mathias}
A.~R.~D.~Mathias. On sequences generic in the sense of Prikry. \emph{J.\ Austral.\ Math.\ Soc.} \textbf{15} (1973), 409--414.

\bibitem{bukovsky}
L.~Bukovsk\'y. Iterated ultrapower and Prikry's forcing. \emph{Comment.\ Math.\ Univ.\ Carolinae} \textbf{18} (1977), 77--85.

\bibitem{dehornoy}
P.~Dehornoy. Iterated ultrapowers and Prikry forcing. \emph{Ann.\ Math.\ Logic} \textbf{15} (1978), 109--160.

\bibitem{sacks}
G.~E.~Sacks. \emph{Higher Recursion Theory}. Springer, 1990. Chapters~I and~III ($\Sigma^1_1$-boundedness; the effective perfect set theorem).

\bibitem{kechris}
A.~S.~Kechris. \emph{Classical Descriptive Set Theory}. Springer, 1995. Section~8 (the topological zero--one law).

\bibitem{reitz}
J.~Reitz. The Ground Axiom. \emph{J.\ Symbolic Logic} \textbf{72}(4) (2007), 1299--1317. \doi{10.2178/jsl/1203350787}.

\end{thebibliography}
\end{document}
